\documentclass[11pt]{article}
\usepackage[margin=1in]{geometry}
\usepackage{fix-cm}
\usepackage[T2A,T1]{fontenc}
\usepackage[utf8]{inputenc}
\usepackage[english,russian]{babel}
\usepackage{PTSerif}
\usepackage{microtype}
\usepackage{amsmath,amssymb}
\usepackage{enumitem}
\usepackage{ragged2e}
\usepackage{graphicx}
\usepackage{placeins}
\usepackage[hidelinks]{hyperref}

\setlength{\parindent}{0pt}
\setlength{\parskip}{0.55\baselineskip}
\setlist[itemize]{leftmargin=*,itemsep=0.45\baselineskip,parsep=0pt}
\setlist[description]{leftmargin=3em,labelwidth=2.5em,itemsep=0.35\baselineskip}
\setlength{\emergencystretch}{3em}

\title{Ингибиторы транспортеров OATP1B1/OATP1B3 и BCRP: молекулярные механизмы, лекарственные и нутрицевтические взаимодействия, клинические ограничения}
\author{Вербицкий~С.\,В.\\[0.3em]
  {\small ORCID: \href{https://orcid.org/0009-0008-5706-0174}{0009-0008-5706-0174}}}
\date{9 сентября 2026 года}

\begin{document}
\maketitle

\section*{Аннотация}

\textbf{Предпосылки.} Ингибирование полипептидов, транспортирующих органические анионы (OATP1B1/OATP1B3), и белка резистентности рака молочной железы (BCRP) изменяет всасывание лекарственных средств, их поступление в гепатоциты и экскрецию независимо от ингибирования цитохрома P450 (CYP) либо одновременно с ним. Перекрытие субстратных спектров этих транспортеров обусловливает чувствительность фармакокинетики транспортер-зависимых статинов к сочетаниям противовирусных, иммуносупрессивных и таргетных препаратов. Область лекарственных и нутрицевтических взаимодействий также охватывает совместный прием препаратов с пищевыми компонентами, полифенолами и концентрированными биологически активными добавками (БАД); экспериментальную активность последних необходимо отличать от установленного клинического эффекта. [1,6,12,79,81,86,95]

\textbf{Цель.} В настоящем обзоре систематизированы молекулярные механизмы лекарственных и нутрицевтических взаимодействий с участием OATP1B1/OATP1B3 и BCRP, экспериментально установленные отношения экспозиции и клинические ограничения по данным поиска 9 сентября 2026 года. Разграничены результаты исследований ингибирования in vitro, клинические фармакокинетические данные и механистические модели; нутрицевтическая феноконверсия рассматривается как рабочая концепция, а не валидированная диагностическая категория. [1--5,79--97]

\textbf{Механизмы.} Блокада синусоидальных OATP ограничивает поступление вещества в печень; ингибирование кишечного BCRP уменьшает эффлюкс в просвет кишки и может повышать биодоступность при приеме внутрь. Ингибирование BCRP на канальцевой мембране гепатоцитов или в тканевых барьерах имеет иные последствия для внутриклеточного удержания вещества. Поэтому одновременное подавление захвата и эффлюкса может повышать плазменную экспозицию без пропорционального усиления фармакологического действия в печени. Розувастатин представляет собой чувствительный комплексный зонд, а не селективный инструмент оценки какой-либо одной из этих групп транспортеров. [6,9,12,31]

\textbf{Клиническая значимость.} Отношения AUC розувастатина составляют 1.88 при сочетании с гемфиброзилом, 2.15 при исследованной схеме глекапревира/пибрентасвира и 7.1 у получавших циклоспорин реципиентов трансплантата. В исследовании взаимодействия с софосбувиром/\allowbreak велпатасвиром/\allowbreak воксилапревиром, включавшем дополнительную дозу воксилапревира, отношения AUC и максимальной концентрации составили 7.39 и 18.88. Такие величины требуют ограничений, специфичных для субстрата и схемы лечения, а не экстраполяции классификаций, разработанных для CYP. [13,17,25,26]

\textbf{Заключение.} Двойная блокада транспортеров является важным, но не универсальным принципом организации безопасной кардиометаболической полифармакотерапии. Обоснованный подход учитывает транспортерные и ферментативные пути, локальную экспозицию, активные метаболиты, функциональный резерв почек и генотип. Для ингибиторов OATP1B/BCRP отсутствуют общепринятая классификация на «сильные» и «умеренные», единый обязательный период отмывки и универсальный порог токсичности, соответствующий двукратному увеличению экспозиции. Колхицин преимущественно является препаратом-субстратом, экспозиция которого изменяется под действием взаимодействующего средства; экспериментальные антитела, связывающие транспортеры, не следует представлять как установленный класс клинически применяемых ингибиторов с внецелевым действием. [1,23,24,30,34,38]

\subsection*{Цель обзора и  количественные обозначения}

Цель обзора --- систематизировать классические межлекарственные взаимодействия (DDI) и лекарственно-пищевые взаимодействия, включая лекарственно-нутрицевтические (DFDI --- рабочее обозначение этого направления в настоящем обзоре), с участием переносчиков органических анионов OATP1B1/OATP1B3 и эффлюксного белка семейства АТФ-связывающих кассетных транспортеров BCRP. Область обзора охватывает терапевтическое применение препаратов у взрослых и их сочетания с пищевыми субстанциями, полифенолами и БАДами; основное внимание уделено статинам как клиническим зондовым субстратам с измеримой экспозицией и значимым риском токсичности. Временная граница включенных данных соответствует 9 сентября 2026 года, а не всему сентябрю 2026 года. [1,79--97]

Настоящая работа представляет собой обзор с целевым отбором доказательных данных; она не является зарегистрированным систематическим обзором и не претендует на исчерпывающий отбор исследований в рамках PRISMA-ScR. В работе не заявлены независимый отбор исследований двумя рецензентами, формальная оценка риска систематической ошибки, объединенная оценка эффекта или полный количественный учет этапов отбора. Обзор подробно рассматривает выбранные классы препаратов и механизмы, но не позволяет утверждать, что выявлены все зарегистрированные или исследуемые ингибиторы. Установленная дата завершения поиска не означает, что все источники опубликованы в 2026 году: более ранние исследования остаются доказательной основой ряда действующих ограничений. Далее отдельно представлены регуляторные положения, предлагаемый клинический алгоритм и пояснительные математические модели.

AUC обозначает площадь под кривой «концентрация--время», а \(C_{\max}\) --- максимальную плазменную концентрацию. Для обозначения направленности взаимодействия далее разграничены препарат-субстрат, фармакокинетика которого изменяется, и взаимодействующее средство, вызывающее это изменение. Одна молекула может выполнять обе роли в разных сочетаниях.

Обозначения: \(R_{\mathrm{AUC}}=\mathrm{AUC}_{\mathrm{with}}/\mathrm{AUC}_{\mathrm{without}}\) и \(R_{C_{\max}}=C_{\max,\mathrm{with}}/C_{\max,\mathrm{without}}\), где индексы \(\mathrm{with}\) и \(\mathrm{without}\) соответствуют наличию и отсутствию взаимодействующего препарата. Отношение 2 означает увеличение на 100\%, а не на 200\%. Отношения геометрических средних и 90\%-ные доверительные интервалы (ДИ) приводятся, если они указаны в исследовании; использование в источнике арифметических процентных изменений оговаривается отдельно. AUC за интервал дозирования, AUC до последней измеримой концентрации и AUC, экстраполированная до бесконечности, не рассматриваются как взаимозаменяемые показатели. Концентрация полумаксимального ингибирования in vitro (\(\mathrm{IC}_{50}\)) зависит от условий анализа и не может автоматически рассматриваться как эквивалент константы ингибирования (\(K_i\)). [1]

\section{Молекулярная физиология OATP1B1, OATP1B3 и BCRP}

\subsection{Мембранная локализация и направленность транспорта определяют характер риска}

OATP1B1, кодируемый геном \textit{SLCO1B1}, и OATP1B3, кодируемый геном \textit{SLCO1B3}, преимущественно локализованы на синусоидальной мембране гепатоцитов и обеспечивают захват веществ из крови в клетки печени. Направление обеспечиваемого ими транспорта противоположно направлению канальцевого эффлюкса. Криоэлектронная микроскопия выявляет архитектуру из 12 трансмембранных спиралей, организованных в два шестиспиральных пучка с центральной субстратной полостью. Чередование состояний с доступностью со стороны внеклеточной среды и цитоплазмы обеспечивает распознавание лиганда снаружи и его последующее высвобождение внутри клетки. Эти переносчики растворенных веществ не гидролизуют АТФ непосредственно; в транспорте участвуют градиенты субстратов, обменные процессы и локальная химическая среда. Структурные наблюдения, указывающие на участие бикарбоната и чувствительность к pH, не обосновывают единую универсальную стехиометрию обмена для всех субстратов. [6]

Структуры комплексов с билирубином, эстрон-3-сульфатом и симепревиром показывают, что гибкость системы распознавания OATP1B1 позволяет связывать химически разнородные лиганды. Следовательно, низкое значение \(\mathrm{IC}_{50}\), измеренное с одним зондовым субстратом, не доказывает столь же выраженного ингибирования транспорта всех лекарственных средств. Степень занятости участка связывания, концентрация зонда, конформационное состояние и мембранное окружение могут изменять наблюдаемое ингибирование. Профили распознавания OATP1B1 и OATP1B3 перекрываются, но не идентичны; нельзя считать равными ни содержание этих белков, ни их долевой вклад в клиренс. [7]

BCRP, кодируемый геном \textit{ABCG2}, представляет собой полутранспортер семейства АТФ-связывающих кассетных белков, образующий функциональный димер. Каждый протомер содержит нуклеотидсвязывающий домен и трансмембранный домен из шести спиралей. АТФ-зависимые конформационные переходы сопрягают захват субстрата с его выведением наружу. Структурные исследования транспорта противоопухолевых препаратов выявляют субстратную полость, доступность которой меняется в ходе транспортного цикла; само по себе связывание субстрата не позволяет установить, эффективно ли переносится лиганд или он стабилизирует непродуктивное состояние. [9]

Клинически значимая локализация ингибирования зависит от пути введения препарата и его распределения. В энтероцитах апикальный BCRP возвращает субстрат в просвет кишечника и может ограничивать всасывание. В гепатоцитах канальцевый BCRP участвует в выведении в желчь. BCRP также задействован в почечном транспорте и работе тканевых барьеров, хотя количественный вклад каждой локализации различается между субстратами. Высокая концентрация ингибитора в кишечнике после приема внутрь не означает сопоставимой концентрации несвязанной фракции в области гематоэнцефалического барьера или внутри гепатоцитов. Поэтому выражение «ингибирование BCRP» без анатомического уточнения не дает полного описания механизма. [1,9,10,28,41]

\subsection{Эндогенные и лекарственные субстраты}

\begin{itemize}
\item \textbf{Билирубин и конъюгированный билирубин.} Белки OATP1B участвуют в печеночном обмене различных форм билирубина, включая повторный захват конъюгированного билирубина, поступившего в синусоидальную кровь. Сочетанная наследственная утрата OATP1B1 и OATP1B3 вызывает синдром Ротора, что подтверждает существование этого пути рециркуляции у человека. Повышение концентрации билирубина, обусловленное изменением транспорта, само по себе не доказывает наличия некроза гепатоцитов; для интерпретации необходимы определение фракций билирубина, активности аминотрансфераз и оценка клинического контекста. Такой механизм повышения концентрации билирубина следует отличать от нарушений его конъюгации. [7,8]
\item \textbf{Производные желчных кислот, стероидные конъюгаты и копропорфирины.} К субстратам OATP относятся отдельные желчные кислоты и их конъюгированные производные, эстронсульфат, эстрадиолглюкуронид и копропорфирины. Ингибирование OATP не равнозначно подавлению основного натрийзависимого транспортера захвата желчных кислот NTCP или канальцевой помпы экспорта солей желчных кислот BSEP. Копропорфирин I и отдельные конъюгированные желчные кислоты могут служить эндогенными индикаторами ингибирования OATP1B, но не являются валидированными самостоятельными маркерами повреждения. [6,12,37]
\item \textbf{Копропорфирины I и III: чувствительность не равнозначна селективности.} Копропорфирин I (CP-I) и копропорфирин III (CP-III) являются субстратами обеих изоформ OATP1B. В трансфицированных клетках значения \(K_m\) для CP-I и CP-III составляли соответственно 0.13 и \(0.22\,\mu\mathrm{M}\) для OATP1B1, 3.25 и \(4.61\,\mu\mathrm{M}\) для OATP1B3; CP-III также переносился OATP2B1 с \(K_m=0.31\,\mu\mathrm{M}\). Эти параметры относятся к экспериментальным системам и не определяют изолированно долю клиренса у человека. В исследовании 356 здоровых добровольцев концентрация CP-I натощак при генетически низкой функции OATP1B1 была выше на 68\%, а CP-III --- на 27\% по сравнению с нормальной функцией. Следовательно, CP-I лучше выявлял снижение функции OATP1B1, однако меньшая чувствительность CP-III к этому снижению не доказывает его селективности в отношении OATP1B3. Использование пары CP-I/CP-III для количественного разделения изоформ требует дополнительной валидации; CP-III нельзя представлять как установленный селективный клинический маркер OATP1B3. [52,53]

\item \textbf{Холецистокинин-октапептид и селективность OATP1B3.} Холецистокинин-октапептид (CCK-8) позволяет экспериментально различать транспортную активность OATP1B3 и OATP1B1: первый переносит этот пептид, тогда как спектры переносимых органических анионов у двух изоформ существенно перекрываются. Сайт-направленный мутагенез выявил критическую роль остатков Y537, S545 и T550 в десятом трансмембранном сегменте OATP1B3. Это демонстрирует, что сходство аминокислотных последовательностей не означает идентичности субстратных карманов; отсутствие выраженного ингибирования транспорта одного зонда OATP1B1 не исключает влияния на другой субстрат OATP1B3. Селективность CCK-8 в такой сравнительной клеточной системе не следует считать доказательством абсолютной специфичности во всех тканях человека. [44]
\item \textbf{Статины.} Несмотря на ограниченный метаболизм с участием CYP, фармакокинетика розувастатина зависит от печеночного захвата и BCRP-опосредованных процессов. Правастатин и питавастатин позволяют дополнительно оценивать процессы, зависящие от захвата, но не являются транспортер-независимыми альтернативами. Для аторвастатина и активных кислотных форм ряда других статинов захват сочетается с метаболизмом и взаимопревращением кислотной и лактонной форм. Поэтому перед сравнением выраженности взаимодействий необходимо указать субстрат, молекулярную форму и измеряемый показатель. [3,13,25,26,31,42]
\item \textbf{Метотрексат.} BCRP дикого типа переносит метотрексат и отдельные короткоцепочечные полиглутаматы в экспериментальных системах. В одном исследовании с мембранными везикулами кажущийся \(K_m\) метотрексата составлял приблизительно \(680\,\mu\mathrm{M}\), а \(V_{\max}\) --- приблизительно \(2400\,\mathrm{pmol}\,\mathrm{mg}^{-1}\,\mathrm{min}^{-1}\): это свидетельствует о низком сродстве при высокой транспортной емкости в данной везикулярной системе, но не представляет собой константу клиренса у человека. Перенос соединений с более длинными полиглутаматными цепями нельзя считать эквивалентным переносу короткоцепочечных форм. Наряду с BCRP необходимо учитывать печеночный захват, почечное выведение, внутриклеточный транспорт и удержание фолатов; отношения экспозиции, установленные для взаимодействий розувастатина, нельзя переносить на метотрексат. [11]
\item \textbf{Урат и противоопухолевые субстраты.} ABCG2 переносит урат, что связывает снижение его функции с гиперурикемией и подагрой, а также экспортирует противоопухолевые средства, такие как топотекан и митоксантрон. Ингибирование эффлюкса из опухолевых клеток может увеличивать удержание препарата, однако подавление защитного транспорта в нормальных тканях одновременно способно снижать терапевтический индекс. Таким образом, биологическая целесообразность ингибирования зависит от ткани и цели лечения. [9,10,19]
\end{itemize}

\paragraph{Репаглинид: сочетание ингибирования захвата и метаболизма.}
Репаглинид расширяет клиническую значимость обзора за пределы миотоксичности статинов: повышение его системной экспозиции способно усиливать и пролонгировать гипогликемическое действие. Это субстрат OATP1B1, CYP2C8 и CYP3A4, а не селективный зонд только печеночного захвата. В перекрестном исследовании 12 здоровых добровольцев после трехдневного приема гемфиброзила по 600 мг дважды в сутки или итраконазола по 100 мг дважды в сутки, с первой дозой 200 мг, однократный репаглинид 0.25 мг давал \(R_{\mathrm{AUC}}=8.1\) и 1.4 соответственно; сочетание обоих ингибиторов давало \(R_{\mathrm{AUC}}=19.4\). Эти клинические отношения не являются количественной оценкой изолированного ингибирования OATP1B1. PBPK-анализ воспроизводил взаимодействие при совместном учете ингибирования OATP1B1 гемфиброзилом и его глюкуронидом, механизм-зависимого ингибирования CYP2C8 глюкуронидом и ингибирования CYP3A4 итраконазолом. Сочетание репаглинида с гемфиброзилом противопоказано; инструкция также рекомендует избегать сочетания с клопидогрелом. Эти ограничения нельзя заменять наблюдением только за мышечной токсичностью или формальным расчетом транспортерного скринингового отношения. [54--56]

\paragraph{Ситаглиптин: границы отнесения к BCRP-субстратам.}
Ситаглиптин следует включать как пример необходимости проверки транспортерной принадлежности, а не как установленный чувствительный субстрат BCRP. Первичное исследование выявило его транспорт с участием OAT3, OATP4C1 и P-gp и описало почечную элиминацию посредством клубочковой фильтрации и активной секреции. Эти результаты не подтверждают клинически значимой зависимости ситаглиптина от BCRP. В проверенной доказательной базе недостаточно оснований приписывать повышение его системной экспозиции специфической блокаде кишечного BCRP; отсутствие такого подтверждения также не следует трактовать как экспериментально доказанное отсутствие транспорта во всех возможных системах. [57]

\paragraph{Прямые оральные антикоагулянты: транспортный фенотип и риск кровотечения.}
Апиксабан и ривароксабан необходимо рассматривать как субстраты нескольких путей распределения и элиминации, включая P-gp, BCRP и CYP3A, а не как селективные чувствительные клинические зонды BCRP. Для апиксабана перенос посредством P-gp и BCRP подтвержден в трансфицированных клетках; однако в этих же экспериментах кетоконазол существенно подавлял P-gp-опосредованный, но не BCRP-опосредованный транспорт апиксабана. При приеме кетоконазола 400 мг один раз в сутки клинические отношения для однократного апиксабана 10 мг составляли \(R_{\mathrm{AUC}_{0-\infty}}=1.99\) и \(R_{C_{\max}}=1.62\). В отдельном исследовании ривароксабана 10 мг и кетоконазола 400 мг один раз в сутки в равновесном состоянии отношения составляли \(R_{\mathrm{AUC}_{\tau}}=2.58\) и \(R_{C_{\max}}=1.72\). Это измеренные эффекты сочетанного вмешательства в несколько путей, а не доказательство соответствующей кратности повышения риска кровотечения или изолированного кишечного BCRP-эффекта. [60,61,77]

Количественная значимость кишечного эффлюкса для апиксабана остается предметом механистической дискуссии: повторный анализ клинических PK-профилей предложил объяснение наблюдаемых взаимодействий преимущественно изменением активности CYP3A4. Эта интерпретация не отменяет результатов клеточных экспериментов и сама по себе не отменяет ограничений назначения. Для оценки риска кровотечения следует учитывать конкретную комбинацию, функцию почек и печени и сопутствующие фармакодинамические факторы; ни \(\mathrm{IC}_{50}\) ингибирования BCRP, ни одно отношение AUC не устанавливают универсальный риск клинического исхода. [60,61,78]

\paragraph{Флувастатин и активная кислота симвастатина.}
Сопоставление в единой экспериментальной системе, опубликованное в 2025 году, подтвердило различия между статинами: активная кислота симвастатина переносилась OATP1B1 с \(K_m=2.1\,\mu\mathrm{M}\), тогда как переноса посредством OATP1B3 и OATP2B1 в этой системе не выявлено. Оба энантиомера флувастатина переносились OATP1B1 и OATP2B1, но не OATP1B3; для OATP2B1 значения \(K_m\) составляли 0.57 и \(2.5\,\mu\mathrm{M}\) для 3S,5R- и 3R,5S-энантиомеров соответственно. Таким образом, внутри подтвержденного OATP1B-компонента захвата этих соединений в данном исследовании выявлен вклад OATP1B1, но не OATP1B3. Это не означает, что OATP1B1 обеспечивает 100\% общего печеночного или системного клиренса. Для флувастатина необходимо учитывать альтернативный захват и метаболизм, прежде всего CYP2C9; для симвастатина --- CYP3A и взаимопревращение лактонной и кислотной форм. [1,34,58,59]

\subsection{Кинетическая интерпретация ингибирования}

Базовая модель конкурентного ингибирования захвата имеет вид
\[
v_{\mathrm{uptake}}=
\frac{V_{\max}C_{u,\mathrm{substrate}}}
{K_m\left(1+I_u/K_i\right)+C_{u,\mathrm{substrate}}}.
\]
Здесь \(I_u\) --- концентрация несвязанной фракции ингибитора у соответствующей мембраны, а \(C_{u,\mathrm{substrate}}\) --- концентрация несвязанной фракции субстрата. При концентрациях субстрата значительно ниже \(K_m\) ингибирование уменьшает кажущийся клиренс захвата с \(V_{\max}/K_m\) приблизительно до \((V_{\max}/K_m)/(1+I_u/K_i)\). Это уравнение представляет собой пояснительную предельную модель, а не утверждение о конкурентном действии всех ингибиторов OATP или BCRP. Смешанное ингибирование, множественные участки распознавания, медленная диссоциация, накопление в мембране и изменение доступности транспортера могут ограничивать применимость ее допущений. [1,6,7]

Предварительная инкубация особенно важна, когда наблюдаемая ингибирующая активность в отношении OATP возрастает с длительностью воздействия. Анализ без предварительной инкубации может недооценивать ингибирование, развивающееся после распределения вещества внутри клетки или в мембране; напротив, длительное воздействие, нарушающее жизнеспособность клеток, может вызывать неспецифическое ингибирование. Ни один из этих вариантов сам по себе не доказывает ковалентной инактивации транспортера. Описание эксперимента должно включать сведения об экспрессии транспортера, типе и концентрации зондового субстрата, продолжительность предварительной инкубации и захвата, неспецифическое связывание, температуру, а также подходящие контрольные клетки или везикулы с истощенным запасом АТФ. Эти параметры необходимо учитывать при любом количественном сравнении ингибирующей активности. [1,12,13]

Для оценки долевого вклада необходимо явно задать знаменатель. В ненасыщенном приближении долю OATP1B в общем внутреннем клиренсе поступления в гепатоцит можно определить как
\[
f_{\mathrm{uptake,OATP1B}}=
\frac{CL_{\mathrm{int,1B1}}+CL_{\mathrm{int,1B3}}}
{CL_{\mathrm{int,1B1}}+CL_{\mathrm{int,1B3}}+
 CL_{\mathrm{int,other}}+CL_{\mathrm{int,passive}}}.
\]
Здесь все слагаемые относятся к однонаправленному поступлению при согласованных условиях учета несвязанной фракции; \(CL_{\mathrm{int,other}}\) включает другие переносчики. Эта определяемая моделью доля не равна доле OATP1B в системном клиренсе. Ее численное значение для флувастатина или симвастатиновой кислоты нельзя получить только из \(K_m\), генотипического отношения AUC или отсутствия транспорта одной изоформой в клеточной линии. Необходимы данные о транспортной емкости, экспрессии, пассивной проницаемости, метаболизме и обратном эффлюксе; REF/RAF-масштабирование и кинетическая модель баланса должны сопровождаться оценкой неопределенности. Поэтому универсальный фиксированный процент OATP1B-опосредованного клиренса этим статинам здесь не присваивается. [1,58,59]

\section{Классификация «сильных» и «умеренных» ингибиторов транспортеров}

\subsection{Почему категории ингибирующей активности требуют оговорок}

Принятые FDA категории ингибиторов ферментов --- сильные при увеличении AUC чувствительного субстрата не менее чем в пять раз и умеренные при увеличении не менее чем в два, но менее чем в пять раз --- не являются универсально валидированной системой классификации ингибиторов OATP1B1, OATP1B3 или BCRP. Зондовые препараты для изучения транспортеров часто используют параллельные пути, а взаимодействия после приема внутрь объединяют изменения всасывания и клиренса. Перечни примеров FDA указывают ингибиторы транспортеров, но не устанавливают общеприменимую иерархию «сильный/умеренный». Поэтому в настоящем обзоре описываются наблюдаемая выраженность взаимодействия и уровень доказательств, а не присваивается препарату неизменная категория клинической ингибирующей активности, независимая от субстрата. [1,38]

\begin{itemize}
\item \textbf{Клинически подтвержденное взаимодействие.} Контролируемое сравнение у человека или регистрационное исследование взаимодействия дает специфичное для субстрата отношение AUC либо максимальной концентрации. Выраженное увеличение экспозиции клинически значимо, но без дополнительных экспериментов не определяет долю эффекта, обусловленную OATP1B1, OATP1B3 или BCRP. [1,20,31]
\item \textbf{Механистически обоснованное взаимодействие.} Данные об ингибировании в клеточных системах и экспозиции у человека, а также прогнозы физиологически обоснованного фармакокинетического моделирования (PBPK) либо механистических статических моделей согласованно указывают на один или несколько путей. Согласованность прогноза с наблюдаемым отношением плазменной AUC повышает правдоподобие предложенного механизма, но не позволяет однозначно определить параметры тканевого клиренса. [12,31]
\item \textbf{Только экспериментальное ингибирование.} Подавление транспорта в клетках с избыточной экспрессией, анализах АТФазной активности, опухолевых моделях или экспериментах со связыванием антител подтверждает биологическое взаимодействие, но не устанавливает клинического предела дозы и не доказывает увеличения AUC у человека. [19,23,24]
\end{itemize}

\subsection{Компоненты антиретровирусной терапии и фармакокинетические усилители}

Влияние ингибиторов протеазы на розувастатин нельзя объяснить только ингибированием CYP3A. В механистическом исследовании лопинавир, ритонавир, атазанавир и дарунавир подавляли BCRP и OATP1B1 с одинаковым ранжированием по ингибирующей активности: лопинавир превосходил ритонавир, ритонавир --- атазанавир, а дарунавир действовал существенно слабее. Для этих препаратов средние значения \(\mathrm{IC}_{50}\) в отношении BCRP находились в диапазоне от 15.5 до \(143\,\mu\mathrm{M}\), а для OATP1B1 --- от 0.220 до \(9.53\,\mu\mathrm{M}\); использовались разные зонды, соответствующие изучаемому транспортеру. Моделирование указывало на кишечный BCRP и печеночный OATP1B1 как основные звенья взаимодействия; вклад OATP1B3 и NTCP был меньше, особенно при анализе атазанавира. Это оценка относительного вклада механизмов, а не доказательство того, что каждая усиленная схема одинаково подавляет все пути захвата. [12]

\begin{itemize}
\item \textbf{Атазанавир\slash ритонавир и атазанавир\slash кобицистат.} NIH приводит для розувастатина \(R_{\mathrm{AUC}}\) приблизительно 3 и \(R_{C_{\max}}\) приблизительно 7 при сочетании с атазанавиром\slash ритонавиром и соответственно 3.4 и 10.6 при сочетании с атазанавиром\slash кобицистатом. В американской инструкции розувастатина отношение AUC при усилении ритонавиром указано как 3.1. Эти оценки относятся к конкретным схемам, а не прямое доказательство фиксированного дополнительного эффекта при замене усилителя. [3,4]
\item \textbf{Лопинавир/ритонавир.} У 15 здоровых добровольцев, включенных в анализ и получавших розувастатин по 20 мг в сутки, совместное применение увеличивало AUC за интервал дозирования с 47.6 до \(98.8\,\mathrm{ng\,h/mL}\), а максимальную концентрацию --- с 4.34 до \(20.2\,\mathrm{ng/mL}\). Отношения геометрических средних составляли соответственно 2.1 (90\% ДИ 1.7--2.6) и 4.7 (3.4--6.4). У одного участника отмечено бессимптомное повышение активности креатинкиназы до 17 верхних границ нормы. Использованная в исследовании доза не соответствует дозе, рекомендуемой в настоящее время для совместного применения. [27]
\item \textbf{Дарунавир/ритонавир и дарунавир/кобицистат.} NIH указывает для розувастатина \(R_{\mathrm{AUC}}=1.48\) и \(R_{C_{\max}}\) приблизительно 2.4 при сочетании с дарунавиром/ритонавиром, тогда как при сочетании с дарунавиром/кобицистатом соответствующие отношения составляли приблизительно 1.9 и 3.8. Непропорциональное увеличение максимальной концентрации остается значимым даже при увеличении AUC менее чем в два раза. [4]
\item \textbf{Ритонавир и кобицистат: комплексное усиление, а не простое сложение.} Фармакокинетические усилители повышают экспозицию ингибитора протеазы и сами могут влиять на транспорт и метаболизм. Наблюдаемый клинический эффект определяется экспозицией усилителя, экспозицией второго компонента и зависимостью препарата-субстрата от каждого затронутого пути. Приведенные отношения характеризуют схемы целиком; деление отношений для схемы с кобицистатом на отношения для схемы с ритонавиром не выделяет причинный эффект усилителя. Для установления такого причинного вклада необходимы сопоставления усиленных и неусиленных схем в одинаковых условиях. [4,12]
\end{itemize}

При выборе терапии профили взаимодействий розувастатина и аторвастатина следует рассматривать раздельно. Аторвастатин дополнительно чувствителен к воздействию на CYP3A, тогда как симвастатин и ловастатин вызывают особую обеспокоенность при сочетании с усиленными ингибиторами протеазы. Статин с минимальным участием CYP в метаболизме не является автоматически статином с минимальным риском взаимодействий; противовирусную схему не следует менять исключительно ради упрощения гиполипидемической терапии без обеспечения сохранения вирусологической эффективности. [4,5]

\begin{itemize}
\item \textbf{Питавастатин и правастатин: выбор альтернативы с учетом схемы.} В таблицах NIH дарунавир/ритонавир снижает AUC питавастатина на 26\%, то есть до отношения 0.74; для рассмотренного сочетания коррекция дозы не требуется. Напротив, AUC правастатина увеличивается на 81\% после однократного приема и на 23\% в равновесном состоянии, что соответствует отношениям 1.81 и 1.23; рекомендуется минимальная эффективная доза с наблюдением. Эти различия между субстратами и режимами приема исключают универсальное правило «заменить на любой статин без CYP3A». [4]
\end{itemize}

\subsection{Противовирусные препараты прямого действия для лечения гепатита C}

\begin{itemize}
\item \textbf{Глекапревир.} Экспериментальные значения \(\mathrm{IC}_{50}\) составляли \(0.017\,\mu\mathrm{M}\) для OATP1B1, \(0.064\,\mu\mathrm{M}\) для OATP1B3, \(2.3\,\mu\mathrm{M}\) для BCRP и \(0.33\,\mu\mathrm{M}\) для P-гликопротеина (P-gp). Ингибирование захвата при наномолярных концентрациях в этих анализах указывает на высокую ингибирующую активность в данных экспериментальных условиях, однако клинический эффект по-прежнему зависит от экспозиции несвязанной фракции в печени и долевого вклада отдельных путей в клиренс субстрата. [13]
\item \textbf{Пибрентасвир.} Входящий в комбинацию ингибитор NS5A подавлял P-gp, BCRP и OATP1B1; опубликованные значения \(\mathrm{IC}_{50}\) составляли соответственно 0.036, 14 и \(1.3\,\mu\mathrm{M}\). Значение для OATP1B1 получено в присутствии 4\% альбумина; его нельзя непосредственно сопоставлять со значением из среды, не содержащей белка или интерпретировать как константу для несвязанного вещества. Профиль пибрентасвира не следует отождествлять с профилем глекапревира. Клинические наблюдения для комбинации подтверждают ингибирование P-gp-, BCRP- и OATP-зависимых процессов фармакокинетики субстратов, но не означают, что каждый компонент одинаково сильно подавляет каждый транспортер. [13,14]
\item \textbf{Глекапревир/пибрентасвир со статинами.} При приеме розувастатина по 5 мг в сутки на фоне исследованной противовирусной схемы 400/120 мг в сутки наблюдались \(R_{\mathrm{AUC}}=2.15\) (90\% ДИ 1.88--2.46) и \(R_{C_{\max}}=5.62\) (4.80--6.59). При приеме аторвастатина по 10 мг в сутки на фоне той же исследованной дозы противовирусных препаратов отношения составляли 8.28 (6.06--11.3) и 22.0 (16.4--29.6). Доза глекапревира в этих исследованиях взаимодействий составляла 400 мг, а не зарегистрированные для взрослых 300 мг данного компонента. Это результаты исследований, положенных в основу инструкции, а не измерения при каждой зарегистрированной схеме. [14]
\item \textbf{Велпатасвир.} При приеме велпатасвира по 100 мг в сутки в сочетании с однократной дозой розувастатина 10 мг наблюдались \(R_{\mathrm{AUC}}=2.69\) (90\% ДИ 2.46--2.94) и \(R_{C_{\max}}=2.61\) (2.32--2.92). Это исследование отдельного компонента обосновывает ограничения совместного применения для софосбувира/велпатасвира; софосбувир в данном сравнении с розувастатином не применялся. При сочетании с софосбувиром/велпатасвиром отношения для аторвастатина составляли 1.54 и 1.68. Различия изменений экспозиции субстратов исключают единый множитель экспозиции для всего класса статинов. [16]
\item \textbf{Тройная терапия, содержащая воксилапревир.} У 19 здоровых участников при приеме софосбувира/\allowbreak велпатасвира/\allowbreak воксилапревира в дозе 400/100/100 мг с дополнительными 100 мг воксилапревира в сутки и однократном приеме розувастатина 10 мг получены \(R_{\mathrm{AUC}}=7.39\) (90\% ДИ 6.68--8.18) и \(R_{C_{\max}}=18.88\) (16.23--21.96). Дополнительный воксилапревир --- существенная особенность дизайна: эти значения нельзя описывать как результат прямого эксперимента с использованием только одной стандартной таблетки тройной комбинации. [17]
\item \textbf{Различия между субстратами при тройной схеме.} В аналогичном исследовании с дополнительным воксилапревиром при приеме правастатина в дозе 40 мг отношения AUC и максимальной концентрации составляли 2.16 и 1.89, то есть были значительно меньше, чем у розувастатина. Это различие согласуется с неодинаковой зависимостью от захвата и эффлюкса, но не позволяет численно выделить вклад BCRP или доказать фармакологическую синергию всех трех противовирусных компонентов. Софосбувир нельзя считать просто третьим эквивалентным ингибитором OATP/BCRP. [17,18]
\end{itemize}

Таким образом, максимальные наблюдавшиеся изменения экспозиции статинов выходят за часто приводимый диапазон увеличения AUC в два--семь раз, а отношения максимальной концентрации могут превышать десять. При механистической интерпретации и выборе терапии необходимо сохранять сведения о препарате-субстрате, дозе противовирусного средства и условиях однократного приема либо равновесного состояния, а не рассматривать «тройную терапию гепатита C» как универсальную категорию экспозиции. [14,17]

\subsection{Онкология и таргетная терапия}

\begin{itemize}
\item \textbf{Иматиниб и гефитиниб.} Оба ингибитора тирозинкиназ взаимодействуют с ABCG2 при субмикромолярных концентрациях в экспериментальных системах. Подавление транспорта и изменение накопления субстрата являются эффектами, отличными от воздействия на их основные киназные мишени. Молекула может транспортироваться при одной концентрации и ингибировать транспорт при другой в зависимости от сродства и скорости транспортного цикла. Эти эксперименты не устанавливают универсальное клиническое отношение AUC розувастатина; такое число здесь не приводится. Преодоление устойчивости опухолевых клеток само по себе не обосновывает ограничение совместного назначения сердечно-сосудистых препаратов. [19]
\item \textbf{Даролутамид.} Этот антагонист андрогенного рецептора не относится к ингибиторам тирозинкиназ. При применении даролутамида по 600 мг два раза в сутки экспозиция розувастатина увеличивалась приблизительно в 5.2 раза, а максимальная концентрация --- примерно в пять раз. Клинические данные и результаты in vitro указывают на участие BCRP и, вероятно, OATP1B1/OATP1B3. Слабое влияние на зондовые субстраты CYP при выраженном взаимодействии с розувастатином делает этот случай показательным примером преимущественно транспортер-опосредованного риска. [20]
\item \textbf{Регорафениб.} В клиническом исследовании взаимодействия этот мультикиназный ингибитор увеличивал AUC розувастатина в 3.8 раза, а максимальную концентрацию --- в 4.6 раза. В его инструкции рассматривается повышение экспозиции субстратов BCRP. Ни это наблюдение, ни соответствующее ограничение дозы не доказывают аналогичного профиля взаимодействий у всех ингибиторов киназ. [21]
\item \textbf{Капматиниб.} Совместное применение увеличивало AUC розувастатина на 108\%, а максимальную концентрацию на 204\%, что эквивалентно отношениям 2.08 и 3.04. Инструкция указывает на риск взаимодействия через BCRP при отсутствии ингибирования OATP1B1 или OATP1B3 in vitro. Это клинически значимый пример, опровергающий утверждение, что каждое взаимодействие розувастатина требует двойной блокады. [22]
\item \textbf{Белумосудил.} Ингибитор ROCK2, применяемый при хронической реакции «трансплантат против хозяина», расширяет карту взаимодействий таргетных средств за пределы классических ингибиторов тирозинкиназ. Инструкция REZUROCK приводит для розувастатина \(R_{\mathrm{AUC}}=4.6\) и \(R_{C_{\max}}=3.6\), рассматривая его как субстрат BCRP/OATP1B1. Эти клинические значения следует отличать от прогнозируемых взаимодействий с зондовыми субстратами ферментов, также приведенных в инструкции: для мидазолама увеличение AUC приблизительно в 1.7 раза является модельным прогнозом, а не результатом прямого клинического сравнения. Вклад каждого транспортера в суммарное взаимодействие с розувастатином отдельно не установлен. [47]
\item \textbf{Момелотиниб.} При многократном применении момелотиниба по 200 мг один раз в сутки и однократном приеме розувастатина 10 мг AUC розувастатина увеличивалась на 170\%, а \(C_{\max}\) --- на 220\%; соответствующие отношения равны 2.70 и 3.20, а не 1.70 и 2.20. Инструкция относит взаимодействие к ингибированию BCRP. Сам момелотиниб одновременно является субстратом OATP1B1/OATP1B3; следовательно, один препарат может выступать ингибитором для одного звена сети и препаратом-субстратом для другого. Это обстоятельство существенно при сочетании таргетной терапии с дополнительным ингибитором печеночного захвата. [48]
\item \textbf{Экспериментальные антитела, связывающие транспортеры.} Антитело 5D3 к ABCG2 распознает внеклеточный эпитоп, чувствительный к конформации белка. Высокие концентрации антитела вызывали частичное ингибирование транспорта в клеточных экспериментах; в одной системе подавление транспорта красителя Hoechst, опосредованного ABCG2 дикого типа, составляло приблизительно 30--40\%. Это прямое мишень-специфическое экспериментальное ингибирование, а не клиническое внецелевое взаимодействие зарегистрированного моноклонального антитела. Выраженность связывания антитела может изменяться без пропорционального изменения транспорта, что дополнительно ограничивает интерпретацию анализа. [23]
\item \textbf{Терапевтические моноклональные антитела и конъюгаты «антитело--лекарственное средство».} Нельзя приписывать противоопухолевым антителам как классу клинически установленное прямое внецелевое ингибирование OATP1B/BCRP. Терапевтические белки могут изменять экспрессию ферментов или транспортеров, зависящую от воспаления; высвобождаемый низкомолекулярный лекарственный компонент конъюгата также может обладать потенциалом взаимодействий. Эти косвенные или опосредованные лекарственным компонентом механизмы требуют отдельной оценки и не эквивалентны непосредственному связыванию транспортера интактным антителом. Подход FDA к терапевтическим белкам основан на оценке риска, а не на предположении о взаимозаменяемости антител как ингибиторов эффлюкса. [24]
\end{itemize}

\paragraph{Азольные противогрибковые препараты как сопутствующая терапия.}
Кетоконазол, итраконазол и вориконазол не образуют однородного класса «сильных ингибиторов BCRP». Итраконазол подавляет BCRP и P-gp in vitro, тогда как вориконазол в сравнительном везикулярном исследовании не достигал более чем 60\%-ного подавления исследованных ABC-транспортеров даже при \(250\,\mu\mathrm{M}\). Сильное ингибирование CYP3A азолом нельзя автоматически переносить на BCRP. При итраконазоле 200 мг в сутки в течение пяти дней и розувастатине 10 мг на четвертый день \(R_{\mathrm{AUC}_{0-ct}}\) составляло 1.39, а \(R_{C_{\max}}\) --- 1.36; здесь \(0-ct\) обозначает интервал до последней общей для сравниваемых периодов точки с количественно определяемой концентрацией. Напротив, кетоконазол 200 мг дважды в сутки в течение семи дней не увеличивал экспозицию однократного розувастатина 80 мг на четвертый день: \(R_{\mathrm{AUC}_{0-t}}=1.016\), \(R_{C_{\max}}=0.954\). Историческая исследовательская доза розувастатина 80 мг не является рекомендуемой лечебной дозой. Эти положительные и отрицательные клинические результаты опровергают универсальное утверждение о резком повышении биодоступности всех BCRP-субстратов под действием любого азола. [68--70]

Для противоопухолевых средств и антикоагулянтов клиническая значимость азолов определяется конкретным сочетанием путей CYP3A, P-gp, BCRP и других систем. Повышение экспозиции апиксабана или ривароксабана при кетоконазоле не позволяет количественно выделить кишечный BCRP из суммарного взаимодействия; отсутствие выраженного BCRP-ингибирования у вориконазола также не означает отсутствия опасных CYP-опосредованных взаимодействий. [60,61,68,77]

\subsection{Сердечно-сосудистые, иммуносупрессивные и другие препараты}

\begin{itemize}
\item \textbf{Циклоспорин.} У реципиентов сердца в стабильном состоянии, получавших иммуносупрессивную терапию с циклоспорином, при приеме розувастатина по 10 мг в сутки AUC и максимальная концентрация были соответственно в 7.1 и 10.6 раза выше значений в группе сравнения, включавшей здоровых добровольцев. Это исследование не представляло собой рандомизированное внутрииндивидуальное сравнение, позволяющее изолированно оценить влияние циклоспорина; трансплантация, сопутствующая терапия и популяционные различия ограничивают возможность разграничить влияние отдельных факторов. В сопутствующем эксперименте по захвату кажущаяся константа сродства розувастатина составляла около \(8.5\,\mu\mathrm{M}\), а циклоспорин ингибировал захват с \(\mathrm{IC}_{50}\) около \(2.2\,\mu\mathrm{M}\) при концентрации розувастатина \(5\,\mu\mathrm{M}\). В других системах получены иные константы ингибирования; значения из разных анализов нельзя объединять в единую оценку активности. [25]
\item \textbf{Гемфиброзил.} В исследовании гемфиброзил по 600 мг или плацебо назначали два раза в сутки в течение семи дней; розувастатин 80 мг принимали однократно утром четвертого дня каждого периода. Для AUC от момента введения до последней количественно определяемой концентрации получено \(R_{\mathrm{AUC}}=1.88\) (90\% ДИ 1.60--2.21) и \(R_{C_{\max}}=2.21\) (1.81--2.69). Доза розувастатина 80 мг была экспериментальной и превышает максимальную суточную дозу, разрешенную к применению в настоящее время. Ингибирование захвата помогает объяснить взаимодействие; характеристика гемфиброзила как универсального мощного двойного ингибитора OATP/BCRP некорректна. [26]
\item \textbf{Ацилглюкуронид гемфиброзила.} Этот метаболит является механизм-зависимым ингибитором CYP2C8, что показывает необходимость совместного рассмотрения транспортерных и ферментативных путей. В микросомальных экспериментах предварительная инкубация снижала его \(\mathrm{IC}_{50}\) для CYP2C8 с 24 до \(1.8\,\mu\mathrm{M}\); опубликованные параметры инактивации включали \(K_I=20\)--\(52\,\mu\mathrm{M}\) и \(k_{\mathrm{inact}}=0.21\,\mathrm{min}^{-1}\). Это параметры инактивации фермента, а не константы OATP или BCRP. Они особенно значимы для препаратов-субстратов, чувствительных к CYP2C8, и не должны использоваться для объяснения всех изменений экспозиции розувастатина. [39]
\item \textbf{Тикагрелор.} В рандомизированном перекрестном исследовании у девяти здоровых добровольцев первую дозу тикагрелора 90 мг или плацебо назначали за один час до однократного приема розувастатина 10 мг; тикагрелор 90 мг или плацебо повторно вводили через 12, 24 и 36 часов после первой дозы. При этой схеме оба отношения, для AUC и максимальной концентрации розувастатина, составили приблизительно 2.6; 90\% ДИ равнялись соответственно 1.8--3.8 и 1.7--4.0. Экспериментальные данные об ингибировании и механистическое моделирование указывали на ингибирование кишечного BCRP как основной механизм взаимодействия. Американская инструкция розувастатина от апреля 2026 года отдельно приводит отношения AUC 2.3 и максимальной концентрации 2.4--2.6 для указанных в ней исследований тикагрелора. Эти оценки следует сохранять с привязкой к источнику, а не усреднять для получения искусственной «консенсусной» оценки. [3,28]
\item \textbf{Тикагрелор в клинических популяциях.} В исследовании после инфаркта миокарда у пациентов, получавших розувастатин по 40 мг в сутки, остаточные концентрации розувастатина при тикагрелоре были приблизительно в 1.9 раза выше, чем при прасугреле или клопидогреле: медианы составляли соответственно 9.7, 5.1 и \(5.0\,\mathrm{ng/mL}\). Различия остаточных концентраций подтверждают значимость взаимодействия в реальной практике, но не являются измеренными отношениями AUC и сами по себе не доказывают причинной роли конкретного транспортера. Исследованную дозу 40 мг нельзя считать совместимой с американским ограничением при тикагрелоре, введенным в апреле 2026 года. [3,29]
\item \textbf{Тафамидис.} Ингибирование BCRP не ограничивается противовирусными и противоопухолевыми средствами. Инструкция тафамидиса сообщает об увеличении AUC розувастатина на 96.75\% и максимальной концентрации на 85.59\%, что соответствует отношениям приблизительно 1.97 и 1.86. В таблице CRESTOR уточнена исследованная схема: тафамидис 61 мг два раза в сутки в дни 1 и 2, затем один раз в сутки в дни 3--9; розувастатин 10 мг однократно. Поэтому описание только поддерживающей суточной дозы 61 мг не должно скрывать начальную нагрузочную фазу. Данные инструкции тафамидиса не указывают на сопоставимый потенциал ингибирования OATP1B1/OATP1B3; это еще один пример преимущественно BCRP-опосредованного взаимодействия. [3,49]
\item \textbf{Фебуксостат в сравнении с аллопуринолом.} В трехпериодном рандомизированном перекрестном исследовании у 10 здоровых добровольцев фебуксостат по 120 мг в сутки в течение четырех дней при однократном приеме розувастатина 10 мг увеличивал AUC в 1.9 раза (90\% ДИ 1.5--2.5), а \(C_{\max}\) --- в 2.1 раза (1.8--2.6), без изменения периода полувыведения или почечного клиренса розувастатина. Аллопуринол по 300 мг в сутки в течение семи дней не изменял его фармакокинетику. В везикулах, содержащих BCRP, \(\mathrm{IC}_{50}\) фебуксостата для транспорта розувастатина составляла \(0.35\,\mu\mathrm{M}\); аллопуринол не подавлял транспорт при концентрациях до \(200\,\mu\mathrm{M}\). Совокупность этих результатов согласуется с кишечным механизмом взаимодействия, но не оправдывает экстраполяцию взаимодействия на весь класс ингибиторов ксантиноксидазы. Исследования доза фебуксостата не является универсальной рекомендацией по его назначению. [50]
\item \textbf{Колхицин.} Колхицин преимущественно является субстратом CYP3A4/P-gp и потенциальным источником аддитивной миотоксичности, а не установленным клинически сильным или умеренным ингибитором OATP1B1/OATP1B3/BCRP. В исследованиях взаимодействий COLCRYS сообщается о среднем увеличении AUC колхицина на 259\% при циклоспорине и на 296\% при ритонавире, что соответствует отношениям 3.59 и 3.96, рассчитанным из арифметических процентных изменений. Эти результаты относятся к колхицину как препарату-субстрату, экспозиция которого изменяется при взаимодействии, и не доказывают, что сам колхицин увеличивает AUC статина в такой степени. Статины могут усиливать нервно-мышечную токсичность даже без доказанного транспортер-опосредованного взаимодействия, изменяющего их собственную фармакокинетику. [30]
\item \textbf{Рифампицин: острое ингибирование обеих изоформ и повторное воздействие.} Однократное воздействие ингибирует OATP1B1 и OATP1B3 и используется в клинических исследованиях захвата. Предварительная инкубация in vitro усиливала ингибирующее действие на обе изоформы, снижая определяемые \(K_i\) примерно в 3 и 2.4 раза соответственно; наблюдалось также уменьшение захвата CCK-8 в культивируемых гепатоцитах. Эти данные подтверждают временную зависимость экспериментального эффекта, но не устанавливают длительное селективное ингибирование OATP1B3 у пациентов. При повторном назначении прямое ингибирование сосуществует с индукцией метаболических и транспортных путей. В исследовании репаглинида 4 мг после семидневного приема рифампицина медианная AUC снижалась приблизительно на 50\% при одновременном приеме с последней дозой и на 80\% при интервале 24 часа. Это клинический пример зависимости суммарного эффекта от расписания приема, а не способ измерения изолированной активности OATP1B3. Ни более низкая \(\mathrm{IC}_{50}\), ни эффект предварительной инкубации не обосновывают характеристику рифампицина как «мощного долгосрочного ингибитора именно OATP1B3». [1,37,64,65]
\end{itemize}

\paragraph{Макролиды: одновременное воздействие на захват и CYP3A4.}
Кларитромицин и эритромицин ингибируют CYP3A4 и способны подавлять OATP1B1/OATP1B3 in vitro, создавая предпосылки для сочетанного ограничения поступления статина в гепатоцит и его метаболизма. При использовании сульфобромофталеина как зонда \(\mathrm{IC}_{50}\) для OATP1B3 составляла 32 и \(34\,\mu\mathrm{M}\) соответственно; эти значения не характеризуют автоматически ингибирование каждого статина при терапевтической экспозиции. В популяционной когорте постоянных пользователей CYP3A4-метаболизируемых статинов старше 65 лет назначение кларитромицина или эритромицина вместо азитромицина ассоциировалось с госпитализацией по поводу рабдомиолиза в течение 30 дней: относительный риск 2.17, 95\% ДИ 1.04--4.53; абсолютная разность рисков 0.02 процентного пункта. Это клиническая ассоциация для конкретной популяции, преимущественно получавшей аторвастатин; исследование не разделяет вклад CYP3A4 и OATP1B и не доказывает неизбежности рабдомиолиза. [62,63]

\paragraph{Статины как взаимодействующие средства: пределы перекрестного ингибирования.}
Насыщаемый путь захвата допускает конкуренцию нескольких субстратов, однако высокая назначенная доза сама по себе не доказывает достаточной несвязанной концентрации для ингибирования OATP. Для аторвастатина существует клинический пример: в исследовании 24 генотипированных добровольцев на фоне 40 мг в сутки в подгруппе \textit{SLCO1B1} c.521TT экспозиция однократного репаглинида 0.25 мг увеличивалась на 18\% по AUC и на 41\% по \(C_{\max}\). Авторы предложили ингибирование OATP1B1 как вероятное объяснение; эти процентные изменения не следует представлять как универсальные отношения геометрических средних или доказательство сильного ингибирования. Для розувастатина сопоставимое утверждение о клинически значимом подавлении захвата сопутствующих субстратов при высоких дозах в проверенных источниках не подтверждено. В частности, исследование взаимодействий с пероральными антидиабетическими препаратами оценивало преимущественно обратное направление --- подавление ими захвата розувастатина in vitro. Направленность причинного воздействия нельзя менять при цитировании. [1,66,67]

\paragraph{Ингибиторы протонной помпы: экспериментальная активность и отрицательный клинический результат.}
Омепразол и пантопразол ингибировали BCRP-опосредованный транспорт метотрексата в мембранных везикулах Sf9 с \(\mathrm{IC}_{50}=36\) и \(13\,\mu\mathrm{M}\) соответственно. Это субстрат- и система-зависимая ингибирующая активность in vitro, а не валидированная клиническая категория «умеренные ингибиторы BCRP». В исследовании 36 здоровых мужчин пантопразол 40 мг не изменял фармакокинетику кишечнорастворимого сульфасалазина 500 мг; фамотидин использовался как контроль влияния желудочно-кишечного pH. Отрицательный результат относится к конкретной лекарственной форме, дозе и дизайну, но не поддерживает предположение о неизбежном клинически значимом BCRP-взаимодействии при применении ИПП. Отрицательный результат с сульфасалазином не отменяет описанного риска взаимодействия ИПП с метотрексатом: экспериментальное обоснование роли BCRP и клинические сообщения о задержке элиминации должны рассматриваться отдельно от предположения о классовом эффекте на все субстраты. Продолжительность назначения не заменяет оценку концентрации у мишени; влияние на кислотность, растворение и другие пути элиминации необходимо отличать от BCRP-опосредованного эффекта. [2,71,72]

\paragraph{Сульфасалазин: субстрат и экспериментальный ингибитор.}
Сульфасалазин, противовоспалительный базисный препарат, включен в ICH M12 как пример клинического субстрата кишечного BCRP. Вместе с тем он подавлял BCRP-, OAT1- и OAT3-опосредованный транспорт метотрексата in vitro при \(\mathrm{IC}_{50}<5\,\mu\mathrm{M}\). В той же работе сопоставление ингибирующих концентраций с ожидаемой экспозицией не предсказывало клинически значимого взаимодействия с метотрексатом. Поэтому наличие субстратных и ингибирующих свойств не делает сульфасалазин универсальным клиническим ингибитором BCRP; чувствительность его применения как зонда зависит, в частности, от лекарственной формы, дозы и вариабельности всасывания. [1,72,73]

\paragraph{Телмисартан: выраженность эффекта зависит от PK-показателя.}
Телмисартан является подтвержденным примером сочетания ингибирующей активности в отношении BCRP in vitro с измеренным взаимодействием у человека. После телмисартана 40 мг в сутки в течение 14 дней при розувастатине 10 мг \(C_{\max}\) увеличивалась на 71\%, а AUC --- на 26\%; исследование включало здоровых добровольцев с однократным приемом розувастатина и пациентов, длительно получавших его. Отдельная PBPK-модель другой схемы воспроизводила преимущественное увеличение \(C_{\max}\) и сокращение \(T_{\max}\) через ингибирование кишечного BCRP. Это модельная интерпретация, а не прямое измерение энтероцитарной экспозиции или доказательство изолированного BCRP-механизма. Данные обосновывают внимание к конкретной комбинации, но не общеклассовую характеристику сартанов как сильных ингибиторов BCRP и не автоматический запрет сочетания. [74,75]

\subsection{Влияние нутрицевтиков и компонентов пищи на транспортные сети OATP1B и BCRP}

\paragraph{Объект оценки: определенный продукт, а не название растения.}
Термин «нутрицевтик» далее используется как описательное обозначение пищевых компонентов и добавок, а не как единая регуляторная категория. Заваренный напиток, обычная пищевая порция, стандартизованный экстракт, очищенный полифенол и композиция с биоусилителем не являются взаимозаменяемыми вмешательствами. Для оценки необходимы содержание индивидуальных веществ, суточная и разовая дозы, состав матрицы, форма высвобождения, прием натощак или с пищей и сопутствующие добавки. Концентрация исходного агликона в кишечнике, его несвязанная концентрация на входе в печень и концентрации циркулирующих конъюгатов относятся к разным компартментам. Низкая системная биодоступность не исключает кишечного взаимодействия, но препятствует прямому переносу высокой внутрипросветной концентрации на синусоидальные OATP1B. [79,80,83,85,87,93--95]

\paragraph{EGCG: зависимость эффекта от субстрата.}
Эпи\-галло\-катехин-3-галлат (EGCG) непосредственно изменяет транспорт с участием OATP in vitro, однако его нельзя характеризовать единым диапазоном \(\mathrm{IC}_{50}=0.5\)--\(2.5\,\mu\mathrm{M}\) и универсальным конкурентным либо смешанным механизмом для OATP1B1/OATP1B3. В первичном исследовании ингибирования захвата эстрон-3-сульфата через OATP1B1 получено \(\mathrm{IC}_{50}=7.8\,\mu\mathrm{M}\); для транспорта Fluo-3 через OATP1B3 --- \(8.4\,\mu\mathrm{M}\), причем ингибирование характеризовалось как неконкурентное. Напротив, транспорт эстрон-3-сульфата через OATP1B3 при низкой концентрации субстрата стимулировался, а транспорт эстрадиол-17\(\beta\)-глюкуронида не изменялся. Сам EGCG являлся субстратом OATP1B3. Таким образом, принадлежность лиганда к ингибиторам не исключает стимуляции другого субстратного потока через ту же изоформу. Универсальная характеристика EGCG как «умеренного ингибитора BCRP» с \(\mathrm{IC}_{50}=5\)--\(10\,\mu\mathrm{M}\) также не подтверждена рассмотренными первичными данными; в эксперименте с дасатинибом выраженного подавления BCRP-опосредованного эффлюкса чайными катехинами не обнаружено. [79,88]

При однократном приеме смеси катехинов Polyphenon E натощак средняя максимальная плазменная концентрация неконъюгированного EGCG была более чем в 3.5 раза выше, чем при приеме с пищей. Термин \textit{free} в этом исследовании означает неконъюгированную химическую форму, а не измеренную несвязанную с белками фракцию. Результат обосновывает необходимость учитывать прием натощак при оценке экспозиции концентрированных экстрактов, но не доказывает достижения ингибирующей концентрации одновременно у кишечного BCRP и печеночных OATP1B. Модель двойной блокады остается условным сценарием, для проверки которого нужны данные о локальной несвязанной экспозиции и чувствительности конкретного субстрата; исследованные высокие дозы катехинов не являются рекомендацией к применению. [79,80]

Клинические результаты с розувастатином не подтверждают обязательного повышения его системной экспозиции под действием EGCG. В исследовании у 13 здоровых добровольцев совместный однократный прием розувастатина 20 мг и EGCG 300 мг сопровождался \(R_{\mathrm{AUC}}=0.81\) для AUC до последней измеренной концентрации (90\% ДИ 0.67--0.97) и \(R_{C_{\max}}=0.85\) (0.70--1.04); последнее различие не было статистически значимым. Отбор крови ограничивался восемью часами. После десятидневного предварительного приема EGCG статистически значимых различий основных PK-показателей не выявлено. В отдельном перекрестном исследовании у 20 мужчин экстракт с EGCG 800 мг/сут в течение 14 дней также снижал экспозицию розувастатина 10 мг. Это согласуется с возможным уменьшением кишечного захвата, но само по себе не идентифицирует транспортер и не доказывает изменения клинической гиполипидемической эффективности. [81,82]

\paragraph{Пиперин: многомишенный биоусилитель, а не селективный зонд BCRP.}
Пиперин, в том числе в составах с торговым обозначением BioPerine, следует оценивать по фактической дозе и составу продукта. В экспериментальном исследовании с силибином он ингибировал BCRP- и MRP2-опосредованный эффлюкс в клеточных системах; увеличение биодоступности силибина показано у крыс, а не у пациентов, получающих розувастатин или антикоагулянты. Другие первичные эксперименты демонстрируют ингибирование P-гликопротеина и CYP3A4, что исключает характеристику пиперина как селективного ингибитора ABCG2. Эти функциональные данные не устанавливают конкретный участок связывания и не доказывают блокирования определенного АТФ-зависимого конформационного перехода BCRP. [83,84]

При достаточной локальной экспозиции уменьшение кишечного эффлюкса способно повысить биодоступность чувствительного субстрата; одновременное подавление CYP3A4 или P-гликопротеина может изменять суммарный эффект. Однако для сочетаний пиперина с розувастатином, апиксабаном и ривароксабаном в рассмотренных первичных источниках не установлены клинические \(R_{\mathrm{AUC}}\), \(R_{C_{\max}}\) или величина дополнительного риска кровотечения. Поэтому «резкий скачок» \(C_{\max}\) не следует представлять как наблюдавшийся результат. Обоснование осторожности состоит в многомишенном потенциале, вариабельности экспозиции и тяжести возможных последствий, а не в доказанном полном выключении кишечного BCRP. [60,61,77,78,83,84]

\paragraph{Силимарин и силибинин: состав смеси, экспозиция и тканевая фармакодинамика.}
Силимарин представляет собой смесь флавонолигнанов, тогда как силибинин включает диастереомеры силибин A и силибин B; эти обозначения нельзя использовать как синонимы единого химического ингибитора. При исследовании захвата эстрадиол-17\(\beta\)-глюкуронида значения \(\mathrm{IC}_{50}\) для силибина A составляли 9.7 и \(2.7\,\mu\mathrm{M}\), а для силибина B --- 8.5 и \(5.0\,\mu\mathrm{M}\) соответственно для OATP1B1 и OATP1B3. Для исследованной смеси силимарина получены более низкие значения 1.3 и \(2.2\,\mu\mathrm{M}\). Перенос активности смеси на другой экстракт без аналитической характеристики состава некорректен. В клиническом исследовании у восьми мужчин силимарин 140 мг три раза в сутки в течение пяти дней не вызывал значимого изменения AUC розувастатина 10 мг. Отрицательный результат относится к этой схеме и не устанавливает безопасность высокодозовых или специально усиленных форм. [85,86]

В кинетической модели баланса уменьшение синусоидального захвата может снижать отношение несвязанных концентраций статина в гепатоците и плазме при одновременном уменьшении печеночного клиренса. Однако повышение плазменной концентрации увеличивает движущую силу остаточного захвата и пассивного поступления; одновременно могут изменяться внутриклеточный метаболизм и эффлюкс. Поэтому снижение внутрипеченочной экспозиции и ослабление гиполипидемического действия не являются неизбежными следствиями ингибирования OATP. Для силимарина сочетание доказанного снижения эффективности статина с повышенной миотоксичностью в рассмотренном клиническом исследовании не установлено. Такой «парадокс» допустим как проверяемая PK/PD-гипотеза, но не как документированный клинический исход. [31,43,85,86]

\paragraph{Цитрусовые: раздельная оценка гликозидов, агликонов и фуранокумаринов.}
Нарингин и нарингенин, а также гесперидин и гесперетин являются соответственно гликозидами и агликонами и не должны смешиваться при переносе экспериментальных значений. Для гесперидина и нарингина в сравнительном исследовании получены \(\mathrm{IC}_{50}\) 8.3 и \(6.9\,\mu\mathrm{M}\) в отношении OATP2B1; половинное ингибирование OATP1B1/OATP1B3 не достигалось при концентрациях до \(100\,\mu\mathrm{M}\). Нобилетин ингибировал OATP1B1 и OATP1B3 с \(\mathrm{IC}_{50}\) 2.1 и \(6.6\,\mu\mathrm{M}\). В отдельной системе BCRP-опосредованного эффлюкса дасатиниба \(\mathrm{IC}_{50}\) нобилетина и тангеретина составляли 1.04 и \(1.19\,\mu\mathrm{M}\). Это данные in vitro, а не клинически измеренная блокада BCRP после употребления апельсинового сока. Они не обосновывают приписывания гесперидину активности гесперетина или нобилетина и не позволяют классифицировать нарингенин как универсальный «умеренный» ингибитор OATP1B1/BCRP. [87,88]

Для грейпфрутового сока существенное значение имеет отдельная группа веществ --- фуранокумарины, влияющие на кишечный CYP3A4. Исследование с удалением фуранокумаринов из сока установило их причинный вклад во взаимодействие с фелодипином. В исследовании аторвастатина 40 мг с повторным приемом двойного по концентрации грейпфрутового сока AUC кислоты аторвастатина за 0--72 часа увеличивалась в 2.5 раза; основным объяснением авторы считали уменьшение CYP3A4-опосредованного пресистемного метаболизма. Это не доказывает клинического синергизма с одновременной блокадой печеночного OATP1B1. Подавление кишечного захвата, напротив, может уменьшать поступление других субстратов. Количественный результат зависит от препарата, вида и состава сока, концентрации и режима употребления; экстракты кожуры и обычный сок неэквивалентны. [87,89,90,96]

\paragraph{Кверцетин: транспортное ингибирование не определяет участок связывания.}
Кверцетин ингибировал транспорт SN-38 посредством ABCG2 дикого типа в мембранных везикулах с \(K_i=0.28\,\mu\mathrm{M}\). Однако уменьшение потока субстрата не доказывает связывания кверцетина непосредственно с нуклеотидсвязывающим доменом и подавления гидролиза АТФ. В кинетическом исследовании конформационных переходов ABCG2 кверцетин использовался как транспортируемый субстрат и ускорял образование ловушечных постгидролитических комплексов, что согласуется со стимуляцией АТФазного цикла. Конкуренция за транспорт и изменение скорости отдельных стадий могут сосуществовать. Поэтому утверждение о прямой блокаде гидролиза АТФ кверцетином требует отдельного доказательства и не используется как установленный механизм пищевого взаимодействия. [91,92]

\paragraph{Куркумин и ресвератрол: роль конъюгации и направленности транспорта.}
Интенсивная пресистемная конъюгация ограничивает экспозицию исходных куркумина и ресвератрола, но не позволяет считать их метаболиты фармакокинетически нейтральными. Ресвератрол-3-O-сульфат транспортировался OATP1B3, а ресвератрол-3-O-4\(^{\prime}\)-O-дисульфат --- OATP1B1 и OATP1B3. Для ресвератрол-4\(^{\prime}\)-O-сульфата и двух исследованных моноглюкуронидов транспорт этими OATP не выявлен. Сульфат куркумина являлся субстратом OATP1B1/OATP1B3/OATP2B1, тогда как транспорт глюкуронида куркумина не обнаружен. Эти результаты не обосновывают обобщение, что любые сульфаты и глюкурониды становятся субстратами OATP1B3. [93,94]

Свойства субстрата и ингибитора конъюгата необходимо оценивать раздельно. В эксперименте 2022 года ресвератрол и ресвератрол-3-сульфат ингибировали OATP1B3 с \(\mathrm{IC}_{50}\) 11.6 и \(11.8\,\mu\mathrm{M}\), тогда как для ресвератрол-3-глюкуронида \(\mathrm{IC}_{50}>25\,\mu\mathrm{M}\); последний отчетливее ингибировал OATP2B1. Следовательно, даже более высокая общая концентрация конъюгата в плазме не доказывает доминирования его вклада в ингибирование печеночного захвата: необходим учет связывания, химической формы, локальной экспозиции и субстрата-зонда. [97]

\paragraph{Конкурентный потенциал циркулирующих конъюгатов.}
Накопление конъюгата может уменьшать захват лекарственного субстрата через синусоидальную мембрану, если оба лиганда конкурируют за взаимно исключающие состояния переносчика и несвязанная концентрация конъюгата становится сопоставимой с его константой ингибирования. Это проверяемая кинетическая гипотеза, а не следствие одного лишь статуса субстрата. Изоформная принадлежность неодинакова: сульфат куркумина и дисульфат ресвератрола транспортируются OATP1B1 и OATP1B3, тогда как транспорт ресвератрол-3-O-сульфата в исследовании [93] выявлен только для OATP1B3. Для этих соединений нельзя априорно постулировать общий участок связывания со всеми статинами или клинически значимое конкурентное вытеснение. Раздельные измерения транспорта и ингибирования необходимы также потому, что \(K_m\) транспортируемого конъюгата не является автоматически его \(K_i\) для другого субстрата. [93,94,97]

Для каждого метаболита \(m\) следует независимо от агликона определять \(K_{i,mj}\), механизм ингибирования изоформы \(j\) и несвязанную экспозицию. Соотношение \(I_{u,m}(t)=f_{u,p,m}C_{p,m}(t)\) описывает несвязанную концентрацию в исследованной плазме; ее перенос на синусоидальную поверхность требует оценки концентрационного градиента, образования и удаления метаболита. Высокая полярность конъюгата не доказывает высокого \(f_{u,p,m}\): ресвератрол-3-сульфат образовывал с альбумином более стабильный комплекс, чем исходный ресвератрол. Поэтому \(f_{u,p,m}\) необходимо измерять, а не приравнивать единице. В опытах ингибирования отдельно учитывают неспецифическое связывание в инкубационной среде, стабильность конъюгата и возможное образование агликона; иначе наблюдаемая активность может быть ошибочно приписана исходно внесенному метаболиту. [1,97]

\paragraph{Внутриклеточное образование и локальная экспозиция метаболитов.}
Для конъюгатов, образующихся в энтероцитах или гепатоцитах, периферическая \(C_{p,m}(t)\) не определяет однозначно внутриклеточный пул. Она не является гарантированной нижней границей локальной экспозиции: соотношение зависит от образования, захвата, связывания, деконъюгации и направленного эффлюкса. В клетках HeLa с экспрессией UGT1A1 подавление MRP4 уменьшало выведение моноглюкуронидов ресвератрола и увеличивало их внутриклеточное накопление; одновременно выявлялся гидролиз конъюгатов обратно до агликона. Это экспериментальное основание для раздельного учета образования и удаления метаболита, но не количественная характеристика его экспозиции в гепатоцитах человека in vivo. [99]

В иллюстративной модели хорошо перемешанного клеточного компартмента при постоянном объеме \(V_{\mathrm{cell}}\), линейных потоках и быстром равновесии связывания баланс можно записать как
\[
\begin{aligned}
\frac{dA_{m,\mathrm{cell}}}{dt}
&=v_{\mathrm{form},m}+CL_{\mathrm{in},u,m}C_{p,u,m}\\
&\quad-\left(CL_{\mathrm{efflux},u,m}+CL_{\mathrm{loss},u,m}\right)
I_{u,m,\mathrm{int}},\\
I_{u,m,\mathrm{int}}
&=f_{u,\mathrm{cell},m}\frac{A_{m,\mathrm{cell}}}{V_{\mathrm{cell}}}.
\end{aligned}
\]
Здесь \(A_{m,\mathrm{cell}}\) --- количество метаболита, \(v_{\mathrm{form},m}\) --- скорость образования в единицах количества вещества за время, а не вектор концентрации; все \(CL\) имеют размерность объема за время. \(C_{p,u,m}\) относится к прилежащей внеклеточной плазме, а \(CL_{\mathrm{loss},u,m}\) объединяет дальнейшую биотрансформацию и деконъюгацию; эффлюкс включает раздельно параметризуемые апикальный и базолатеральный пути. При постоянных входных потоках и положительном суммарном клиренсе стационарное решение имеет вид
\[
I_{u,m,\mathrm{int,ss}}=
\frac{v_{\mathrm{form},m}+CL_{\mathrm{in},u,m}C_{p,u,m}}
{CL_{\mathrm{efflux},u,m}+CL_{\mathrm{loss},u,m}}.
\]
Запись \(I_{u,m,\mathrm{int,ss}}=f_{u,\mathrm{cell},m}v_{\mathrm{form},m}/CL_{\mathrm{efflux},\mathrm{tot},m}\) допустима лишь при отсутствии внешнего притока и других потерь, причем \(CL_{\mathrm{efflux},\mathrm{tot},m}\) должен быть определен относительно общей клеточной концентрации. Поскольку \(CL_{\mathrm{efflux},\mathrm{tot},m}=f_{u,\mathrm{cell},m}CL_{\mathrm{efflux},u,m}\), дополнительное умножение на \(f_{u,\mathrm{cell},m}\) при использовании клиренса по несвязанной концентрации означало бы повторный учет связывания. При насыщении эффлюкса, изменяющемся образовании или субклеточной компартментализации требуется динамическая, а не стационарная модель. [31,99]

Внутриклеточная концентрация также не равна концентрации у наружного участка связывания OATP: для оценки ингибирования необходимо указать сторону доступа и состояние переносчика. Внутренний пул сульфата куркумина или моноглюкуронидов ресвератрола нельзя автоматически подставлять в модель конкурентной блокады синусоидального захвата. Низкая периферическая экспозиция может приводить к недооценке локального отношения \(I_u/K_i\), но не к «ложноотрицательному \(K_i\)»: константа ингибирования определяется отдельным экспериментом. Статус субстрата и ингибитора каждой изоформы сохраняет ограничения, изложенные выше. [6,93,94,97,99]

\paragraph{Сторона доступа и условная модель транс-ингибирования.}
Асимметрия ингибирования с внешней и внутренней стороны клетки имеет экспериментальную основу: для циклоспорина A данные клеток HEK293 с экспрессией OATP1B1 и кинетическое моделирование поддерживают сочетание внутриклеточного накопления, cis- и trans-ингибирования с неодинаковыми константами. Это основание для проверки соответствующего механизма у конъюгатов, но не доказательство их связывания с определенным «цитоплазматическим доменом» OATP1B1/1B3. [103]

Если независимо подтверждено снижение входного потока под действием внутриклеточного метаболита, для одной изоформы \(j\) допустима условная факторизованная модель
\[
\begin{aligned}
CL_{\mathrm{in},j}^{\mathrm{cond}}
&=CL_{\mathrm{in},j}^{\mathrm{cis}}\,T_j,\\
T_j&=\left(1+\sum_m
\frac{I_{u,m,\mathrm{int}}}{K_{i,\mathrm{trans},mj}}
\right)^{-1}.
\end{aligned}
\]
Здесь \(CL_{\mathrm{in},j}^{\mathrm{cis}}\) уже учитывает внешнее ингибирование, а \(K_{i,\mathrm{trans},mj}\) параметризуется отдельно по эффекту внутренней несвязанной экспозиции на транспорт определенного зонда. Для единственного метаболита возникает предложенный множитель \((1+I_{u,m,\mathrm{int}}/K_{i,\mathrm{trans},mj})^{-1}\). Факторизация означает дополнительно принятое независимое уменьшение доступной транспортной активности; она не следует только из нахождения лигандов по разные стороны мембраны и не устанавливает аллостерический механизм. Если внутреннее и внешнее связывание сопряжены конформационным циклом, необходима общая схема состояний переносчика, а не механическое перемножение факторов. Сумма по \(m\) также предполагает взаимно исключающее конкурентное действие внутренних лигандов. [6,103]

Для оценки параметров нужны контролируемая загрузка метаболита, измерение его несвязанного клеточного пула и внешней концентрации во времени, а также независимая оценка cis-ингибирования. Иначе выход метаболита из клетки с последующей внешней блокадой неотличим от предполагаемого trans-эффекта. Параметр внутреннего ингибирования не следует считать стандартным внешним конкурентным \(K_i\); для рассматриваемых конъюгатов его величина и сама применимость модели требуют проверки. Измерение только общей клеточной концентрации или одной конечной точки AUC не идентифицирует эти параметры. [93,94,97,103]

Для куркумина имеются прямые клинические данные о взаимодействии с сульфасалазином. У восьми здоровых участников прием куркумина 2 г за 30 минут до сульфасалазина 2 г повышал AUC за 0--24 часа и \(C_{\max}\). По отношениям арифметических средних из таблицы 1 первичной публикации \(R_{\mathrm{AUC}}=19.0/5.89\approx3.23\), а \(R_{C_{\max}}=24.5/6.58\approx3.72\); это не отношения геометрических средних. В текстовом описании результатов порядок коэффициентов для микродозы и терапевтической дозы не согласуется с таблицей; здесь использованы табличные данные. В сопутствующем везикулярном опыте \(\mathrm{IC}_{50}=1.60\pm0.93\,\mu\mathrm{M}\), а \(K_i=0.70\pm0.41\,\mu\mathrm{M}\) рассчитан с использованием ранее определенного \(K_m\). Низкие плазменные концентрации куркумина у участников согласуются с преимущественно кишечным действием, но не доказывают абсолютной селективности к BCRP. Величины этого взаимодействия нельзя переносить на розувастатин, метотрексат или НОАК. [95]

\paragraph{Суммарный эффект транспортной сети.}
При одинаковой пероральной дозе и линейной фармакокинетике исходное соотношение имеет вид
\[
R_{\mathrm{AUC}}=
\frac{F_{\mathrm{with}}}{F_{\mathrm{without}}}
\frac{CL_{\mathrm{without}}}{CL_{\mathrm{with}}}.
\]
Ингибирование кишечного эффлюкса может повышать \(F\), ингибирование кишечного захвата --- снижать \(F\), а подавление печеночного захвата --- изменять \(CL\) и распределение в органе-мишени. Эти эффекты могут усиливать или частично компенсировать друг друга. Данное уравнение является модельным балансом и не идентифицирует механизм по одному отношению AUC. Из него также нельзя вывести \(R_{C_{\max}}\) без учета скорости всасывания и распределения. Нутрицевтическое взаимодействие следует определять для конкретной пары «характеризованный продукт--субстрат», а не для абстрактного класса «полифенолы». [1,31,81,82,86,95]

\section{Фармакокинетические последствия: изменение AUC и экспозиции в органах}

\subsection{Разграничение биодоступности при приеме внутрь и системного клиренса}

При линейной, стационарной фармакокинетике после приема внутрь дозы \(D\)
\[
\mathrm{AUC}_{0-\infty}=\frac{F D}{CL},
\qquad
R_{\mathrm{AUC}}=
\frac{F_{\mathrm{with}}}{F_{\mathrm{without}}}
\frac{CL_{\mathrm{without}}}{CL_{\mathrm{with}}},
\]
если дозы одинаковы. Здесь \(F\) --- абсолютная биодоступность, а \(CL\) --- системный клиренс. Та же зависимость «доза/клиренс» применима к AUC за интервал дозирования в равновесном состоянии. Биодоступность при приеме внутрь можно представить как \(F=F_aF_gF_h\), где множители характеризуют всасывание и доли вещества, избежавшие кишечных и печеночных потерь при первом прохождении. Ингибирование кишечного BCRP может изменять эффективное всасывание, тогда как ингибирование OATP способно изменять печеночную экстракцию и последующий системный клиренс. Их вклады нельзя определить по одной AUC. Модели PBPK позволяют представить эти процессы раздельно, но для определения их величин требуется независимая информация. [1,31]

В качестве пояснительного расчета, не отражающего наблюдавшееся взаимодействие, увеличение биодоступности в 1.5 раза при снижении клиренса до 40\% исходного дает \(R_{\mathrm{AUC}}=1.5/0.40=3.75\). Умножение следует из уравнения и не доказывает молекулярной синергии. Напротив, если два ингибитора действуют на один и тот же насыщаемый путь захвата, перемножение их отдельно измеренных отношений AUC может существенно завышать прогноз совместного эффекта.

Упрощенная модель долевого вклада в клиренс позволяет оценить предельное увеличение экспозиции:
\[
R_{\mathrm{AUC}}\approx
\left[(1-f_T)+\frac{f_T}{1+I_u/K_i}\right]^{-1},
\]
где \(f_T\) --- исходная доля общего клиренса, приходящаяся на ингибируемый путь, при допущении неизменной биодоступности и отсутствия воздействия на параллельные пути. При полном ингибировании предельное отношение составляет \(1/(1-f_T)\). Следовательно, \(f_T=0.5\) дает двукратный предел, а \(f_T=0.8\) --- пятикратный. Это модельные примеры, а не измеренные доли транспортерного клиренса конкретного статина. Поскольку захват может определять доступность нескольких последующих путей элиминации, эффективный вклад транспортера не обязательно равен независимо суммируемой доле билиарного или метаболического клиренса. [31]

\subsection{Максимальная концентрация и экспозиция несвязанной фракции}

Максимальная концентрация зависит от скорости всасывания, биодоступности, распределения, клиренса и времени приема препарата. Для однокамерной модели перорального введения со всасыванием первого порядка
\[
C(t)=\frac{F D k_a}{V(k_a-k_e)}
\left(e^{-k_e t}-e^{-k_a t}\right),
\qquad
t_{\max}=\frac{\ln(k_a/k_e)}{k_a-k_e}.
\]
Здесь \(k_a\) и \(k_e\) обозначают константы скорости всасывания и элиминации, а \(V\) --- кажущийся объем распределения. Эта пояснительная модель показывает, почему не существует универсального алгебраического пересчета между отношениями AUC и максимальной концентрации. Ускорение поступления всосавшегося вещества в системный кровоток или уменьшение ранней печеночной экстракции может непропорционально увеличивать максимальную концентрацию. Сопоставление отношений максимальной концентрации и AUC 5.62 и 2.15 при глекапревире/пибрентасвире, а также 18.88 и 7.39 при схеме с дополнительным воксилапревиром иллюстрирует это различие, не позволяя однозначно определить лежащие в его основе константы скорости. [14,17]

Если связывание с белками остается линейным и доля несвязанного вещества постоянна в течение рассматриваемого интервала, плазменная экспозиция несвязанной фракции удовлетворяет соотношению \(\mathrm{AUC}_u=f_u\mathrm{AUC}_{\mathrm{total}}\), где \(f_u\) --- доля вещества, не связанного с белками плазмы, а индекс \(\mathrm{total}\) указывает, что AUC рассчитана по суммарной концентрации связанного и несвязанного вещества. При изменении \(f_u\) изменение общей AUC не обязательно отражает изменение экспозиции фармакологически доступной фракции. Концентрация альбумина, вытеснение из связи с белками и нарушения функции органов могут усложнять интерпретацию. Тем не менее наблюдаемое увеличение общей экспозиции нельзя объяснять безвредным вытеснением из связи с белками без данных о концентрации несвязанной фракции и клиренсе. Для ингибиторов с высокой степенью связывания экспозиция несвязанной фракции в портальной крови или внутри клетки особенно важна при интерпретации выраженного ингибирования, наблюдаемого в клеточных системах. [1,25,31]

\subsection{Соотношение печеночной и плазменной экспозиции}

Упрощенная модель баланса вещества в гепатоците имеет вид
\[
\frac{dA_{\mathrm{hep}}}{dt}
=PS_{\mathrm{in}}C_{u,p}
-\left(PS_{\mathrm{out}}+CL_{\mathrm{met}}+CL_{\mathrm{bile}}\right)C_{u,\mathrm{hep}},
\]
а при допущении квазистационарного состояния
\[
K_{p,uu,\mathrm{hep}}=
\frac{C_{u,\mathrm{hep}}}{C_{u,p}}
=\frac{PS_{\mathrm{in}}}
{PS_{\mathrm{out}}+CL_{\mathrm{met}}+CL_{\mathrm{bile}}}.
\]
\(A_{\mathrm{hep}}\) обозначает количество вещества в гепатоците, а \(C_{u,p}\) и \(C_{u,\mathrm{hep}}\) --- концентрации несвязанной фракции в плазме и гепатоците соответственно. \(PS_{\mathrm{in}}\) и \(PS_{\mathrm{out}}\) обозначают клиренсы поступления в клетку и выхода из нее через синусоидальную мембрану; \(CL_{\mathrm{met}}\) и \(CL_{\mathrm{bile}}\) --- клиренсы внутриклеточного метаболизма и выведения в желчь. Схема не учитывает ограничение перфузией, зональную неоднородность, кинетику связывания и концентрационно-зависимый транспорт. Она приведена для демонстрации направленности зависимостей, а не для оценки концентраций в печени человека. [31]

Позитронно-эмиссионная томография с радиоактивно меченным розувастатином у человека использовалась для изучения печеночных концентраций и гепатобилиарного транспорта в отсутствие и в присутствии циклоспорина. Такие исследования позволяют ограничить диапазоны параметров, неидентифицируемых по одним лишь плазменным данным захвата и билиарного клиренса, однако условия и путь введения радиотрейсера необходимо учитывать при сопоставлении с обычными исследованиями взаимодействий после приема внутрь. Они не дают универсального пересчета плазменной AUC в экспозицию несвязанного вещества в печени. [43]

Уменьшение поступления может снижать отношение концентраций несвязанной фракции «печень/плазма» при одновременном увеличении плазменной AUC. Напротив, уменьшение канальцевого выведения может увеличивать внутриклеточное удержание. Если меняются оба процесса, итоговая печеночная экспозиция зависит от относительной величины этих изменений и от изменения поступления из плазмы. Таким образом, более высокая плазменная концентрация статина не гарантирует пропорционально более выраженного печеночного эффекта по снижению липопротеинов низкой плотности (ЛПНП). Менее выраженное снижение уровня ЛПНП, наблюдавшийся в исследовании взаимодействия с лопинавиром/ритонавиром, согласуется с этим предположением, однако изменения липидного обмена и дизайн исследования не позволяют считать его доказательством сниженной концентрации в гепатоцитах. [27,31]

\subsection{Увеличение экспозиции, повреждение мышц и почечные исходы}

Формулировка «увеличение в два--семь раз» описывает важный эмпирический диапазон, а не валидированную границу между терапевтическими и токсическими концентрациями. Клиническое значение отношений AUC определяется исходной дозой, абсолютной концентрацией, длительностью экспозиции, восприимчивостью пациента и рассматриваемой молекулярной формой. По этим исследованиям нельзя установить универсальный плазменный порог розувастатина, при котором развивается рабдомиолиз. Небольшие исследования с участием здоровых добровольцев также не позволяют оценить частоту редкой тяжелой токсичности у пожилых пациентов с множественными заболеваниями. [25--29,34]

\begin{itemize}
\item \textbf{Скелетная мускулатура.} Повышенная системная экспозиция статина может увеличивать его поступление во внепеченочные ткани и риск миотоксичности. Распределение и метаболизм препарата в мышечной ткани, нарушения мевалонатного пути, функция митохондрий и индивидуальная восприимчивость усложняют простую зависимость между экспозицией и повреждением. Выраженное повышение активности креатинкиназы, наблюдавшееся при сочетании с лопинавиром/ритонавиром, служит сигналом возможной миотоксичности, а не оценкой частоты рабдомиолиза. [27,34]
\item \textbf{Повреждение почек вследствие рабдомиолиза.} Тяжелое повреждение мышц может вызвать миоглобин-ассоциированное острое повреждение почек. Почечная дисфункция затем способна дополнительно нарушать элиминацию препарата-субстрата, формируя клинически значимую петлю обратной связи. Циклоспорин или другой сопутствующий нефротоксичный препарат может независимо повреждать почки; поэтому отнесение каждого нежелательного почечного исхода к ингибированию BCRP механистически необоснованно. [3,25]
\item \textbf{Протеинурия или гематурия без повреждения мышц.} В крупном наблюдательном сравнении применение розувастатина ассоциировалось с умеренно более высокими рисками гематурии, протеинурии и почечной недостаточности, требующей заместительной терапии, по сравнению с аторвастатином. Отношения рисков составляли соответственно 1.08, 1.17 и 1.15. Исследование обосновывает осторожность при высоких дозах и тяжелой болезни почек, но не устанавливает, что эти исходы вызваны ингибированием OATP/BCRP, и не определяет двукратное увеличение AUC как порог «бессимптомной нефропатии». [32]
\item \textbf{Секреция креатинина и фильтрация.} Кобицистат может повышать концентрацию креатинина в сыворотке за счет ингибирования канальцевой секреции без эквивалентного снижения измеренной скорости клубочковой фильтрации. Исследование клиренса йогексола непосредственно разграничило эти процессы. Поэтому изолированное повышение креатинина следует интерпретировать с учетом времени его появления, анализа мочи, сопутствующих нефротоксичных воздействий и, при клинической необходимости, альтернативной оценки фильтрации; прогрессирующее повреждение нельзя без дополнительной оценки объяснять доброкачественным эффектом усилителя. [33]
\end{itemize}

Для метотрексата и противоопухолевых субстратов удержание в тканях может быть одновременно фармакологически полезным и токсичным. Отсутствие валидированного у человека отношения экспозиции, селективного для конкретного транспортера, является основанием для субстрат-специфического мониторинга и специализированной оценки, а не основанием переносить доли снижения доз статинов на другие лекарственные средства. Клеточную кинетику эффлюкса метотрексата нельзя напрямую преобразовать в схему терапии спасения в онкологии или в прогноз вероятности повреждения почек. [9,11,24]

\section{Снижение клинического риска и регуляторные ограничения дозирования: документы 2024--2026 годов}

\subsection{Временные рамки и статус регуляторных документов}

В данном разделе сопоставлены документы 2024--2026 годов; речь идет о синтезе данных за указанный период, а не о едином руководстве FDA--EMA--NIH. FDA выпустило окончательное руководство ICH M12 в августе 2024 года. EMA указывает, что ICH M12 вступило в силу в Европейском союзе 30 ноября 2024 года, заменив прежнее руководство по взаимодействиям; впоследствии были опубликованы материалы по внедрению. Скрининговые пороги на этапе разработки препаратов гармонизированы в большей степени, чем ограничения назначения отдельных лекарств. Лечебные рекомендации NIH, генотип-ориентированные рекомендации CPIC и официально утвержденная информация о препаратах имеют разные цели и не должны смешиваться. [1,2,4,34]

27 мая 2026 года NIH объявили о прекращении использования отдельного раздела по статинам для первичной профилактики и переходе к междисциплинарному руководству ACC/AHA и других профессиональных обществ по дислипидемии, опубликованному в марте 2026 года. NIH продолжили направлять клиницистов к своим таблицам межлекарственных взаимодействий для выбора и контроля сочетаний статинов с антиретровирусными препаратами. Следовательно, обзор с датой завершения в сентябре 2026 года не должен представлять раздел по профилактическому применению статинов от февраля 2024 года как неизменный действующий подход. Приведенные далее ограничения, специфичные для взаимодействий, по-прежнему опираются на доступные таблицы NIH и актуальную информацию о препаратах. [5]

Дата 30 ноября 2024 года относится к вступлению ICH M12 в силу в Европейском союзе, а не к первоначальному выпуску окончательного руководства FDA, состоявшемуся в августе 2024 года. Стратегия внедрения EMA уточняет, что отдельные темы, не полностью охваченные M12, продолжают регулироваться применимыми положениями прежнего руководства до замены специализированными документами. Это особенно важно при рассмотрении ИПП: изменение желудочно-кишечной растворимости или моторики нельзя подменять анализом ингибирования транспортеров. [1,2]

\subsection{Скрининговые критерии ICH M12 --- не клинические пороги токсичности}

\begin{itemize}
\item \textbf{OATP1B1/OATP1B3.} Базовый скрининг не позволяет исключить клиническое ингибирование при \(C_{\max,\mathrm{inlet},u}/\mathrm{IC}_{50,u}\geq0.1\). В прежней записи через отношение этому соответствует \(R=1+C_{\max,\mathrm{inlet},u}/\mathrm{IC}_{50,u}\geq1.1\). Параметр экспозиции здесь --- расчетная концентрация несвязанной фракции на входе в печень, а не просто общая максимальная концентрация в периферической плазме. [1]
\item \textbf{Кишечные BCRP/P-gp после приема внутрь.} Базовый критерий составляет
\[
(D/250\,\mathrm{mL})/\mathrm{IC}_{50,u}\geq10
\]
при использовании согласованных молярных единиц. Отношение дозы к объему является консервативным суррогатом кишечной концентрации, а не измеренной концентрацией в энтероците. [1]
\item \textbf{Парентеральные ингибиторы и метаболиты, образующиеся после всасывания.} M12 описывает альтернативный скрининг P-gp/BCRP с использованием \(C_{\max,u}/\mathrm{IC}_{50,u}\geq0.02\), когда кишечный критерий на основе дозы неприменим. Эти пороги указывают на необходимость дополнительной оценки; они не предписывают фиксированного снижения дозы и не классифицируют препарат как клинически «сильный» ингибитор. [1]
\end{itemize}

Для оценки концентрации на входе в печень M12 использует
\[
C_{\max,\mathrm{inlet},u}
=f_{u,p}\left(C_{\max}+\frac{F_a F_g k_a D}{Q_h R_B}\right).
\]
\begin{itemize}
\item \textbf{Параметры экспозиции на входе в печень.} \(Q_h\) обозначает печеночный кровоток, \(R_B\) --- отношение концентраций в цельной крови и плазме, \(k_a\) --- константу скорости всасывания. Дозу и концентрации необходимо выражать в согласованных единицах. При отсутствии данных консервативная оценка допускает \(F_a=F_g=1\) и \(k_a=0.1\,\mathrm{min}^{-1}\). Это расчетные предпосылки скрининга, а не измеренные свойства каждого ингибитора. [1]
\item \textbf{Нижняя граница учета несвязанной фракции.} Если достоверность измерения \(f_{u,p}<0.01\) не подтверждена, M12 предписывает использовать \(f_{u,p}=0.01\). Такой подход ограничивает риск недооценки взаимодействия препарата с высокой степенью связывания с белками; он не означает, что его истинная несвязанная фракция обязательно составляет 1\%. [1]
\end{itemize}

Прогнозы PBPK, механистические статические модели, эндогенные биомаркеры или клиническое исследование взаимодействия могут уточнить положительный результат скрининга. Концентрация ингибитора и \(\mathrm{IC}_{50}\) должны определяться с использованием совместимых подходов к учету связывания. Проверка пригодности модели требует большего, чем подгонка к одному значению AUC розувастатина: сравнения с использованием различных зондов и ингибиторов должны проверять предполагаемый вклад отдельных путей. Копропорфирин I может дополнять оценку OATP, но не измеряет ингибирование BCRP и не подтверждает безопасность конкретной дозы статина. [1,31,37]

Проверку CP-I целесообразно дополнять CP-III как вторым эндогенным показателем, но не как независимым селективным измерителем OATP1B3. Для интерпретации необходимы стандартизация отбора, обработки и хранения образцов, исходные значения, генотип \textit{SLCO1B1} и согласованная модель нескольких путей транспорта. Различия ответов двух изомеров сами по себе не позволяют решить обратную задачу и единственным образом восстановить активность OATP1B1 и OATP1B3. [52,53]

Опубликованный в 2026 году анализ данных 17 компаний включал 58 клинических исследований ингибирования транспортеров с 42 взаимодействующими соединениями. В большинстве исследований увеличение экспозиции было менее двукратным; более выраженные взаимодействия чаще сопровождали вмешательство в несколько путей, а статические регуляторные критерии характеризовались низкой частотой ложноотрицательных и высокой частотой ложноположительных результатов для P-gp и OATP. Эти результаты поддерживают использование скрининга как консервативного этапа отбора, а не как классификации клинической силы ингибитора. Выборка исследований разработчиков не позволяет оценить частоту взаимодействий во всей кардиометаболической популяции или отменить ограничения для уже известных опасных сочетаний. [76]

\paragraph{Перенос принципов ICH M12 на пищевые смеси.}
Применение скрининговых критериев к БАДам требует компонентного анализа: дозу отдельного соединения нельзя заменять массой всей растительной смеси, а \(\mathrm{IC}_{50}\) одного агликона --- константой для всех его конъюгатов. Для OATP1B существенна несвязанная экспозиция на входе в печень; для кишечного BCRP расчет «доза/250 мл» является консервативным скринингом, а не измеренной концентрацией у внутриклеточного участка связывания. Положительный скрининг обосновывает дальнейшее исследование, но не присвоение экстракту категории клинически «сильного» ингибитора и не прогноз определенного увеличения риска кровотечения или рабдомиолиза. ICH M12 не валидирует нутрицевтическую феноконверсию как диагностическую категорию и не заменяет оценку конкретного продукта и субстрата. [1,79,85,87,95,97]

\paragraph{Расчет кишечного индекса для куркумина.}
В качестве воспроизводимого арифметического примера используем дозу самого куркумина \(D_{\mathrm{mass}}=2\,\mathrm{g}\), соответствующую исследованию [95], а не массу произвольного экстракта куркумы. При молярной массе \(M=368.38\,\mathrm{g\,mol}^{-1}\) номинальная концентрация для условного объема \(0.250\,\mathrm{L}\) составляет [98]
\[
\begin{aligned}
D_{\mathrm{mol}}&=\frac{D_{\mathrm{mass}}}{M}
=\frac{2}{368.38}\,\mathrm{mol}\approx5.43\,\mathrm{mmol},\\
I_{\mathrm{gut,nom}}&=\frac{D_{\mathrm{mol}}}{0.250\,\mathrm{L}}
\approx21.7\,\mathrm{mM}=2.17\times10^{4}\,\mu\mathrm{M},\\
J_{\mathrm{gut,nom}}&=\frac{I_{\mathrm{gut,nom}}}{\mathrm{IC}_{50}}
\approx\frac{2.17\times10^{4}}{1.60}\approx1.36\times10^{4}.
\end{aligned}
\]
Здесь \(J_{\mathrm{gut,nom}}\) --- безразмерный индекс, рассчитанный с экспериментальной везикулярной \(\mathrm{IC}_{50}=1.60\,\mu\mathrm{M}\); он примерно в \(1.36\times10^{3}\) раз превышает скрининговую границу 10. Строгое сопоставление с ICH M12 требует \(\mathrm{IC}_{50,u}\); подстановка номинальной \(\mathrm{IC}_{50}\) предполагает пренебрежимо малое неспецифическое связывание в опыте. Если это допущение не подтверждено, приведенное число остается номинальной оценкой, а не измерением индекса на основе несвязанных концентраций. [1,95]

Граница 10 является критерием, при превышении которого базовый скрининг не позволяет исключить взаимодействие, а не порогом клинической безопасности. Расчет не учитывает растворение, преципитацию, деградацию, время кишечного транзита и доступ ингибитора к мембране; он не доказывает неизбежности взаимодействия и не прогнозирует его величину. Наблюдаемое \(R_{C_{\max}}=24.5/6.58\approx3.72\) для сульфасалазина подтверждается отдельными клиническими данными таблицы 1 [95], а не выводится из \(J_{\mathrm{gut,nom}}\). Низкая плазменная экспозиция куркумина совместима с локальным кишечным воздействием, но ни она, ни высокий индекс сами по себе не устанавливают селективности механизма относительно BCRP. [1,95]

\paragraph{Растворимость: отдельный сценарий, а не замена критерия M12.}
Номинальный расчет «доза/250 мл» является консервативным суррогатом, а не строгой физической верхней границей концентрации во всех отделах кишечника. ICH M12 не предписывает заменять его минимумом из дозовой концентрации и растворимости. Дополнительная равновесная модель однородной среды может использовать
\[
I_{\mathrm{gut,eff}}=
\min\!\left(I_{\mathrm{gut,nom}},C_{\mathrm{sat,gut}}\right),
\qquad
J_{\mathrm{gut,eff}}=
\frac{I_{\mathrm{gut,eff}}}{\mathrm{IC}_{50,u}},
\]
если \(C_{\mathrm{sat,gut}}\) экспериментально определена для соответствующих pH, температуры, среды и твердой фазы, а растворенная концентрация сопоставима с несвязанной формой в опыте ингибирования. Это дополнительный анализ чувствительности; перенос на него порога 10 не превращает новую метрику в валидированный регуляторный критерий. [1]

Термодинамическую растворимость свободного агликона необходимо отличать от кажущейся растворимости в мицеллах и временной пересыщенности. Для куркумина экспериментально показаны зависимость стабильности от среды и необходимость различать преципитацию и химическую деградацию. Аморфный комплекс куркумина с хитозаном в присутствии гидроксипропилметилцеллюлозы поддерживал пересыщенность примерно в пять раз выше насыщенной растворимости в течение трех часов in vitro. Поэтому простое отсечение по равновесной растворимости не является универсальным описанием лекарственных форм и пищевых смесей. [101,102]

Если исключительно для условного сценария принять \(C_{\mathrm{sat,gut}}=20\,\mu\mathrm{M}\) и \(\mathrm{IC}_{50,u}\approx1.60\,\mu\mathrm{M}\), то
\[
J_{\mathrm{gut,eff}}=\frac{20}{1.60}=12.5.
\]
Число \(20\,\mu\mathrm{M}\) здесь не представлено как установленная универсальная термодинамическая растворимость куркумина в ЖКТ; соответствующего измерения для продукта и среды исследования [95] в рассмотренных источниках не установлено. При тех же допущениях граница \(J=10\) соответствует \(I_{\mathrm{gut,eff}}=16\,\mu\mathrm{M}\), поэтому вывод чувствителен к растворимости, связыванию и неопределенности \(\mathrm{IC}_{50}\). Значение 12.5 не исключает взаимодействия в данном сценарии, но не доказывает ни его неизбежности, ни специфичности относительно BCRP. Измеренное изменение \(C_{\max}\) и механистическая атрибуция требуют самостоятельных данных; итоговая модель должна воспроизводить клинический результат без подбора растворимости только ради пересечения порога. [1,95,101,102]

\paragraph{Временная пересыщенность и кинетика преципитации.}
При неизменных pH, температуре, среде и твердой фазе \(C_{\mathrm{sat,gut}}\) остается равновесной характеристикой; релаксация пересыщенного раствора изменяет фактическую растворенную концентрацию, а не саму термодинамическую растворимость. В качестве частной кинетической аппроксимации при постоянном объеме и отсутствии продолжающегося растворения, всасывания, деградации и транзитного оттока можно записать
\[
\begin{aligned}
\frac{dC_{\mathrm{diss}}}{dt}
&=-k_{\mathrm{prec}}\left(C_{\mathrm{diss}}-C_{\mathrm{sat,gut}}\right),\\
C_{\mathrm{diss}}(t)
&=C_{\mathrm{sat,gut}}
\left[1+(\beta-1)e^{-k_{\mathrm{prec}}t}\right],\\
\beta&=\frac{C_{\mathrm{diss}}(0)}{C_{\mathrm{sat,gut}}}\geq1,
\qquad k_{\mathrm{prec}}\geq0.
\end{aligned}
\]
\(\beta\) обозначает начальную кратность пересыщения, а не долю аморфного вещества или универсальный коэффициент аморфизации; \(k_{\mathrm{prec}}\) имеет размерность обратного времени. При \(\beta=1\) пересыщения нет, а при \(k_{\mathrm{prec}}>0\) избыток экспоненциально исчезает. Для \(k_{\mathrm{prec}}=0\) модель сохраняет начальное пересыщение. Индукционный период нуклеации, нелинейный рост частиц и влияние матрицы требуют иной кинетики; наличие пересыщения в исследовании [102] не валидирует именно экспоненциальную форму для клинического продукта. [101,102]

В более общем балансе растворенного вещества следует учитывать
\[
\frac{d(V_{\mathrm{gut}}C_{\mathrm{diss}})}{dt}
=R_{\mathrm{diss}}-R_{\mathrm{prec}}-R_{\mathrm{abs}}
-R_{\mathrm{deg}}-R_{\mathrm{out}},
\]
где все \(R\) имеют размерность количества вещества за время, а \(R_{\mathrm{prec}}\geq0\) должен обращаться в нуль при отсутствии пересыщения. Начальное количество \(V_{\mathrm{gut}}\beta C_{\mathrm{sat,gut}}\) не может превышать доступное количество вещества в выбранном компартменте. При изменении условий вдоль кишечника растворимость также может зависеть от времени через эти условия, но это отдельный эффект, не тождественный преципитации.

При согласованном учете связывания можно рассчитать дополнительный временной показатель \(J_{\mathrm{gut,eff}}(t)=I_{u,\mathrm{gut}}(t)/\mathrm{IC}_{50,u}\), не подменяя им стандартный индекс ICH M12. В частном сценарии \(I_{u,\mathrm{gut}}=C_{\mathrm{diss}}\) и постоянной \(\mathrm{IC}_{50,u}\) начальное значение в \(\beta\) раз выше равновесного. Это допускает исследование ранней локальной экспозиции ингибитора, но не доказывает «внутрикишечной перегрузки плазмы»: плазменный профиль субстрата требует отдельного сопряженного баланса абсорбции, распределения и элиминации. Формула не гарантирует отсутствия недооценки \(J\) или воспроизведения клинического \(C_{\max}\) без независимой оценки параметров. [1,95,101,102]

\subsection{Ограничения доз и недопустимые сочетания в конкретных схемах}

Далее приведены ограничения совместного применения у взрослых, а не рекомендации по начальной дозе для каждого пациента. Выбор малой дозы не отменяет противопоказание или рекомендацию избегать сочетания. При одновременном наличии нескольких ограничений действует наиболее строгое применимое правило; если в его пределах невозможно достичь целевого гиполипидемического эффекта, следует изменить схему, а не превышать предел дозы. [3--5,14--18]

\begin{itemize}
\item \textbf{Циклоспорин и розувастатин.} Американская таблица коррекции доз ограничивает розувастатин дозой 5 мг один раз в сутки. Этот предел не означает предпочтительности комбинации: инструкция также содержит предостережение относительно совместного применения. В национальной общей характеристике препарата (SmPC) Rosuvastatin Viatris, опубликованной ирландским регулятором HPRA 21 марта 2025 года, сочетание с циклоспорином прямо противопоказано. Американский предел нельзя переносить на эту европейскую инструкцию; историческое обсуждение EMA не заменяет действующую национальную SmPC. [3,40,45]
\item \textbf{Глекапревир/пибрентасвир и розувастатин.} Американская инструкция допускает не более 10 мг в сутки; информация EMA о Maviret ограничивает розувастатин дозой 5 мг в сутки. Американская инструкция требует снижения дозы правастатина на 50\%, тогда как европейский предел правастатина составляет 20 мг в сутки. В США аторвастатин, ловастатин и симвастатин не рекомендуются совместно с Mavyret; европейские ограничения для аторвастатина и симвастатина являются противопоказаниями. [14,15]
\item \textbf{Софосбувир/велпатасвир и розувастатин.} Доза розувастатина не должна превышать 10 мг в сутки. Для аторвастатина инструкция Epclusa рекомендует тщательное наблюдение за нежелательными реакциями со стороны мышц, а не применение к нему предела дозы, установленного для розувастатина. [16]
\item \textbf{Софосбувир/\allowbreak велпатасвир/\allowbreak воксилапревир и статины.} Совместное применение розувастатина не рекомендуется американской инструкцией Vosevi и противопоказано согласно SmPC EMA. В рамках ограничений Vosevi правастатин может применяться в дозе не более 40 мг в сутки. Значительно менее выраженное взаимодействие с правастатином не доказывает безопасности другого статина или дополнительных неисследованных ингибиторов. [17,18]
\item \textbf{Усиленные схемы атазанавира и лопинавира.} NIH ограничивают розувастатин дозой 10 мг в сутки при атазанавире/ритонавире или атазанавире/кобицистате. Американская инструкция розувастатина устанавливает начальную дозу 5 мг и максимальную дозу 10 мг в сутки при атазанавире/ритонавире и лопинавире/ритонавире. [3,4]
\item \textbf{Усиленный дарунавир.} NIH ограничивают розувастатин дозой 20 мг в сутки при дарунавире/кобицистате. Для дарунавира/ритонавира рекомендуется минимальная эффективная доза с мониторингом, без переноса на эту схему максимума, установленного для сочетания с кобицистатом. NIH ограничивают аторвастатин дозой 20 мг в сутки при любой из двух усиленных схем дарунавира; аторвастатин не следует сочетать с атазанавиром/кобицистатом. Симвастатин и ловастатин противопоказаны с перечисленными усиленными ингибиторами протеазы. [4]
\item \textbf{Гемфиброзил.} По возможности следует избегать совместного применения розувастатина. Если сочетание неизбежно, американская инструкция предписывает начинать прием розувастатина с 5 мг в сутки и не превышать 10 мг в сутки. Комбинация также сопряжена с риском аддитивной мышечной токсичности, выходящий за пределы исключительно фармакокинетического объяснения. [3,26]
\item \textbf{Таргетная терапия.} Американский предел дозы розувастатина составляет 5 мг в сутки при даролутамиде и 10 мг в сутки при регорафенибе или капматинибе. Эти ограничения инструкции препарата-субстрата необходимо учитывать вместе с рекомендациями инструкции самого противоопухолевого средства по исключению сочетаний, снижению доз или мониторингу чувствительных субстратов транспортеров. [3,20--22]
\item \textbf{Белумосудил и момелотиниб.} В американской инструкции CRESTOR максимальная суточная доза розувастатина составляет 5 мг с белумосудилом и 10 мг с момелотинибом. Эти ограничения специфичны для названного субстрата и не определяют коэффициентов снижения доз для всех препаратов соответствующего класса. [3,47,48]
\item \textbf{Тафамидис и фебуксостат.} Американская инструкция CRESTOR рекомендует избегать сочетания с тафамидисом; если оно неизбежно, начинать прием розувастатина с 5 мг и не превышать 20 мг в сутки. С фебуксостатом установлен максимум 20 мг в сутки. Наблюдаемое отношение AUC немного менее двух не отменяет прямого ограничения инструкции. [3,49,50]
\item \textbf{Тикагрелор: проверенное изменение 2026 года.} Американская инструкция розувастатина от апреля 2026 года устанавливает максимум 20 мг в сутки при тикагрелоре и рекомендует рассмотреть дозу менее 20 мг у пациентов с повышенным риском рабдомиолиза. Более ранняя рекомендация, ограничивавшаяся мониторингом, не полностью отражает проверенную версию инструкции. [3]
\item \textbf{Тяжелое нарушение функции почек.} Для розувастатина у пациентов, не получающих гемодиализ, при указанном в инструкции клиренсе креатинина ниже \(30\,\mathrm{mL/min/1.73\,m^2}\) начальная доза по американской инструкции составляет 5 мг, а максимальная --- 10 мг в сутки. Если ограничение, обусловленное взаимодействием, ниже, оно сохраняет приоритет. В проверенной SmPC HPRA тяжелое нарушение функции почек при клиренсе креатинина менее 30 мл/мин, напротив, является противопоказанием для всех доз. [3,45]
\item \textbf{Колхицин с ингибиторами.} COLCRYS противопоказан с сильными ингибиторами CYP3A4 или P-gp у пациентов с нарушением функции почек или печени. У пациентов без этих нарушений необходима коррекция дозы с учетом показания и конкретного ингибитора. Снижение дозы статина само по себе не устраняет потенциально смертельное взаимодействие, при котором изменяется экспозиция колхицина. [30]
\end{itemize}

\subsection{Период отмывки: разграничение документированных правил и экстраполяции}

Общего правила FDA, EMA или NIH, предписывающего фиксированный двух-, трех-, семи- или четырнадцатидневный период отмывки для всех ингибиторов OATP1B/BCRP, не существует. Разнесение приема доз на несколько часов нельзя считать способом устранения взаимодействия, обусловленного сохраняющейся системной экспозицией ингибитора, активными метаболитами или длительным связыванием ингибитора с транспортером. Период отмывки в исследовании предназначен для ограничения переноса эффекта между периодами; он не становится автоматически клинической рекомендацией по возобновлению терапии. [1,4,14--18]

Для обратимого ингибитора с элиминацией первого порядка остаточная концентрация после \(n\) периодов полувыведения составляет \(2^{-n}\) исходной. После пяти периодов полувыведения остается приблизительно 3.1\%, но эта концентрация может сохранять ингибирующее действие, если исходная концентрация ингибитора значительно превышала \(K_i\). Пояснительный расчет времени до возобновления терапии имеет вид
\[
t=t_{1/2}\log_2\left(\frac{I_{u,0}}{I_{u,\mathrm{target}}}\right),
\]
при допущении единого эффективного периода полувыведения, отсутствия продолжающегося поступления, значимого активного метаболита и при обратимом равновесном ингибировании. Это механистический инструмент планирования, а не валидированный алгоритм назначения. Максимальную концентрацию на входе в печень во время активного всасывания и кишечный суррогат на основе дозы нельзя просто экстраполировать как тканевые концентрации после отмены. [1,31]

\begin{itemize}
\item \textbf{Указанный в инструкции период сохраняющегося взаимодействия.} Для COLCRYS таблица коррекции дозы при взаимодействиях применяется, если перечисленные ингибиторы принимаются в настоящее время или их прием завершен в предшествующие 14 дней. Это конкретное правило для колхицина; оно не является универсальным требованием полностью отменять колхицин на 14 дней или стандартом отмывки OATP/BCRP. [30]
\item \textbf{Документированный обязательный интервал: системное применение фузидовой кислоты.} В SmPC HPRA розувастатин не применяют во время системной терапии фузидовой кислотой и в течение семи дней после ее прекращения. Это конкретное правило предотвращения миопатии, а не универсальная отмывка OATP/BCRP: инструкция прямо указывает, что фармакокинетический, фармакодинамический либо смешанный механизм взаимодействия не установлен. Предусмотренные инструкцией исключительные случаи длительной терапии требуют индивидуального решения и тщательного врачебного наблюдения. Следовательно, наличие обязательного интервала у одной комбинации не позволяет назначать такой же интервал другим ингибиторам. [45]
\item \textbf{Временная отмена статина при ограниченном по длительности противовирусном курсе.} Если принято решение о перерыве, назначающий врач должен документировать последнюю дозу статина, дату завершения противовирусной терапии, клинически значимое сохранение исходного вещества и метаболитов в организме, сердечно-сосудистый риск перерыва и план возобновления. Не следует вводить обязательный интервал после препаратов прямого противовирусного действия, если инструкция его не предусматривает. Совместимая альтернатива может быть предпочтительнее отказа от гиполипидемической терапии на весь противовирусный курс. [14--18]
\item \textbf{Сохраняющееся ингибирование или индукция ферментов.} Сопутствующий препарат, вызывающий взаимодействие, может влиять на восстановление активности CYP независимо от восстановления функции транспортеров. Инактивация CYP2C8 глюкуронидом гемфиброзила и индукция при повторном применении рифампицина показывают, почему исчезновение исходного вещества не обязательно восстанавливает все пути элиминации. [1,37,39]
\end{itemize}

\subsection{Фармакогенетическая стратификация}

\begin{itemize}
\item \textbf{\textit{SLCO1B1} c.521T\(>\)C, rs4149056.} Замена p.Val174Ala встречается в гаплотипах со сниженной функцией и может нарушать печеночный захват статинов. В исследовании SEARCH с симвастатином 80 мг шансы миопатии увеличивались в 4.5 раза на каждый аллель C (95\% ДИ 2.6--7.7), а у гомозигот CC отношение шансов по сравнению с TT составляло 16.9 (4.7--61.1). Это генотип-ассоциированные отношения шансов для конкретного статина и дозы, а не отношения AUC розувастатина или универсальные множители риска. [35]
\item \textbf{\textit{ABCG2} c.421C\(>\)A, rs2231142.} Замена p.Gln141Lys снижает функцию ABCG2. В соответствующей системе интерпретации CPIC генотипы CC, CA и AA соотносятся соответственно с нормальной, сниженной и низкой функцией. Функциональные последствия включают изменение эффлюкса субстратов; генетические эксперименты по транспорту урата подтверждают функциональную значимость снижения активности, однако для оценки эффектов в отношении конкретных лекарств требуются фармакокинетические данные у человека. [10,34]
\item \textbf{Количественный эффект \textit{SLCO1B1} c.521T\(>\)C.} В исследовании с участием 32 здоровых добровольцев после однократного приема розувастатина 10 мг у носителей CC (\(n=4\)) AUC за 0--48 часов и \(C_{\max}\) были соответственно на 65\% и 79\% выше, чем у носителей TT (\(n=16\)); отношения составляют 1.65 и 1.79. Для аторвастатина 20 мг в той же работе отношение AUC за 0--48 часов между CC и TT составляло 2.44. Это прямое свидетельство субстрат-зависимого генетического эффекта, а не основание считать все статины одинаково зависимыми от OATP1B1. Малочисленная группа гомозигот ограничивает точность оценки. [51]
\item \textbf{Количественный эффект \textit{ABCG2} c.421C\(>\)A.} В отдельном фармакокинетическом исследовании с участием 32 добровольцев после приема розувастатина в дозе 20 мг у носителей AA (\(n=4\)) AUC, экстраполированная до бесконечности, была на 144\% выше, а \(C_{\max}\) --- на 131\% выше, чем у CC (\(n=16\)); соответствующие отношения равны 2.44 и 2.31. По сравнению с CA (\(n=12\)) AUC у AA была выше на 100\%. Эти результаты относятся к различиям генотипов в конкретном исследовании, а не к эффекту добавления лекарственного ингибитора; перемножать их с отношениями для циклоспорина или противовирусной схемы без валидированной модели нельзя. [46]
\item \textbf{Начальные дозы розувастатина с учетом генотипа.} CPIC рекомендует начальную дозу не более 20 мг при низкой функции ABCG2 и нормальной функции SLCO1B1 либо при низкой функции SLCO1B1 и нормальной/сниженной функции ABCG2. При сочетании низкой функции ABCG2 со сниженной или низкой функцией SLCO1B1 рекомендация начинать не более чем с 10 мг отнесена к категории «необязательная». Это рекомендации по начальной дозе, а не разрешение превышать более низкий предел, обусловленный межлекарственным взаимодействием. [34]
\item \textbf{Объем тестирования.} CPIC определяет, как использовать имеющиеся результаты, но не предписывает универсальное тестирование перед назначением статина. Результат исследования только варианта c.521 не позволяет полностью исключить другие аллели SLCO1B1 со сниженной функцией. В статье, включенной в майский выпуск журнала за 2026 год и впервые опубликованной онлайн 1 декабря 2025 года, авторы публикации CPIC, описали обновления таблицы функциональных характеристик аллелей с учетом данных участников преимущественно африканского происхождения из регионов к югу от Сахары. Эти обновления уточняют функциональную аннотацию, но сами по себе не формируют нового набора рекомендаций по дозам статинов. Полная интерпретация гаплотипов и актуальная лабораторная таблица соответствий предпочтительнее предположений, основанных лишь на происхождении пациента. [34,36]
\end{itemize}

Фармакологическое ингибирование может изменять функциональный фенотип, ожидаемый на основании генотипа, то есть вызывать феноконверсию; при этом генетические и лекарственные эффекты не обязательно перемножаются. У пациента с исходно низкой активностью транспортера относительный эффект ингибитора может быть меньше при сохранении высокой абсолютной экспозиции; напротив, подавление оставшегося параллельного пути может создавать существенный дополнительный риск. Валидированного универсального калькулятора, позволяющего рассчитать безопасную дозу с одновременным учетом генотипов c.521 и c.421, нескольких ингибиторов и нарушений функции органов, не существует. Это ограничение применения имеющихся данных в клинической практике, а не основание игнорировать генетическую информацию или сведения о взаимодействиях. [12,31,34,36]

\paragraph{Нутрицевтическая феноконверсия: рабочая концепция и ограничения.}
Под нутрицевтической феноконверсией можно понимать временное изменение измеряемой транспортерной функции под действием пищевого компонента или БАД, вследствие которого функциональный ответ отклоняется от ожидаемого по генотипу. При достаточной локальной концентрации ингибитора фенотипически нормальный транспортер может обеспечивать меньший поток субстрата, что в отношении отдельной PK-конечной точки способно напоминать снижение функции, ассоциированное с \textit{SLCO1B1} c.521T\(>\)C или \textit{ABCG2} c.421C\(>\)A. Однако «медленный переносчик» не является валидированной общей диагностической категорией для потребителей полифенолов. Локальное кишечное ингибирование неэквивалентно генетическому изменению транспортера во всех тканях, а функциональная аналогия не означает изменения генотипа или доказанной эквивалентности клинического риска. [34,46,51,79,95]

В качестве иллюстративной, а не валидированной клинической модели для независимых параллельных путей при малой концентрации субстрата и быстром равновесии обратимого конкурентного ингибирования можно записать
\[
CL_{\mathrm{uptake,eff}}=
CL_{\mathrm{passive}}+
\sum_j \frac{g_j\,CL_{\mathrm{uptake},j}^{\mathrm{ref}}}
{1+I_{u,j}/K_{i,j}},
\]
где \(g_j\) обозначает относительную функциональную активность пути с учетом генотипа, \(I_{u,j}\) --- локальную несвязанную концентрацию ингибитора, а \(K_{i,j}\) --- константу для соответствующей пары «ингибитор--субстрат». Модель не включает стимуляцию транспорта, смешанное ингибирование, насыщение, конкуренцию нескольких ингибиторов и изменение экспрессии. Она показывает, почему генетические и нутрицевтические отношения AUC нельзя автоматически перемножать: остаточный клиренс представляет собой сумму путей, а системная экспозиция дополнительно зависит от биодоступности и последующей элиминации. В исследовании экстракта зеленого чая вариант \textit{ABCG2} c.421C\(>\)A влиял на исходную экспозицию розувастатина, но значимого влияния на величину взаимодействия с экстрактом не обнаружено; малая выборка ограничивает обобщение. [31,34,82]

\paragraph{Многокомпонентное ингибирование и масштабирование захвата.}
Для смеси обратимых конкурентных ингибиторов, связывание которых с одной изоформой \(j\) взаимно исключается, при сохранении допущений исходной модели знаменатель транспортного слагаемого заменяется на \(1+\sum_m I_{u,m}/K_{i,mj}\):
\[
CL_{\mathrm{uptake,eff}}=
CL_{\mathrm{passive}}+
\sum_j\frac{g_j\,CL_{\mathrm{uptake},j}^{\mathrm{ref}}}
{1+\displaystyle\sum_m I_{u,m}/K_{i,mj}}.
\]
Индекс \(m\) относится к отдельному компоненту или активному метаболиту; \(I_{u,m}\) и \(K_{i,mj}\) должны быть выражены в согласованных единицах и относиться к несвязанной фракции. Аддитивность в знаменателе следует из баланса взаимно исключающих состояний свободного и занятого переносчика. Произведение \(\prod_m(1+I_{u,m}/K_{i,mj})\) не является равнозначной альтернативой: уже для двух лигандов оно содержит дополнительный член \(I_{u,1}I_{u,2}/(K_{i,1j}K_{i,2j})\). Такая форма требует иной модели независимого одновременного связывания с мультипликативным уменьшением остаточной активности; ее нельзя вводить только на основании многокомпонентности экстракта. При смешанном ингибировании, кооперативности или транспорте конкурирующих лигандов выбор схемы и параметров требует отдельной экспериментальной проверки. Эти уравнения представляют пояснительный вывод при заданных допущениях, а не валидированную модель силимарина. [1,59,85]

Для силибина A и силибина B экспериментальные \(\mathrm{IC}_{50}\) ингибирования OATP1B3 составляли соответственно 2.7 и \(5.0\,\mu\mathrm{M}\); это показатели ингибирующей активности в конкретном опыте, а не непосредственно измеренные константы связывания. Если \(I_{u,m}=w_m I_{u,\mathrm{mix}}\) при фиксированных долях \(w_m\), \(\sum_m w_m=1\), конкурентная модель допускает условную константу смеси
\[
\frac{1}{K_{i,\mathrm{mix},j}}=
\sum_m\frac{w_m}{K_{i,mj}}.
\]
Однако доли несвязанных компонентов на входе в печень не обязаны совпадать с массовыми долями в препарате и могут изменяться во времени. Поэтому константа «суммарной смеси» без состава, компонентной PK и условий опыта не идентифицирует ее действие in vivo; условная эмпирическая константа пригодна лишь для проверенных состава и условий. [85]

Коэффициент относительной активности транспортера (RAF) масштабирует активность соответствующего пути между экспериментальной системой и гепатоцитами, а не заменяет компонентные константы ингибирования. Неодинаковая активность диастереомеров не делает сам RAF математически неопределимым, но препятствует обоснованному прогнозу остаточного захвата, если неизвестны их несвязанные концентрации и \(K_{i,mj}\). Даже корректно масштабированный \(CL_{\mathrm{uptake,eff}}\) остается параметром входа в клетку, а не полным печеночным клиренсом: последний требует учета обратного эффлюкса, метаболизма, билиарного выведения и кровотока в общей кинетической модели баланса. [31,59,85]

\paragraph{Пассивный перенос и мембранные эффекты смеси.}
Исходная модель предполагает неизменность \(CL_{\mathrm{passive}}\); это скалярный параметр, а не вектор диффузии. В опытах с фосфатидилхолиновыми бислоями силибин взаимодействовал преимущественно с поверхностной областью мембраны, тогда как изменения свойств ее более глубоких областей были ограниченными. Такие наблюдения обосновывают проверку мембранных эффектов, но не устанавливают увеличения пассивного входа статинов в гепатоциты при клинической экспозиции силимарина. [100]

Для анализа чувствительности, а не обязательной коррекции, при фиксированном составе смеси можно ввести локальную линейную аппроксимацию
\[
CL_{\mathrm{passive}}(I_{\mathrm{mix},u})=
CL_{\mathrm{passive},0}\left(1+\alpha I_{\mathrm{mix},u}\right).
\]
\(I_{\mathrm{mix},u}\) обозначает явно определенный локальный показатель несвязанной экспозиции смеси, а не массу принятого экстракта; \(\alpha\) имеет размерность обратной концентрации и зависит от субстрата, состава смеси и мембранной системы. Линейная форма применима лишь в исследованном диапазоне малых возмущений, где \(|\alpha I_{\mathrm{mix},u}|\ll1\) и клиренс остается положительным. Знак \(\alpha\) не задан априори; при отсутствии измерений базовым остается \(\alpha=0\), а положительные и отрицательные значения служат сценариями, не доказанными эффектами силимарина. Изменение микровязкости само по себе не определяет знак изменения проницаемости. [100]

\paragraph{Численная нижняя граница пассивного клиренса.}
Линейная аппроксимация сама по себе не имеет сингулярности, но при экстраполяции и отрицательном \(\alpha\) способна давать нефизический отрицательный клиренс. Для численного анализа можно явно задать нижнюю границу, а не «положительный супремум»:
\[
\begin{aligned}
CL_{\mathrm{passive}}(I_{\mathrm{mix},u})
&=\max\!\left\{CL_{\min},
CL_{\mathrm{passive},0}(1+\alpha I_{\mathrm{mix},u})\right\},\\
CL_{\min}&=\gamma CL_{\mathrm{passive},0},
\qquad 0<\gamma\leq1.
\end{aligned}
\]
При положительном \(CL_{\mathrm{passive},0}\) это обеспечивает положительность расчетного коэффициента, но не обосновывает физиологического значения \(\gamma\). Без экспериментальных данных \(\gamma\) является параметром численной регуляризации и анализа чувствительности. При \(\gamma=1\) отрицательная мембранная поправка полностью отсекается; функция имеет излом на границе отсечения и не ограничена сверху при положительном \(\alpha\). Долю стохастических реализаций, попавших на отсечение, следует регистрировать: частое срабатывание указывает на выход за область применимости исходной аппроксимации, а не подтверждает плато проницаемости. Сами концентрации и исходный клиренс должны задаваться неотрицательными и положительным соответственно.

Нижняя граница не доказывает сохранения реальной проницаемости при повреждении мембраны и не делает экстраполяцию биофизически валидной. Неотрицательность коэффициента совместно с корректной движущей силой исключает пассивный поток против градиента в принятой модели; двунаправленность задается разностью химических потенциалов, а не операцией \(\max\). Строго положительный универсальный минимум не является термодинамическим требованием: нулевая проницаемость также не нарушает условия отсутствия переноса против градиента. Эта регуляризация не представлена как установленное действие силимарина. [100]

Термодинамически пассивный путь сохраняет двунаправленность. Для одной нейтральной проникающей формы при одинаковых условиях равновесного распределения по обе стороны мембраны чистый поток выражается как
\[
J_{\mathrm{passive}}=
CL_{\mathrm{passive}}(I_{\mathrm{mix},u})
\left(C_{u,\mathrm{ext}}-C_{u,\mathrm{int}}\right),
\]
где \(J_{\mathrm{passive}}\) имеет размерность количества вещества за время. При различиях ионизации или электрохимического потенциала необходима запись через соответствующие активности и движущую силу. Подавление OATP не превращает пассивный путь в односторонний вход; его усиление увеличивает также возможность обратного выхода. Полное ингибирование и насыщение OATP избытком субстрата неэквивалентны. Следовательно, игнорирование мембранного эффекта может смещать оценку в любую сторону, а утверждение об обязательной недооценке остаточного захвата не следует из модели. Эта зависимость является пояснительным физическим ограничением, а не измеренной характеристикой нутрицевтического взаимодействия.

Для подтверждения феноконверсии в исследовательском протоколе нужны исходная оценка, количественная характеристика продукта, повторное измерение функции во время экспозиции и оценка восстановления после отмены. CP-I может дополнять анализ печеночного OATP1B, но не является показателем кишечного BCRP; CP-III не обеспечивает селективного измерения OATP1B3. Само по себе употребление экстракта не служит основанием присваивать пациенту категорию сниженной функции CPIC, изменять интерпретацию генетического теста или переносить на него дозовые ограничения для определенного генотипа. [34,52,53]

\subsection{Предлагаемый алгоритм применения данных в клинической практике}

Следующая схема представляет собой практический синтез для врачей и фармацевтических специалистов, а не самостоятельный утвержденный регуляторный протокол. Приоритет отдается профилактике, а не лабораторному наблюдению после уже состоявшегося назначения противопоказанной комбинации. [3--5,14--18,30,34]

\begin{itemize}
\item \textbf{Проверить полную схему терапии до назначения нового препарата.} Зафиксировать конкретный статин, дозу, усилитель, противовирусные компоненты, эпизодически принимаемые препараты, добавки и время приема последних доз. Определить роль каждого средства: субстрат, ингибитор, индуктор, источник аддитивной миотоксичности, нефротоксин или сочетание ролей. Указать, подтверждено ли взаимодействие клинически, прогнозируется ли моделью или остается неизученным.
\item \textbf{Сначала учесть противопоказания и ограничения доз.} До начала лечения устранить противопоказанные сочетания и комбинации, которых рекомендовано избегать. Использовать действующую информацию о препарате соответствующей юрисдикции, затем применять наиболее строгое из оставшихся ограничений, связанных с взаимодействиями, почечной функцией и фармакогенетикой. Не компенсировать наличие нескольких неисследованных ингибиторов произвольным дополнительным процентным снижением дозы.
\item \textbf{Зафиксировать исходное клиническое состояние и лабораторные показатели.} Оценить функцию почек, наличие заболеваний печени, мышечные симптомы, предшествующую непереносимость статинов и обратимые сопутствующие факторы, например неконтролируемый гипотиреоз. Выполнить исходные лабораторные исследования, соответствующие субстрату и профилю риска; исходное определение активности креатинкиназы особенно полезно, если симптомы, прежняя токсичность или значительная нагрузка взаимодействующими препаратами могут усложнить последующую интерпретацию. Рутинное серийное определение активности креатинкиназы у каждого бессимптомного пациента не является универсальным требованием.
\item \textbf{Дать инструкции по действиям при появлении симптомов.} Впервые возникшие выраженная слабость, мышечная боль, темная моча, уменьшение диуреза или общее недомогание требуют своевременной оценки. При подозрении на миопатию или рабдомиолиз необходимо прекратить прием статина, предположительно вызвавшего реакцию, и провести клиническое обследование, включая, по показаниям, определение активности креатинкиназы, оценку функции почек, определение электролитов и анализ мочи. Ранее нормальная активность креатинкиназы не гарантирует безопасности последующей терапии.
\item \textbf{Повторно оценивать состояние после изменения схемы и отмены.} Организовать ранний контроль, когда экспозиция ингибитора предположительно приближается к соответствующему равновесному состоянию; сроки индивидуализировать, а не устанавливать единый день обследования для всех пациентов. После коррекции дозы повторно оценить липидный ответ: сниженный системный клиренс не гарантирует сохранения фармакологического эффекта в печени. Отслеживать каждую временную отмену, чтобы предотвратить неоправданно длительную отмену препарата или преждевременное возобновление его приема в высокой дозе.
\item \textbf{Использовать специализированные алгоритмы для конкретного субстрата.} Метотрексат, противоопухолевые препараты и иммуносупрессия после трансплантации требуют собственных стандартов мониторинга и коррекции доз. Отношения AUC, полученные для статинов, мониторинг мышечных симптомов и генотипические ограничения доз не заменяют эти стандарты. При неопределенности сопоставления экспозиций совместимая альтернативная схема предпочтительнее невалидированного кинетического расчета.
\end{itemize}

\paragraph{Дополнительный контур оценки БАДов и пищевых концентратов.}
В предлагаемом алгоритме полная сверка терапии должна включать обязательное для данного процесса документирование всех высокодозовых экстрактов, биоусилителей и регулярно употребляемых концентрированных или специализированных соков. Следует записывать торговое название, производителя, по возможности номер партии или фотографию маркировки, массу экстракта и содержание активных компонентов в одной порции, число порций, время приема, связь с пищей, длительность и цель применения. Отдельно фиксируют пиперин в составе комбинированного продукта, даже если он не вынесен в его название. Это предложение по качеству лекарственного анамнеза, а не утверждение о специальной юридической обязанности, установленной ICH M12. Изменение продукта или его дозы рассматривают как изменение схемы терапии. [1,80,82--87]

Интервал 2--4 часа между добавкой и лекарством не следует считать универсальной гарантией отсутствия взаимодействия, однако неверно утверждать, что такое разнесение всегда неэффективно. В клиническом исследовании грейпфрутового сока и фексофенадина прием сока одновременно или за два часа до лекарства снижал экспозицию, тогда как при интервале четыре часа значимого эффекта не обнаружено. Для CYP3A4-зависимых взаимодействий фуранокумаринов механизм иной, и исчезновение вещества из просвета кишечника не означает немедленного восстановления ферментативной активности. Для перечисленных катехинов и БАДов длительное внутриклеточное удержание и медленная диссоциация комплексов с OATP/BCRP не доказаны как общий клинический механизм. Продолжительность риска должна определяться данными о конкретном сочетании, а не предположением о едином времени восстановления всех транспортеров. [89,96]

Пока для конкретного продукта и препарата не установлена совместимость, изменение времени приема не должно подменять пересмотр необходимости добавки. При высоких дозах розувастатина, включая дозы более 20 мг/сут, при приеме аторвастатина, метотрексата, апиксабана или ривароксабана предлагается не начинать необязательные концентрированные БАДы с пиперином, EGCG или силимарином без предварительной оценки лечащим врачом или клиническим фармакологом; при уже начатом совместном приеме следует обсудить исключение добавки на период оценки. Если ее польза не обоснована, состав или экспозиция неопределенны либо присутствуют дополнительные взаимодействия, предпочтительно исключить добавку, а не снижать доказанно эффективную лекарственную терапию по невалидированной модели. Это предосторожностная стратегия, а не универсальное зарегистрированное противопоказание для трех растительных веществ. Порог розувастатина 20 мг не является установленной границей нутрицевтического взаимодействия и не означает отсутствия риска при меньшей дозе. [1,3,34,60,61,81--86]

Субстрат-специфичная оценка для статинов включает не только мышечную переносимость, но и последующий липидный ответ: экстракт может уменьшать, а не увеличивать системную экспозицию. Для метотрексата первостепенны доза, путь введения, почечная функция и предусмотренный соответствующим протоколом мониторинг; при внутривенном введении повышение кишечной биодоступности вообще не является механизмом взаимодействия. Для НОАК оценивают совокупное воздействие на P-гликопротеин, CYP3A4, почечное выведение и гемостаз; BCRP-специфичный риск нельзя выводить из результатов с сульфасалазином или клеточных опытов с пиперином. Не следует самостоятельно отменять антикоагулянт или иное назначенное лекарство из-за выявленного приема БАД; приоритетом служит профессиональная оценка добавки и схемы лечения. [60,61,77,78,83,84,95]

Необоснованное расширение запрета на обычные пищевые порции, все виды чая, все цитрусовые или все формы силимарина также нежелательно. Необходимо сохранять различие между пищей и концентратом, клинически проверенной схемой и неизученной комбинацией, предосторожностью и доказанным противопоказанием. Отрицательное исследование не доказывает безопасности любых доз, но положительный эксперимент in vitro не устанавливает необходимости полного отказа от целой группы продуктов. [79--90]

\paragraph{Недостаточный липидный ответ при сопутствующем приеме БАД.}
При документированном приеме концентратов с пиперином, EGCG или силимарином и недостаточном ответе на розувастатин либо аторвастатин предлагается сначала проверить приверженность, фактическую дозу и длительность терапии, питание, вторичные причины дислипидемии и остальные взаимодействия. Если добавка необязательна и взаимодействие представляется правдоподобным, лечащий врач может согласовать ее полное временное исключение с заранее определенной повторной оценкой липидного ответа. По возможности остальные условия сохраняют неизменными, если это клинически безопасно. Такая диагностическая отмена является предлагаемой мерой предосторожности, а не обязательным требованием ICH M12 или доказательством причинной роли любого из трех экстрактов. [1,3,81,82,85,86]

Универсальное правило «не менее пяти периодов полувыведения нативного полифенола» не определяет время восстановления транспортной функции. В простейшей модели моноэкспоненциальной элиминации без продолжающегося всасывания и образования активных метаболитов
\[
\frac{I_u(5t_{1/2})}{I_u(0)}=2^{-5}=0.03125,
\qquad
\frac{I_u(5t_{1/2})}{K_i}=\frac{I_u(0)}{32K_i}.
\]
Даже остаток 3.125\% может быть значим при высоком исходном \(I_u(0)/K_i\); плазменный период полувыведения агликона не описывает автоматически длительность действия метаболитов, тканевой экспозиции или восстановления функции. Срок исключения определяют по продукту и клинической ситуации, а срок оценки ЛПНП --- также по динамике фармакодинамического ответа. Для розувастатина максимальное снижение ЛПНП обычно достигается примерно к четвертой неделе лечения; использование сопоставимого срока наблюдения после устранения предполагаемого взаимодействия является клинической экстраполяцией, а не установленным сроком отмывки БАД. [1,3]

\paragraph{Восстановление кишечного транспорта и время оценки липидного ответа.}
Ингибирование кишечного BCRP может изменять долю всосавшегося вещества \(F_a\), форму входного потока и наблюдаемые \(C_{\max}\) и \(T_{\max}\). Однако из изменения пика нельзя однозначно вывести изменение единственной константы \(k_a\). В принятом разложении \(F=F_aF_gF_h\) величина \(F_g\) --- доля, избежавшая кишечного метаболизма, а не фракция кишечной экстракции; последнее понятие соответствовало бы \(1-F_g\). Взаимодействие эффлюкса и метаболизма может вторично изменять \(F_g\), но изолированное ингибирование BCRP не требует обязательного изменения этого параметра. [1,31,95]

При обратимом связывании и неизменном количестве функционального белка восстановление активности не требует обязательного ресинтеза ABCG2. В простейшей модели занятости участка
\[
\frac{d\theta}{dt}=k_{\mathrm{on}}I_{u,\mathrm{site}}(1-\theta)
-k_{\mathrm{off}}\theta,
\]
где \(\theta\) --- доля занятых участков, после исчезновения доступного ингибитора \(\theta(t)=\theta(0)e^{-k_{\mathrm{off}}t}\). Эта схема иллюстрирует зависимость восстановления от локальной экспозиции и диссоциации, не подменяя многостадийный транспортный цикл ABCG2 одностадийным связыванием. Обновление белкового пула становится отдельным ограничением, если показаны снижение экспрессии, утрата белка или необратимая инактивация. Для условного восстановления пула \(E\) при постоянном синтезе можно записать
\[
\begin{aligned}
\frac{dE}{dt}&=k_{\mathrm{syn}}-k_{\mathrm{deg}}E,\\
E(t)&=E_0+\bigl[E(0)-E_0\bigr]e^{-k_{\mathrm{deg}}t},
\qquad E_0=\frac{k_{\mathrm{syn}}}{k_{\mathrm{deg}}}.
\end{aligned}
\]
Здесь предполагается прекращение повреждающего воздействия и восстановление исходных параметров синтеза и деградации. Время уменьшения исходного дефицита до доли \(\varepsilon\) равно \(t_{\varepsilon}=\ln(1/\varepsilon)/k_{\mathrm{deg}}\); оно требует измеренного \(k_{\mathrm{deg}}\) и выбранного критерия восстановления, а не произвольного календарного интервала. Данные у человека, обосновывающие обязательные 7--14 дней полного ресинтеза кишечного ABCG2 после приема перечисленных БАДов, в рассмотренных источниках не установлены. Поэтому такой интервал не вводится как физиологическая константа или универсальное клиническое требование. [1,95,96] Приведенное уравнение является пояснительной моделью баланса белкового пула, а не измеренной в этих исследованиях кинетикой его восстановления.

\paragraph{Мембранный резервуар и повторное связывание.}
Локальная концентрация около мембранной мишени не обязана совпадать со среднеобъемной водной концентрацией: такой градиент непосредственно измерен для флуоресцентного производного пропранолола у живых клеток. Это экспериментальный пример микрофармакокинетики, а не количественное доказательство удержания пиперина в энтероцитарной мембране или определенной длительности ингибирования BCRP. Коэффициент \(K_{\mathrm{lip}}=C_{\mathrm{lip,eq}}/C_{\mathrm{w,eq}}\) характеризует равновесное распределение при заданных условиях; он не определяет по отдельности скорости обмена и диссоциации комплекса с белком. Поэтому равенство \(k_{\mathrm{off,eff}}=k_{\mathrm{off}}/K_{\mathrm{lip}}\) не является общим кинетическим законом. [83,84,104,105]

Для иллюстрации мембранного резервуара при линейном обмене, постоянном объеме липидной фазы и пренебрежимо малом истощении резервуара связыванием с мишенью можно использовать
\[
\begin{aligned}
\frac{dA_{\mathrm{lip}}}{dt}
&=CL_{\mathrm{ex}}\left(C_{\mathrm{w}}-
\frac{C_{\mathrm{lip}}}{K_{\mathrm{lip}}}\right)
-k_{\mathrm{loss,lip}}A_{\mathrm{lip}},\\
C_{\mathrm{lip}}&=\frac{A_{\mathrm{lip}}}{V_{\mathrm{lip}}}.
\end{aligned}
\]
Здесь \(CL_{\mathrm{ex}}\) --- объемный коэффициент обмена, \(K_{\mathrm{lip}}>0\) --- безразмерное отношение согласованных концентраций, \(C_{\mathrm{w}}\) --- концентрация доступного вещества в смежной водной фазе, а \(k_{\mathrm{loss,lip}}\) --- константа других потерь. В пределе постоянно поддерживаемой нулевой \(C_{\mathrm{w}}\) мембранный пул убывает с константой
\[
k_{\mathrm{release}}=
\frac{CL_{\mathrm{ex}}}{K_{\mathrm{lip}}V_{\mathrm{lip}}}
+k_{\mathrm{loss,lip}}.
\]
Это скорость опустошения резервуара в данной модели, а не исправленный \(k_{\mathrm{off}}\) белкового комплекса. Исчезновение вещества из химуса не устанавливает нулевую концентрацию во всех водных микрокомпартментах энтероцита.

Если мембранный пул поддерживает ненулевую \(I_{u,\mathrm{site}}(t)\), в уравнении занятости сохраняется член \(k_{\mathrm{on}}I_{u,\mathrm{site}}(1-\theta)\); повторное связывание может замедлять наблюдаемое снижение занятости без обязательного изменения микроскопического \(k_{\mathrm{off}}\). Связь между общей \(C_{\mathrm{lip}}\) и доступной концентрацией у мишени требует отдельной параметризации с учетом расположения лиганда в бислое и доступа к участку связывания. Необходимо раздельно измерять распределение, скорость выхода из мембраны и диссоциацию при подавлении повторного связывания. Пока таких данных для конкретной пары «нутрицевтик--BCRP» нет, мембранный резервуар остается проверяемым объяснением, а не основанием для универсального периода отмывки. [104,105]

Диагностическая отмена должна раздельно учитывать исчезновение ингибитора и активных метаболитов из соответствующего компартмента, обратимость воздействия и более медленную динамику липидного ответа. Срок повторного измерения ЛПНП не служит суррогатом времени обновления ABCG2; необоснованное ожидание «полного ресинтеза» не должно подменять оценку конкретного сочетания или задерживать необходимую коррекцию терапии. [1,3]

Недостаточное снижение ЛПНП, даже при измеренном повышении плазменной экспозиции статина, совместимо с нарушением печеночного захвата, но не является его прямым специфичным клиническим доказательством. Концентрация в плазме не заменяет измерение экспозиции в гепатоците; для причинного вывода необходимы временная связь, исключение альтернативных объяснений и согласованные изменения после отмены. Улучшение липидного ответа после исключения добавки поддерживает гипотезу взаимодействия, но само по себе не идентифицирует OATP1B или BCRP. До уточнения причины не следует компенсировать предполагаемую интерференцию эмпирическим повышением дозы на основании невалидированной модели; при этом обследование не должно неоправданно задерживать необходимую коррекцию доказательной гиполипидемической терапии с учетом сердечно-сосудистого риска. [3,31,43,81,82,86]

\section{Заключение и пробелы в оценке взаимодействий при полифармакотерапии}

Оценка взаимодействий, сосредоточенная исключительно на CYP, недостаточна, однако метаболизм с участием CYP сохраняет свое значение. Взаимодействия «розувастатин--ингибитор протеазы», «циклоспорин--статин», «даролутамид--розувастатин» и «противовирусный препарат--статин» показывают, что захват и эффлюкс могут оказывать определяющее влияние на экспозицию. Капматиниб и тикагрелор также демонстрируют, что клинически значимые взаимодействия с преимущественным участием BCRP не требуют столь же значимого ингибирования OATP. Метаболиты гемфиброзила и сочетания аторвастатина с усиленными схемами подчеркивают еще одно положение: анализ транспортеров не заменяет анализ ферментов. [12,20,22,25,28,39]

Следовательно, оценка лекарственных и нутрицевтических взаимодействий должна охватывать всю сеть фармакокинетических процессов: поступление из кишечника, синусоидальный захват, внутриклеточный метаболизм, канальцевый экспорт, почечное выведение, тканевое распределение и восстановление функции соответствующих путей после отмены ингибитора. Двойная блокада OATP/BCRP --- важное сочетание механизмов в этой системе, а не единственное объяснение токсичности при кардиометаболической терапии. Решения по безопасности должны быть привязаны к конкретному препарату-субстрату, вызывающей взаимодействие схеме, дозе, фенотипу пациента, функции органов и регуляторной юрисдикции. [1--5,31]

\begin{itemize}
\item \textbf{Неселективность зондов.} Розувастатин полезен именно потому, что выявляет клинически значимые комбинированные эффекты, однако эта чувствительность ограничивает возможность отнести наблюдаемый эффект к конкретному транспортеру. Исследования должны сочетать лекарственные зонды с эндогенными маркерами, валидированными клеточными анализами и независимой проверкой пригодности моделей, а не интерпретировать любое увеличение AUC розувастатина как двойную блокаду. [12,31,37]
\item \textbf{Недостаток исследований с несколькими препаратами, вызывающими взаимодействие.} Большинство клинических оценок получено для одной взаимодействующей схемы и одного препарата-субстрата в контролируемых условиях. Они не определяют распределение экспозиции при одновременной противовирусной, иммуносупрессивной, антиагрегантной и гиполипидемической терапии. Механистические модели взаимодействий должны учитывать перекрывающиеся пути и насыщение, а не просто перемножать попарные отношения. [4,14,17,31]
\item \textbf{Ограниченность измерений тканевой экспозиции и экспозиции несвязанной фракции.} Плазменную AUC проще определить, чем экспозицию в гепатоцитах, скелетных мышцах, опухоли или почечных канальцах. Поэтому вклад противоположно направленных изменений поступления и эффлюкса трудно разграничить даже при точном описании системной экспозиции. Нельзя судить об эффективности и повреждении на уровне отдельных тканей только по общей плазменной концентрации. [9,11,25,31]
\item \textbf{Недостаточная генетическая и популяционная представленность.} Исследования распространенных вариантов не охватывают все функциональные аллели, популяционные группы, возрастные диапазоны и фенотипы нарушений функции органов. Обновление аннотации SLCO1B1 в 2026 году иллюстрирует и прогресс, и необходимость поддерживать актуальность таблиц интерпретации. Проспективные доказательства того, что поддержка решений с совместным учетом генотипа и взаимодействий предотвращает тяжелые нежелательные исходы, остаются менее многочисленными, чем данные о фармакокинетических ассоциациях. [34--36]
\item \textbf{Неоднородность доказательств токсичности.} Контролируемые исследования экспозиции, отдельные случаи повышения активности креатинкиназы, ассоциации с почечными исходами в наблюдательных исследованиях и сообщения о редких событиях отвечают на разные вопросы. На их основании нельзя устанавливать необоснованный универсальный порог концентрации или утверждать неизбежность развития рабдомиолиза либо нефропатии при увеличении AUC в два--семь раз. [25--29,32]
\item \textbf{Зависимость клинических ограничений от версии документа.} Различия между американским и европейским пределами розувастатина при глекапревире/пибрентасвире, изменение политики NIH в мае 2026 года и ограничение розувастатина при тикагрелоре от апреля 2026 года показывают, что запись о взаимодействии должна включать источник, юрисдикцию, дату редакции и дату обращения. Статический перечень без этих полей может стать небезопасным даже при сохранении правильности исходных фармакокинетических значений. [3,5,14,15]
\end{itemize}

Клиническая цель --- сохранение терапевтической пользы при экспозиции, допустимость которой обоснована имеющимися данными, а не максимальное подавление каждого транспортера. Для ее достижения необходимы исключение запрещенных сочетаний, соблюдение конкретных ограничений доз, планирование терапии с учетом остаточного ингибирования и осторожная интерпретация генотипа и тканевой экспозиции. Подход, учитывающий транспортеры, устраняет пробелы скрининга, учитывающего только CYP, не заменяя одну чрезмерно упрощенную модель другой. [1--5,24,30,34]

\paragraph{Открытая наука и воспроизводимость последующих обновлений.}
Для последующего обновления обзора предлагается открытая версионируемая таблица извлечения данных с отдельными полями для роли препарата, экспериментальной системы, субстрата-зонда, предварительной инкубации, связывания, дозы, пути и времени введения, генотипа, функции органов, определения AUC, типа среднего, доверительного интервала и клинического исхода. Измеренные \(\mathrm{IC}_{50}\) и \(K_i\) следует отделять от параметров, подобранных по клиническим данным; к моделям необходимо прилагать уравнения, параметры, версии программ и анализ чувствительности. Отрицательные результаты для конкретных сочетаний должны сохраняться наравне с положительными. Это предлагаемые требования к воспроизводимости, а не заявление о выполненной регистрации протокола, независимом двойном отборе или уже опубликованном открытом наборе данных. [1,2,52--78]

Для пищевого контура дополнительно предлагается сохранять химическую характеристику партии, соотношение агликонов и конъюгатов, форму дозирования, прием натощак, состав сопутствующей пищи, аналитическое определение «свободной» формы и разделение родительского соединения и метаболитов. Если текст статьи и таблицы дают несогласующиеся численные отношения, необходимо фиксировать расхождение и способ извлечения, не подменяя его недокументированным пересчетом. Разработка модели совместного генетического и нутрицевтического воздействия должна включать отрицательные клинические результаты и независимую проверку прогноза; ее калибровка на одном субстрате не подтверждает переносимость на другие препараты. [79--97]

\section*{Список литературы}

\begin{description}[before={\microtypesetup{expansion=false}\setlength{\RaggedRightRightskip}{0pt plus 4em}\RaggedRight}]
\item[{[1]}] Управление по контролю качества пищевых продуктов и лекарственных средств США (FDA). \textit{M12: исследования лекарственных взаимодействий. Руководство для промышленности}. Окончательная редакция, август 2024 года. ICH M12; документ FDA 161199. Разделы об ингибировании транспортеров, клинической интерпретации и оценке на основе моделей. \href{https://www.fda.gov/media/161199/download}{Официальный источник}.
\item[{[2]}] Европейское агентство по лекарственным средствам (EMA). \textit{Руководство ICH M12 по исследованиям лекарственных взаимодействий --- этап 5}; \textit{Стратегия внедрения руководства ICH M12 по исследованиям лекарственных взаимодействий}. EMA/CHMP/ICH/652460/2022; EMA/460496/2024. Дата вступления в силу: 30 ноября 2024 года; уведомление о внедрении от 2 декабря 2024 года размещено 13 февраля 2025 года. \href{https://www.ema.europa.eu/en/ich-m12-drug-interaction-studies-scientific-guideline}{Официальный источник}. \href{https://www.ema.europa.eu/en/documents/scientific-guideline/implementation-strategy-ich-guideline-m12-drug-interaction-studies_en.pdf}{Стратегия внедрения}.
\item[{[3]}] AstraZeneca. \textit{Crestor (розувастатин): инструкция по медицинскому применению в США}. Редакция от апреля 2026 года; обновление DailyMed от 29 апреля 2026 года. Разделы 2.5--2.6, 5.1, 5.4, 7.1 и 12.3. Идентификатор набора DailyMed: 325a5d0e-9a72-4015-9fcd-1655fb504cee. \href{https://dailymed.nlm.nih.gov/dailymed/drugInfo.cfm?setid=325a5d0e-9a72-4015-9fcd-1655fb504cee}{Официальный источник}.
\item[{[4]}] Национальные институты здоровья США (NIH). \textit{Межлекарственные взаимодействия: ингибиторы протеазы и другие препараты}. Рекомендации по применению антиретровирусных средств у взрослых и подростков с ВИЧ-инфекцией. Таблица 24a; обновление от 12 сентября 2024 года; дата обращения: 8 сентября 2026 года. \href{https://clinicalinfo.hiv.gov/en/guidelines/hiv-clinical-guidelines-adult-and-adolescent-arv/drug-interactions-protease-inhibitors}{Официальный источник}.
\item[{[5]}] Национальные институты здоровья США (NIH). \textit{Обновленная информация о терапии статинами для первичной профилактики атеросклеротических сердечно-сосудистых заболеваний у людей с ВИЧ-инфекцией}. Заявление опубликовано 27 мая 2026 года. \href{https://clinicalinfo.hiv.gov/en/guidelines/hiv-clinical-guidelines-adult-and-adolescent-arv/statin-therapy-people-hiv}{Официальный источник}.
\item[{[6]}] Ciuta AD, Nosol K, Kowal J и др. Структура транспортеров лекарственных средств человека OATP1B1 и OATP1B3. \textit{\foreignlanguage{english}{Nature Communications}}. 2023;14:5774. DOI: \href{https://doi.org/10.1038/s41467-023-41552-8}{\nolinkurl{10.1038/s41467-023-41552-8}}. PMID: \href{https://pubmed.ncbi.nlm.nih.gov/37723174/}{37723174}.
\item[{[7]}] Shan Z, Yang X, Liu H и др. Криоэлектронно-микроскопические структуры полипептида человека OATP1B1, транспортирующего органические анионы. \textit{\foreignlanguage{english}{Cell Research}}. 2023;33:940--951. DOI: \href{https://doi.org/10.1038/s41422-023-00870-8}{\nolinkurl{10.1038/s41422-023-00870-8}}. PMID: \href{https://pubmed.ncbi.nlm.nih.gov/37674011/}{37674011}.
\item[{[8]}] van de Steeg E, Stranecky V, Hartmannova H и др. Полный дефицит OATP1B1 и OATP1B3 вызывает синдром Ротора у человека, нарушая повторный захват конъюгированного билирубина печенью. \textit{\foreignlanguage{english}{Journal of Clinical Investigation}}. 2012;122:519--528. DOI: \href{https://doi.org/10.1172/JCI59526}{\nolinkurl{10.1172/JCI59526}}. PMID: \href{https://pubmed.ncbi.nlm.nih.gov/22232210/}{22232210}.
\item[{[9]}] Orlando BJ, Liao M. ABCG2 переносит противоопухолевые препараты посредством переключения из закрытого состояния в открытое. \textit{\foreignlanguage{english}{Nature Communications}}. 2020;11:2264. DOI: \href{https://doi.org/10.1038/s41467-020-16155-2}{\nolinkurl{10.1038/s41467-020-16155-2}}. PMID: \href{https://pubmed.ncbi.nlm.nih.gov/32385283/}{32385283}.
\item[{[10]}] Woodward OM, Kottgen A, Coresh J и др. Выявление транспортера урата ABCG2 с распространенным функциональным полиморфизмом, вызывающим подагру. \textit{\foreignlanguage{english}{Proceedings of the National Academy of Sciences USA}}. 2009;106:10338--10342. DOI: \href{https://doi.org/10.1073/pnas.0901249106}{\nolinkurl{10.1073/pnas.0901249106}}. PMID: \href{https://pubmed.ncbi.nlm.nih.gov/19506252/}{19506252}.
\item[{[11]}] Volk EL, Schneider E. Белок резистентности рака молочной железы дикого типа (BCRP/ABCG2) является транспортером полиглутаматов метотрексата. \textit{\foreignlanguage{english}{Cancer Research}}. 2003;63:5538--5543. PMID: \href{https://pubmed.ncbi.nlm.nih.gov/14500392/}{14500392}.
\item[{[12]}] Elsby R, Coghlan H, Edgerton J и др. Механистические исследования in vitro показывают, что клинические взаимодействия ингибиторов протеазы с розувастатином обусловлены ингибированием кишечного BCRP и печеночного OATP1B1 при минимальном вкладе OATP1B3, NTCP и OAT3. \textit{\foreignlanguage{english}{Pharmacology Research \& Perspectives}}. 2023;11:e01060. DOI: \href{https://doi.org/10.1002/prp2.1060}{\nolinkurl{10.1002/prp2.1060}}. PMID: \href{https://pubmed.ncbi.nlm.nih.gov/36811234/}{36811234}.
\item[{[13]}] Kosloski MP, Bow DAJ, Kikuchi R и др. Перенос результатов исследований ингибирования транспорта in vitro на клинические межлекарственные взаимодействия глекапревира и пибрентасвира. \textit{\foreignlanguage{english}{Journal of Pharmacology and Experimental Therapeutics}}. 2019;370:278--287. DOI: \href{https://doi.org/10.1124/jpet.119.256966}{\nolinkurl{10.1124/jpet.119.256966}}. PMID: \href{https://pubmed.ncbi.nlm.nih.gov/31167814/}{31167814}.
\item[{[14]}] AbbVie. \textit{Mavyret (глекапревир/пибрентасвир): инструкция по медицинскому применению в США}. Редакция от июня 2025 года. Разделы 7 и 12.3; таблицы доз и экспозиции в исследованиях взаимодействий. \href{https://www.rxabbvie.com/pdf/mavyret_pi.pdf}{Официальный источник}.
\item[{[15]}] Европейское агентство по лекарственным средствам (EMA). \textit{Maviret: информация о препарате в составе европейского публичного оценочного отчета (EPAR)}. EMEA/H/C/004430. Запись с информацией о препарате обновлена 15 июля 2026 года. Разделы SmPC 4.3 и 4.5. \href{https://www.ema.europa.eu/en/medicines/human/EPAR/maviret}{Официальный источник}.
\item[{[16]}] Gilead Sciences. \textit{Epclusa (софосбувир/велпатасвир): инструкция по медицинскому применению в США}. Редакция от апреля 2022 года; версия на сайте производителя просмотрена 8 сентября 2026 года. Разделы 7 и 12.3. \href{https://www.gilead.com/-/media/files/pdfs/medicines/liver-disease/epclusa/epclusa_pi.pdf}{Официальный источник}.
\item[{[17]}] Gilead Sciences. \textit{Vosevi (софосбувир/\allowbreak велпатасвир/\allowbreak воксилапревир): инструкция по медицинскому применению в США}. Размещенная на сайте производителя редакция от ноября 2019 года; дата обращения: 8 сентября 2026 года. Разделы 7 и 12.3, включая исследования взаимодействий с дополнительным воксилапревиром. \href{https://www.gilead.com/-/media/files/pdfs/medicines/liver-disease/vosevi/vosevi_pi.pdf}{Официальный источник}.
\item[{[18]}] Европейское агентство по лекарственным средствам (EMA). \textit{Vosevi: информация о препарате в составе EPAR}. EMEA/H/C/004350. Запись с информацией о препарате обновлена 26 июля 2024 года. Разделы SmPC 4.3 и 4.5. \href{https://www.ema.europa.eu/en/medicines/human/EPAR/vosevi}{Официальный источник}.
\item[{[19]}] Ozvegy-Laczka C, Hegedus T, Varady G и др. Высокоаффинное взаимодействие ингибиторов тирозинкиназ с транспортером множественной лекарственной устойчивости ABCG2. \textit{\foreignlanguage{english}{Molecular Pharmacology}}. 2004;65:1485--1495. DOI: \href{https://doi.org/10.1124/mol.65.6.1485}{\nolinkurl{10.1124/mol.65.6.1485}}. PMID: \href{https://pubmed.ncbi.nlm.nih.gov/15155841/}{15155841}.
\item[{[20]}] Zurth C, Koskinen M, Fricke R и др. Потенциал межлекарственных взаимодействий даролутамида: исследования in vitro и клинические исследования. \textit{\foreignlanguage{english}{European Journal of Drug Metabolism and Pharmacokinetics}}. 2019;44:747--759. DOI: \href{https://doi.org/10.1007/s13318-019-00577-5}{\nolinkurl{10.1007/s13318-019-00577-5}}. PMID: \href{https://pubmed.ncbi.nlm.nih.gov/31571146/}{31571146}.
\item[{[21]}] Bayer. \textit{Stivarga (регорафениб): инструкция по медицинскому применению в США}. Версия на сайте производителя с обновлениями от февраля 2026 года. Разделы 7 и 12.3; исследование взаимодействия с субстратом BCRP. \href{https://dailymed.nlm.nih.gov/dailymed/drugInfo.cfm?setid=824f19c9-0546-4a8a-8d8f-c4055c04f7c7}{Официальный источник}.
\item[{[22]}] Novartis. \textit{Tabrecta (капматиниб): инструкция по медицинскому применению в США}. Версия на сайте производителя просмотрена 8 сентября 2026 года. Разделы 7.2 и 12.3. \href{https://www.novartis.com/us-en/sites/novartis_us/files/tabrecta.pdf}{Официальный источник}.
\item[{[23]}] Ozvegy-Laczka C, Laczko R, Hegedus C и др. Взаимодействие с моноклональным антителом 5D3 регулируется внутримолекулярными перестройками, но не образованием ковалентного димера транспортера множественной лекарственной устойчивости человека ABCG2. \textit{\foreignlanguage{english}{Journal of Biological Chemistry}}. 2008;283:26059--26070. DOI: \href{https://doi.org/10.1074/jbc.M803230200}{\nolinkurl{10.1074/jbc.M803230200}}. PMID: \href{https://pubmed.ncbi.nlm.nih.gov/18644784/}{18644784}.
\item[{[24]}] Управление по контролю качества пищевых продуктов и лекарственных средств США (FDA). \textit{Оценка межлекарственных взаимодействий терапевтических белков: руководство для промышленности}. Июнь 2023 года. Документ FDA 140909. \href{https://www.fda.gov/media/140909/download}{Официальный источник}.
\item[{[25]}] Simonson SG, Raza A, Martin PD и др. Фармакокинетика розувастатина у реципиентов сердца, получающих иммуносупрессивную терапию с циклоспорином. \textit{\foreignlanguage{english}{Clinical Pharmacology \& Therapeutics}}. 2004;76:167--177. DOI: \href{https://doi.org/10.1016/j.clpt.2004.03.010}{\nolinkurl{10.1016/j.clpt.2004.03.010}}. PMID: \href{https://pubmed.ncbi.nlm.nih.gov/15289793/}{15289793}.
\item[{[26]}] Schneck DW, Birmingham BK, Zalikowski JA и др. Влияние гемфиброзила на фармакокинетику розувастатина. \textit{\foreignlanguage{english}{Clinical Pharmacology \& Therapeutics}}. 2004;75:455--463. DOI: \href{https://doi.org/10.1016/j.clpt.2003.12.014}{\nolinkurl{10.1016/j.clpt.2003.12.014}}. PMID: \href{https://pubmed.ncbi.nlm.nih.gov/15116058/}{15116058}.
\item[{[27]}] Kiser JJ, Gerber JG, Predhomme JA и др. Межлекарственное взаимодействие лопинавира/ритонавира и розувастатина у здоровых добровольцев. \textit{\foreignlanguage{english}{Journal of Acquired Immune Deficiency Syndromes}}. 2008;47:570--578. DOI: \href{https://doi.org/10.1097/QAI.0b013e318160a542}{\nolinkurl{10.1097/QAI.0b013e318160a542}}. PMID: \href{https://pubmed.ncbi.nlm.nih.gov/18176327/}{18176327}.
\item[{[28]}] Lehtisalo M, Tarkiainen EK, Neuvonen M и др. Тикагрелор увеличивает экспозицию розувастатина --- субстрата белка резистентности рака молочной железы. \textit{\foreignlanguage{english}{Clinical Pharmacology \& Therapeutics}}. 2024;115:71--79; впервые опубликовано в интернете 3 октября 2023 года. DOI: \href{https://doi.org/10.1002/cpt.3067}{\nolinkurl{10.1002/cpt.3067}}. PMID: \href{https://pubmed.ncbi.nlm.nih.gov/37786998/}{37786998}.
\item[{[29]}] Dermota T, Jug B, Trontelj J, Bozic Mijovski M. Тикагрелор ассоциирован с повышенными концентрациями розувастатина в крови у пациентов, перенесших инфаркт миокарда. \textit{\foreignlanguage{english}{Clinical Pharmacokinetics}}. 2025;64:565--571. DOI: \href{https://doi.org/10.1007/s40262-025-01489-1}{\nolinkurl{10.1007/s40262-025-01489-1}}. PMID: \href{https://pubmed.ncbi.nlm.nih.gov/40029501/}{40029501}.
\item[{[30]}] \textit{COLCRYS (колхицин): инструкция по медицинскому применению в США}. DailyMed, дата обращения: 8 сентября 2026 года. Разделы 2.4, 4, 5, 7 и 12.3. Идентификатор набора DailyMed: 56c130d1-7581-4152-99a2-0014ee9366c0. \href{https://dailymed.nlm.nih.gov/dailymed/drugInfo.cfm?setid=56c130d1-7581-4152-99a2-0014ee9366c0}{Официальный источник}.
\item[{[31]}] Wang Q, Zheng M, Leil T. Исследование транспортер-опосредованных межлекарственных взаимодействий с использованием физиологически обоснованной фармакокинетической модели розувастатина. \textit{\foreignlanguage{english}{CPT: Pharmacometrics \& Systems Pharmacology}}. 2017;6:228--238. DOI: \href{https://doi.org/10.1002/psp4.12168}{\nolinkurl{10.1002/psp4.12168}}. PMID: \href{https://pubmed.ncbi.nlm.nih.gov/28296193/}{28296193}.
\item[{[32]}] Shin JI, Fine DM, Sang Y и др. Связь применения розувастатина с риском гематурии и протеинурии. \textit{\foreignlanguage{english}{Journal of the American Society of Nephrology}}. 2022;33:1767--1777. DOI: \href{https://doi.org/10.1681/ASN.2022020135}{\nolinkurl{10.1681/ASN.2022020135}}. PMID: \href{https://pubmed.ncbi.nlm.nih.gov/35853713/}{35853713}. PMCID: \href{https://pmc.ncbi.nlm.nih.gov/articles/PMC9529194/}{PMC9529194}.
\item[{[33]}] German P, Liu HC, Szwarcberg J и др. Влияние кобицистата на скорость клубочковой фильтрации у лиц с нормальной и нарушенной функцией почек. \textit{\foreignlanguage{english}{Journal of Acquired Immune Deficiency Syndromes}}. 2012;61:32--40. DOI: \href{https://doi.org/10.1097/QAI.0b013e3182645648}{\nolinkurl{10.1097/QAI.0b013e3182645648}}. PMID: \href{https://pubmed.ncbi.nlm.nih.gov/22732469/}{22732469}.
\item[{[34]}] Cooper-DeHoff RM, Niemi M, Ramsey LB и др. Рекомендации Консорциума по внедрению клинической фармакогенетики по генотипам SLCO1B1, ABCG2 и CYP2C9 и статин-ассоциированным скелетно-мышечным симптомам. \textit{\foreignlanguage{english}{Clinical Pharmacology \& Therapeutics}}. 2022;111:1007--1021. DOI: \href{https://doi.org/10.1002/cpt.2557}{\nolinkurl{10.1002/cpt.2557}}. PMID: \href{https://pubmed.ncbi.nlm.nih.gov/35152405/}{35152405}. PMCID: \href{https://pmc.ncbi.nlm.nih.gov/articles/PMC9035072/}{PMC9035072}. Комбинированные рекомендации для розувастатина: таблица 5.
\item[{[35]}] Исследовательская группа SEARCH. Варианты SLCO1B1 и статин-индуцированная миопатия: полногеномное исследование. \textit{\foreignlanguage{english}{New England Journal of Medicine}}. 2008;359:789--799. DOI: \href{https://doi.org/10.1056/NEJMoa0801936}{\nolinkurl{10.1056/NEJMoa0801936}}. PMID: \href{https://pubmed.ncbi.nlm.nih.gov/18650507/}{18650507}.
\item[{[36]}] Oni-Orisan A, Yee SW, Haldar T и др. Обновление таблицы функциональных характеристик аллелей SLCO1B1 Консорциума по внедрению клинической фармакогенетики (CPIC) с использованием данных участников преимущественно африканского происхождения из регионов к югу от Сахары. \textit{\foreignlanguage{english}{Clinical Pharmacology \& Therapeutics}}. 2026;119:1133--1135; впервые опубликовано в интернете 1 декабря 2025 года. DOI: \href{https://doi.org/10.1002/cpt.70153}{\nolinkurl{10.1002/cpt.70153}}. PMID: \href{https://pubmed.ncbi.nlm.nih.gov/41321195/}{41321195}. PMCID: \href{https://pmc.ncbi.nlm.nih.gov/articles/PMC13055011/}{PMC13055011}.
\item[{[37]}] Mori D, Kimoto E, Rago B и др. Дозозависимое ингибирование OATP1B рифампицином у здоровых добровольцев: комплексная оценка кандидатных биомаркеров и зондовых препаратов OATP1B. \textit{\foreignlanguage{english}{Clinical Pharmacology \& Therapeutics}}. 2020;107:1004--1013; впервые опубликовано в интернете в 2019 году. DOI: \href{https://doi.org/10.1002/cpt.1695}{\nolinkurl{10.1002/cpt.1695}}. PMID: \href{https://pubmed.ncbi.nlm.nih.gov/31628668/}{31628668}.
\item[{[38]}] Управление по контролю качества пищевых продуктов и лекарственных средств США (FDA). \textit{Для медицинских специалистов: перечень FDA с примерами лекарственных средств, взаимодействующих с ферментами CYP и транспортерными системами}. Дата обращения: 8 сентября 2026 года. \href{https://www.fda.gov/drugs/drug-interactions-labeling/healthcare-professionals-fdas-examples-drugs-interact-cyp-enzymes-and-transporter-systems}{Официальный источник}.
\item[{[39]}] Ogilvie BW, Zhang D, Li W и др. Глюкуронирование превращает гемфиброзил в мощный метаболизм-зависимый ингибитор CYP2C8: значение для межлекарственных взаимодействий. \textit{\foreignlanguage{english}{Drug Metabolism and Disposition}}. 2006;34:191--197. DOI: \href{https://doi.org/10.1124/dmd.105.007633}{\nolinkurl{10.1124/dmd.105.007633}}. PMID: \href{https://pubmed.ncbi.nlm.nih.gov/16299161/}{16299161}.
\item[{[40]}] Европейское агентство по лекарственным средствам (EMA). \textit{Vosevi: европейский публичный оценочный отчет EPAR}. EMA/441550/2017; EMEA/H/C/004350. 2017; впервые размещен на сайте EMA 22 сентября 2017 года. Обсуждение тактики ведения при взаимодействиях розувастатина и европейского противопоказания сочетания циклоспорина с розувастатином; действующая национальная SmPC розувастатина сохраняет приоритет. \href{https://www.ema.europa.eu/documents/assessment-report/vosevi-epar-public-assessment-report_en.pdf}{Официальный источник}.
\item[{[41]}] Maliepaard M, Scheffer GL, Faneyte IF и др. Субклеточная локализация и распределение транспортера BCRP в нормальных тканях человека. \textit{\foreignlanguage{english}{Cancer Research}}. 2001;61:3458--3464. PMID: \href{https://pubmed.ncbi.nlm.nih.gov/11309308/}{11309308}.
\item[{[42]}] Ho RH, Tirona RG, Leake BF и др. Транспортеры лекарственных средств и желчных кислот в печеночном захвате розувастатина: функция, экспрессия и фармакогенетика. \textit{\foreignlanguage{english}{Gastroenterology}}. 2006;130:1793--1806. DOI: \href{https://doi.org/10.1053/j.gastro.2006.02.034}{\nolinkurl{10.1053/j.gastro.2006.02.034}}. PMID: \href{https://pubmed.ncbi.nlm.nih.gov/16697742/}{16697742}.
\item[{[43]}] Billington S, Shoner S, Lee S и др. ПЭТ-визуализация печеночных концентраций и гепатобилиарного транспорта \([^{11}\mathrm{C}]\)розувастатина у человека в отсутствие и в присутствии циклоспорина A. \textit{\foreignlanguage{english}{Clinical Pharmacology \& Therapeutics}}. 2019;106:1056--1066. DOI: \href{https://doi.org/10.1002/cpt.1506}{\nolinkurl{10.1002/cpt.1506}}. PMID: \href{https://pubmed.ncbi.nlm.nih.gov/31102467/}{31102467}.
\item[{[44]}] Gui C, Hagenbuch B. Аминокислотные остатки десятого трансмембранного домена полипептида OATP1B3, транспортирующего органические анионы, критически важны для транспорта холецистокинина-октапептида. \textit{\foreignlanguage{english}{Biochemistry}}. 2008;47:9090--9097. DOI: \href{https://doi.org/10.1021/bi8008455}{\nolinkurl{10.1021/bi8008455}}. PMID: \href{https://pubmed.ncbi.nlm.nih.gov/18690707/}{18690707}.
\item[{[45]}] Управление по регулированию медицинской продукции Ирландии (HPRA). \textit{Rosuvastatin Viatris 5 мг, таблетки, покрытые пленочной оболочкой: общая характеристика лекарственного препарата}. Историческая редакция от 21 марта 2025 года. Регистрационный номер PA23266/003/001; CRN00DG4K. Разделы 4.2--4.5. Дата обращения: 8 сентября 2026 года. \href{https://www.hpra.ie/find-a-medicine/for-human-use/authorised-medicines/details/18914}{Официальная карточка препарата}.
\item[{[46]}] Keskitalo JE, Zolk O, Fromm MF, Kurkinen KJ, Neuvonen PJ, Niemi M. Полиморфизм ABCG2 существенно влияет на фармакокинетику аторвастатина и розувастатина. \textit{\foreignlanguage{english}{Clinical Pharmacology \& Therapeutics}}. 2009;86:197--203. DOI: \href{https://doi.org/10.1038/clpt.2009.79}{\nolinkurl{10.1038/clpt.2009.79}}. PMID: \href{https://pubmed.ncbi.nlm.nih.gov/19474787/}{19474787}.
\item[{[47]}] Kadmon Pharmaceuticals. \textit{REZUROCK (белумосудил): инструкция по медицинскому применению в США}. DailyMed, версия 11, опубликована 9 апреля 2026 года; дата обращения: 8 сентября 2026 года. Разделы 7 и 12.1--12.3. Идентификатор набора: 102e4ef4-7f84-4e34-8df1-479c24d1575d. \href{https://dailymed.nlm.nih.gov/dailymed/drugInfo.cfm?setid=102e4ef4-7f84-4e34-8df1-479c24d1575d}{Официальный источник}.
\item[{[48]}] \textit{OJJAARA (момелотиниб): инструкция по медицинскому применению в США}. DailyMed, дата обращения: 8 сентября 2026 года. Разделы 7 и 12.3. Идентификатор набора: 4a672fd1-eb7b-4aa9-9b23-845e4b5dc400. \href{https://dailymed.nlm.nih.gov/dailymed/drugInfo.cfm?setid=4a672fd1-eb7b-4aa9-9b23-845e4b5dc400}{Официальный источник}.
\item[{[49]}] \textit{VYNDAQEL/VYNDAMAX (тафамидис): инструкция по медицинскому применению в США}. DailyMed, дата обращения: 8 сентября 2026 года. Разделы 7 и 12.3. Идентификатор набора: 1b4121ee-a733-4456-a917-be2603477839. \href{https://dailymed.nlm.nih.gov/dailymed/drugInfo.cfm?setid=1b4121ee-a733-4456-a917-be2603477839}{Официальный источник}.
\item[{[50]}] Lehtisalo M, Keskitalo JE, Tornio A и др. Фебуксостат, но не аллопуринол, значительно повышает плазменные концентрации розувастатина --- субстрата белка резистентности рака молочной железы. \textit{\foreignlanguage{english}{Clinical and Translational Science}}. 2020;13:1236--1243. DOI: \href{https://doi.org/10.1111/cts.12809}{\nolinkurl{10.1111/cts.12809}}. PMID: \href{https://pubmed.ncbi.nlm.nih.gov/32453913/}{32453913}.
\item[{[51]}] Pasanen MK, Fredrikson H, Neuvonen PJ, Niemi M. Различное влияние полиморфизма SLCO1B1 на фармакокинетику аторвастатина и розувастатина. \textit{\foreignlanguage{english}{Clinical Pharmacology \& Therapeutics}}. 2007;82:726--733. DOI: \href{https://doi.org/10.1038/sj.clpt.6100220}{\nolinkurl{10.1038/sj.clpt.6100220}}. PMID: \href{https://pubmed.ncbi.nlm.nih.gov/17473846/}{17473846}.
\item[{[52]}] Bednarczyk D, Boiselle C. Транспорт копропорфиринов I и III, опосредованный полипептидами OATP. \textit{\foreignlanguage{english}{Xenobiotica}}. 2016;46:457--466. DOI: \href{https://doi.org/10.3109/00498254.2015.1085111}{\nolinkurl{10.3109/00498254.2015.1085111}}. PMID: \href{https://pubmed.ncbi.nlm.nih.gov/26383540/}{26383540}.
\item[{[53]}] Neuvonen M, Tornio A, Hirvensalo P, Backman JT, Niemi M. Характеристики плазменных копропорфиринов I и III как биомаркеров OATP1B1 у человека. \textit{\foreignlanguage{english}{Clinical Pharmacology \& Therapeutics}}. 2021;110:1622--1632. DOI: \href{https://doi.org/10.1002/cpt.2429}{\nolinkurl{10.1002/cpt.2429}}. PMID: \href{https://pubmed.ncbi.nlm.nih.gov/34580865/}{34580865}.
\item[{[54]}] Niemi M, Backman JT, Neuvonen M, Neuvonen PJ. Влияние гемфиброзила, итраконазола и их сочетания на фармакокинетику и фармакодинамику репаглинида: потенциально опасное взаимодействие гемфиброзила и репаглинида. \textit{\foreignlanguage{english}{Diabetologia}}. 2003;46:347--351. DOI: \href{https://doi.org/10.1007/s00125-003-1034-7}{\nolinkurl{10.1007/s00125-003-1034-7}}. PMID: \href{https://pubmed.ncbi.nlm.nih.gov/12687332/}{12687332}.
\item[{[55]}] Kudo T, Hisaka A, Sugiyama Y, Ito K. Анализ увеличения концентрации репаглинида под действием гемфиброзила и итраконазола с учетом ингибирования печеночного захвата и метаболических ферментов. \textit{\foreignlanguage{english}{Drug Metabolism and Disposition}}. 2013;41:362--371. DOI: \href{https://doi.org/10.1124/dmd.112.049460}{\nolinkurl{10.1124/dmd.112.049460}}. PMID: \href{https://pubmed.ncbi.nlm.nih.gov/23139378/}{23139378}.
\item[{[56]}] \textit{Repaglinide tablets: инструкция по медицинскому применению в США}. DailyMed; разделы 2.3, 4, 5.1, 7 и 12.3. Идентификатор набора: d7673e22-0dbb-4395-ad60-943d249203a1. Дата обращения: 9 сентября 2026 года. \href{https://dailymed.nlm.nih.gov/dailymed/drugInfo.cfm?setid=d7673e22-0dbb-4395-ad60-943d249203a1}{Официальный источник}.
\item[{[57]}] Chu XY, Bleasby K, Yabut J и др. Транспорт ингибитора DPP-4 ситаглиптина посредством OAT3, OATP4C1 и P-гликопротеина человека. \textit{\foreignlanguage{english}{Journal of Pharmacology and Experimental Therapeutics}}. 2007;321:673--683. DOI: \href{https://doi.org/10.1124/jpet.106.116517}{\nolinkurl{10.1124/jpet.106.116517}}. PMID: \href{https://pubmed.ncbi.nlm.nih.gov/17314201/}{17314201}.
\item[{[58]}] Kiander W, Sinokki A, Neuvonen M, Kidron H, Niemi M. Сравнительный захват статинов печеночными полипептидами, транспортирующими органические анионы. \textit{\foreignlanguage{english}{European Journal of Pharmaceutical Sciences}}. 2025;209:107073. DOI: \href{https://doi.org/10.1016/j.ejps.2025.107073}{\nolinkurl{10.1016/j.ejps.2025.107073}}. PMID: \href{https://pubmed.ncbi.nlm.nih.gov/40107569/}{40107569}.
\item[{[59]}] Kunze A, Huwyler J, Camenisch G, Poller B. Прогнозирование OATP1B1- и OATP1B3-опосредованного печеночного захвата статинов по данным экспрессии и активности транспортеров. \textit{\foreignlanguage{english}{Drug Metabolism and Disposition}}. 2014;42:1514--1521. DOI: \href{https://doi.org/10.1124/dmd.114.058412}{\nolinkurl{10.1124/dmd.114.058412}}. PMID: \href{https://pubmed.ncbi.nlm.nih.gov/24989890/}{24989890}.
\item[{[60]}] Frost CE, Byon W, Song Y и др. Влияние кетоконазола и дилтиазема на фармакокинетику апиксабана, перорального прямого ингибитора фактора Xa. \textit{\foreignlanguage{english}{British Journal of Clinical Pharmacology}}. 2015;79:838--846. DOI: \href{https://doi.org/10.1111/bcp.12541}{\nolinkurl{10.1111/bcp.12541}}. PMID: \href{https://pubmed.ncbi.nlm.nih.gov/25377242/}{25377242}.
\item[{[61]}] Mueck W, Kubitza D, Becka M. Совместное применение ривароксабана с препаратами, использующими общие пути элиминации: фармакокинетические эффекты у здоровых добровольцев. \textit{\foreignlanguage{english}{British Journal of Clinical Pharmacology}}. 2013;76:455--466. DOI: \href{https://doi.org/10.1111/bcp.12075}{\nolinkurl{10.1111/bcp.12075}}. PMID: \href{https://pubmed.ncbi.nlm.nih.gov/23305158/}{23305158}.
\item[{[62]}] Seithel A, Eberl S, Singer K и др. Влияние макролидных антибиотиков на OATP1B1- и OATP1B3-опосредованный захват органических анионов и лекарств. \textit{\foreignlanguage{english}{Drug Metabolism and Disposition}}. 2007;35:779--786. DOI: \href{https://doi.org/10.1124/dmd.106.014407}{\nolinkurl{10.1124/dmd.106.014407}}. PMID: \href{https://pubmed.ncbi.nlm.nih.gov/17296622/}{17296622}.
\item[{[63]}] Patel AM, Shariff S, Bailey DG и др. Токсичность статинов при совместном назначении макролидов: популяционное когортное исследование. \textit{\foreignlanguage{english}{Annals of Internal Medicine}}. 2013;158:869--876. DOI: \href{https://doi.org/10.7326/0003-4819-158-12-201306180-00004}{\nolinkurl{10.7326/0003-4819-158-12-201306180-00004}}. PMID: \href{https://pubmed.ncbi.nlm.nih.gov/23778904/}{23778904}.
\item[{[64]}] Pahwa S, Alam K, Crowe A и др. Предварительная инкубация с рифампицином и дазатинибом усиливает ингибирование OATP1B1- и OATP1B3-опосредованного транспорта. \textit{\foreignlanguage{english}{Journal of Pharmaceutical Sciences}}. 2017;106:2123--2135. DOI: \href{https://doi.org/10.1016/j.xphs.2017.03.022}{\nolinkurl{10.1016/j.xphs.2017.03.022}}. PMCID: \href{https://pmc.ncbi.nlm.nih.gov/articles/PMC5511785/}{PMC5511785}. PMID: \href{https://pubmed.ncbi.nlm.nih.gov/28373111/}{28373111}.
\item[{[65]}] Bidstrup TB, Stilling N, Damkier P и др. Рифампицин, по-видимому, действует и как индуктор, и как ингибитор метаболизма репаглинида. \textit{\foreignlanguage{english}{European Journal of Clinical Pharmacology}}. 2004;60:109--114. DOI: \href{https://doi.org/10.1007/s00228-004-0746-z}{\nolinkurl{10.1007/s00228-004-0746-z}}. PMID: \href{https://pubmed.ncbi.nlm.nih.gov/15034704/}{15034704}.
\item[{[66]}] Kalliokoski A, Backman JT, Kurkinen KJ, Neuvonen PJ, Niemi M. Влияние гемфиброзила и аторвастатина на фармакокинетику репаглинида в зависимости от полиморфизма SLCO1B1. \textit{\foreignlanguage{english}{Clinical Pharmacology \& Therapeutics}}. 2008;84:488--496. DOI: \href{https://doi.org/10.1038/clpt.2008.74}{\nolinkurl{10.1038/clpt.2008.74}}. PMID: \href{https://pubmed.ncbi.nlm.nih.gov/19238654/}{19238654}.
\item[{[67]}] van de Steeg E, Greupink R, Schreurs M и др. Взаимодействия розувастатина и пероральных антидиабетических препаратов на уровне OATP1B1. \textit{\foreignlanguage{english}{Drug Metabolism and Disposition}}. 2013;41:592--601. DOI: \href{https://doi.org/10.1124/dmd.112.049023}{\nolinkurl{10.1124/dmd.112.049023}}. PMID: \href{https://pubmed.ncbi.nlm.nih.gov/23248200/}{23248200}.
\item[{[68]}] Lempers VJC, van den Heuvel JJMW, Russel FGM и др. Ингибирующее действие противогрибковых препаратов на ABC-транспортеры P-gp, MRP1--MRP5, BCRP и BSEP. \textit{\foreignlanguage{english}{Antimicrobial Agents and Chemotherapy}}. 2016;60:3372--3379. DOI: \href{https://doi.org/10.1128/AAC.02931-15}{\nolinkurl{10.1128/AAC.02931-15}}. PMID: \href{https://pubmed.ncbi.nlm.nih.gov/27001813/}{27001813}.
\item[{[69]}] Cooper KJ, Martin PD, Dane AL и др. Влияние итраконазола на фармакокинетику розувастатина. \textit{\foreignlanguage{english}{Clinical Pharmacology \& Therapeutics}}. 2003;73:322--329. DOI: \href{https://doi.org/10.1016/S0009-9236(02)17633-8}{\nolinkurl{10.1016/S0009-9236(02)17633-8}}. PMID: \href{https://pubmed.ncbi.nlm.nih.gov/12709722/}{12709722}.
\item[{[70]}] Cooper KJ, Martin PD, Dane AL и др. Отсутствие влияния кетоконазола на фармакокинетику розувастатина у здоровых добровольцев. \textit{\foreignlanguage{english}{British Journal of Clinical Pharmacology}}. 2003;55:94--99. DOI: \href{https://doi.org/10.1046/j.1365-2125.2003.01720.x}{\nolinkurl{10.1046/j.1365-2125.2003.01720.x}}. PMID: \href{https://pubmed.ncbi.nlm.nih.gov/12534645/}{12534645}.
\item[{[71]}] Breedveld P, Zelcer N, Pluim D и др. Механизм фармакокинетического взаимодействия метотрексата и бензимидазолов: возможная роль BCRP. \textit{\foreignlanguage{english}{Cancer Research}}. 2004;64:5804--5811. DOI: \href{https://doi.org/10.1158/0008-5472.CAN-03-4062}{\nolinkurl{10.1158/0008-5472.CAN-03-4062}}. PMID: \href{https://pubmed.ncbi.nlm.nih.gov/15313923/}{15313923}.
\item[{[72]}] Adkison KK, Vaidya SS, Lee DY и др. Пероральный сульфасалазин как клинический зонд BCRP: влияние варианта C421A и совместного приема пантопразола. \textit{\foreignlanguage{english}{Journal of Pharmaceutical Sciences}}. 2010;99:1046--1062. DOI: \href{https://doi.org/10.1002/jps.21860}{\nolinkurl{10.1002/jps.21860}}. PMID: \href{https://pubmed.ncbi.nlm.nih.gov/19569219/}{19569219}.
\item[{[73]}] Elsby R, Fox L, Stresser D и др. Оценка in vitro риска транспортер-опосредованного взаимодействия AZD9056 с метотрексатом. \textit{\foreignlanguage{english}{European Journal of Pharmaceutical Sciences}}. 2011;43:41--49. DOI: \href{https://doi.org/10.1016/j.ejps.2011.03.006}{\nolinkurl{10.1016/j.ejps.2011.03.006}}. PMID: \href{https://pubmed.ncbi.nlm.nih.gov/21440623/}{21440623}.
\item[{[74]}] Hu M, Lee HK, To KKW и др. Телмисартан увеличивает системную экспозицию розувастатина после однократного и повторного приема и ингибирует его ABCG2-опосредованный транспорт in vitro. \textit{\foreignlanguage{english}{European Journal of Clinical Pharmacology}}. 2016;72:1471--1478. DOI: \href{https://doi.org/10.1007/s00228-016-2130-1}{\nolinkurl{10.1007/s00228-016-2130-1}}. PMID: \href{https://pubmed.ncbi.nlm.nih.gov/27651239/}{27651239}.
\item[{[75]}] Bae SH, Park WS, Han S и др. PBPK-прогнозирование кишечных BCRP-опосредованных взаимодействий розувастатина у корейских участников. \textit{\foreignlanguage{english}{Korean Journal of Physiology \& Pharmacology}}. 2018;22:321--329. DOI: \href{https://doi.org/10.4196/kjpp.2018.22.3.321}{\nolinkurl{10.4196/kjpp.2018.22.3.321}}. PMID: \href{https://pubmed.ncbi.nlm.nih.gov/29719454/}{29719454}.
\item[{[76]}] Sane RS, Mitra P, Lai Y и др. Оценка клинического ингибирования транспортеров: результаты рабочей группы Международного консорциума инноваций и качества в фармацевтической разработке. \textit{\foreignlanguage{english}{Drug Metabolism and Disposition}}. 2026;54:100255; опубликовано в интернете 24 февраля 2026 года. DOI: \href{https://doi.org/10.1016/j.dmd.2026.100255}{\nolinkurl{10.1016/j.dmd.2026.100255}}. PMID: \href{https://pubmed.ncbi.nlm.nih.gov/41894930/}{41894930}.
\item[{[77]}] Zhang D, He K, Herbst JJ и др. Характеристика эффлюксных транспортеров, участвующих в распределении и элиминации апиксабана. \textit{\foreignlanguage{english}{Drug Metabolism and Disposition}}. 2013;41:827--835. DOI: \href{https://doi.org/10.1124/dmd.112.050260}{\nolinkurl{10.1124/dmd.112.050260}}. PMID: \href{https://pubmed.ncbi.nlm.nih.gov/23382458/}{23382458}.
\item[{[78]}] Sodhi JK, Liu S, Benet LZ. Кишечные эффлюксные транспортеры P-gp и BCRP не имеют клинической значимости для распределения и элиминации апиксабана. \textit{\foreignlanguage{english}{Pharmaceutical Research}}. 2020;37:208. DOI: \href{https://doi.org/10.1007/s11095-020-02927-4}{\nolinkurl{10.1007/s11095-020-02927-4}}. PMID: \href{https://pubmed.ncbi.nlm.nih.gov/32996065/}{32996065}.

\item[{[79]}] Roth M, Timmermann BN, Hagenbuch B. Взаимодействия катехинов зеленого чая с полипептидами, транспортирующими органические анионы. \textit{\foreignlanguage{english}{Drug Metabolism and Disposition}}. 2011;39:920--926. DOI: \href{https://doi.org/10.1124/dmd.110.036640}{\nolinkurl{10.1124/dmd.110.036640}}. PMCID: \href{https://pmc.ncbi.nlm.nih.gov/articles/PMC3082372/}{PMC3082372}. PMID: \href{https://pubmed.ncbi.nlm.nih.gov/21278283/}{21278283}.
\item[{[80]}] Chow HH, Hakim IA, Vining DR и др. Влияние условий приема на пероральную биодоступность катехинов зеленого чая после однократного введения Polyphenon E здоровым добровольцам. \textit{\foreignlanguage{english}{Clinical Cancer Research}}. 2005;11:4627--4633. DOI: \href{https://doi.org/10.1158/1078-0432.CCR-04-2549}{\nolinkurl{10.1158/1078-0432.CCR-04-2549}}. PMID: \href{https://pubmed.ncbi.nlm.nih.gov/15958649/}{15958649}.
\item[{[81]}] Kim TE, Ha N, Kim Y и др. Влияние эпигаллокатехин-3-галлата, основного компонента зеленого чая, на фармакокинетику розувастатина у здоровых добровольцев. \textit{\foreignlanguage{english}{Drug Design, Development and Therapy}}. 2017;11:1409--1416. DOI: \href{https://doi.org/10.2147/DDDT.S130050}{\nolinkurl{10.2147/DDDT.S130050}}. PMID: \href{https://pubmed.ncbi.nlm.nih.gov/28533679/}{28533679}.
\item[{[82]}] Zeng W, Hu M, Lee HK и др. Влияние экстракта зеленого чая и соевых изофлавонов на фармакокинетику розувастатина у здоровых добровольцев. \textit{\foreignlanguage{english}{Frontiers in Nutrition}}. 2022;9:850318. DOI: \href{https://doi.org/10.3389/fnut.2022.850318}{\nolinkurl{10.3389/fnut.2022.850318}}. PMCID: \href{https://pmc.ncbi.nlm.nih.gov/articles/PMC8987933/}{PMC8987933}. PMID: \href{https://pubmed.ncbi.nlm.nih.gov/35399656/}{35399656}.
\item[{[83]}] Bi X, Yuan Z, Qu B, Zhou H, Liu Z, Xie Y. Пиперин повышает биодоступность силибина посредством ингибирования эффлюксных транспортеров BCRP и MRP2. \textit{\foreignlanguage{english}{Phytomedicine}}. 2019;54:98--108. DOI: \href{https://doi.org/10.1016/j.phymed.2018.09.217}{\nolinkurl{10.1016/j.phymed.2018.09.217}}. PMID: \href{https://pubmed.ncbi.nlm.nih.gov/30668388/}{30668388}.
\item[{[84]}] Bhardwaj RK, Glaeser H, Becquemont L, Klotz U, Gupta SK, Fromm MF. Пиперин, основной компонент черного перца, ингибирует P-гликопротеин и CYP3A4 человека. \textit{\foreignlanguage{english}{Journal of Pharmacology and Experimental Therapeutics}}. 2002;302:645--650. DOI: \href{https://doi.org/10.1124/jpet.102.034728}{\nolinkurl{10.1124/jpet.102.034728}}. PMID: \href{https://pubmed.ncbi.nlm.nih.gov/12130727/}{12130727}.
\item[{[85]}] Kock K, Xie Y, Hawke RL, Oberlies NH, Brouwer KLR. Взаимодействие флавонолигнанов силимарина с полипептидами, транспортирующими органические анионы. \textit{\foreignlanguage{english}{Drug Metabolism and Disposition}}. 2013;41:958--965. DOI: \href{https://doi.org/10.1124/dmd.112.048272}{\nolinkurl{10.1124/dmd.112.048272}}. PMID: \href{https://pubmed.ncbi.nlm.nih.gov/23401473/}{23401473}. PMCID: \href{https://pmc.ncbi.nlm.nih.gov/articles/PMC3629808/}{PMC3629808}.
\item[{[86]}] Deng JW, Shon JH, Shin HJ и др. Влияние добавки силимарина на фармакокинетику розувастатина. \textit{\foreignlanguage{english}{Pharmaceutical Research}}. 2008;25:1807--1814. DOI: \href{https://doi.org/10.1007/s11095-007-9492-0}{\nolinkurl{10.1007/s11095-007-9492-0}}. PMID: \href{https://pubmed.ncbi.nlm.nih.gov/18236139/}{18236139}.
\item[{[87]}] Bajraktari-Sylejmani G, Weiss J. Потенциальный риск пищевых взаимодействий: полиметоксифлавоны и флаваноны цитрусовых как ингибиторы OATP1B1, OATP1B3 и OATP2B1. \textit{\foreignlanguage{english}{European Journal of Drug Metabolism and Pharmacokinetics}}. 2020;45:809--815. DOI: \href{https://doi.org/10.1007/s13318-020-00634-4}{\nolinkurl{10.1007/s13318-020-00634-4}}. PMID: \href{https://pubmed.ncbi.nlm.nih.gov/32661908/}{32661908}.
\item[{[88]}] Fleisher B, Unum J, Shao J, An G. Компоненты фруктовых соков взаимодействуют с дасатинибом посредством ингибирования BCRP: новый механизм взаимодействия напитков и лекарств. \textit{\foreignlanguage{english}{Journal of Pharmaceutical Sciences}}. 2015;104:266--275. DOI: \href{https://doi.org/10.1002/jps.24289}{\nolinkurl{10.1002/jps.24289}}. PMID: \href{https://pubmed.ncbi.nlm.nih.gov/25418056/}{25418056}.
\item[{[89]}] Paine MF, Widmer WW, Hart HL и др. Грейпфрутовый сок без фуранокумаринов подтверждает их роль во взаимодействии грейпфрутового сока и фелодипина. \textit{\foreignlanguage{english}{American Journal of Clinical Nutrition}}. 2006;83:1097--1105. DOI: \href{https://doi.org/10.1093/ajcn/83.5.1097}{\nolinkurl{10.1093/ajcn/83.5.1097}}. PMID: \href{https://pubmed.ncbi.nlm.nih.gov/16685052/}{16685052}.
\item[{[90]}] Lilja JJ, Kivisto KT, Neuvonen PJ. Грейпфрутовый сок увеличивает сывороточные концентрации аторвастатина и не влияет на правастатин. \textit{\foreignlanguage{english}{Clinical Pharmacology \& Therapeutics}}. 1999;66:118--127. DOI: \href{https://doi.org/10.1053/cp.1999.v66.100453001}{\nolinkurl{10.1053/cp.1999.v66.100453001}}. PMID: \href{https://pubmed.ncbi.nlm.nih.gov/10460065/}{10460065}.
\item[{[91]}] Yoshikawa M, Ikegami Y, Sano K и др. Транспорт SN-38 транспортером ABCG2 человека дикого типа и его ингибирование природным флавоноидом кверцетином. \textit{\foreignlanguage{english}{Journal of Experimental Therapeutics and Oncology}}. 2004;4:25--35. PMID: \href{https://pubmed.ncbi.nlm.nih.gov/15255290/}{15255290}.
\item[{[92]}] Gyongy Z, Mocsar G, Hegedus E и др. Связывание нуклеотидов как ключевой регулятор конформационных переходов ABCG2. \textit{\foreignlanguage{english}{eLife}}. 2023;12:e83976. DOI: \href{https://doi.org/10.7554/eLife.83976}{\nolinkurl{10.7554/eLife.83976}}. PMCID: \href{https://pmc.ncbi.nlm.nih.gov/articles/PMC9917445/}{PMC9917445}. PMID: \href{https://pubmed.ncbi.nlm.nih.gov/36763413/}{36763413}.
\item[{[93]}] Riha J, Brenner S, Bohmdorfer M и др. Ресвератрол и его основные сульфатные конъюгаты являются субстратами OATP: влияние на рост клеток рака молочной железы ZR-75-1. \textit{\foreignlanguage{english}{Molecular Nutrition \& Food Research}}. 2014;58:1830--1842. DOI: \href{https://doi.org/10.1002/mnfr.201400095}{\nolinkurl{10.1002/mnfr.201400095}}. PMID: \href{https://pubmed.ncbi.nlm.nih.gov/24996158/}{24996158}.
\item[{[94]}] Jaerapong N, Jamil QA, Riha J и др. Участие OATP в захвате куркумина и его основных метаболитов клетками рака молочной железы человека: влияние на противоопухолевую активность. \textit{\foreignlanguage{english}{Oncology Reports}}. 2019;41:2558--2566. DOI: \href{https://doi.org/10.3892/or.2019.7011}{\nolinkurl{10.3892/or.2019.7011}}. PMID: \href{https://pubmed.ncbi.nlm.nih.gov/30816509/}{30816509}.
\item[{[95]}] Kusuhara H, Furuie H, Inano A и др. Исследование фармакокинетического взаимодействия сульфасалазина у здоровых добровольцев и влияния куркумина как ингибитора BCRP in vivo. \textit{\foreignlanguage{english}{British Journal of Pharmacology}}. 2012;166:1793--1803. DOI: \href{https://doi.org/10.1111/j.1476-5381.2012.01887.x}{\nolinkurl{10.1111/j.1476-5381.2012.01887.x}}. PMID: \href{https://pubmed.ncbi.nlm.nih.gov/22300367/}{22300367}. Таблица 1: количественные данные для микродозы и терапевтической дозы.
\item[{[96]}] Glaeser H, Bailey DG, Dresser GK и др. Экспрессия кишечных транспортеров лекарственных средств и влияние грейпфрутового сока у человека. \textit{\foreignlanguage{english}{Clinical Pharmacology \& Therapeutics}}. 2007;81:362--370. DOI: \href{https://doi.org/10.1038/sj.clpt.6100056}{\nolinkurl{10.1038/sj.clpt.6100056}}. PMID: \href{https://pubmed.ncbi.nlm.nih.gov/17215845/}{17215845}.
\item[{[97]}] Poor M, Kaci H, Bodnarova S и др. Взаимодействия ресвератрола и его метаболитов --- ресвератрол-3-сульфата, ресвератрол-3-глюкуронида и дигидроресвератрола --- с сывороточным альбумином, ферментами CYP и транспортерами OATP. \textit{\foreignlanguage{english}{Biomedicine \& Pharmacotherapy}}. 2022;151:113136. DOI: \href{https://doi.org/10.1016/j.biopha.2022.113136}{\nolinkurl{10.1016/j.biopha.2022.113136}}. PMID: \href{https://pubmed.ncbi.nlm.nih.gov/35594715/}{35594715}.

\item[{[98]}] Управление по контролю качества пищевых продуктов и лекарственных средств США (FDA). \textit{Global Substance Registration System (GSRS): Curcumin}. UNII: IT942ZTH98. Молекулярная формула \(\mathrm{C}_{21}\mathrm{H}_{20}\mathrm{O}_{6}\); молярная масса 368.38 г/моль. Дата обращения: 9 сентября 2026 года. \href{https://precision.fda.gov/ginas/app/ui/substances/IT942ZTH98}{Официальный источник}.

\item[{[99]}] Wang S, Li F, Quan E, Dong D, Wu B. Характеристика эффлюкса глюкуронидов ресвератрола в клетках HeLa с экспрессией UGT1A1: клеточная фармакокинетическая модель для оценки вклада MRP4. \textit{\foreignlanguage{english}{Drug Metabolism and Disposition}}. 2016;44:485--488. DOI: \href{https://doi.org/10.1124/dmd.115.067710}{\nolinkurl{10.1124/dmd.115.067710}}. PMID: \href{https://pubmed.ncbi.nlm.nih.gov/26758854/}{26758854}.
\item[{[100]}] Wesolowska O, Lania-Pietrzak B, Kuzdzal M и др. Влияние силибина на биофизические свойства фосфолипидных бислоев. \textit{\foreignlanguage{english}{Acta Pharmacologica Sinica}}. 2007;28:296--306. DOI: \href{https://doi.org/10.1111/j.1745-7254.2007.00487.x}{\nolinkurl{10.1111/j.1745-7254.2007.00487.x}}. PMID: \href{https://pubmed.ncbi.nlm.nih.gov/17241534/}{17241534}.
\item[{[101]}] Naksuriya O, van Steenbergen MJ, Torano JS, Okonogi S, Hennink WE. Кинетическое исследование деградации свободного куркумина и куркумина в полимерных мицеллах. \textit{\foreignlanguage{english}{The AAPS Journal}}. 2016;18:777--787. DOI: \href{https://doi.org/10.1208/s12248-015-9863-0}{\nolinkurl{10.1208/s12248-015-9863-0}}. PMID: \href{https://pubmed.ncbi.nlm.nih.gov/27038456/}{27038456}. PMCID: \href{https://pmc.ncbi.nlm.nih.gov/articles/PMC5256596/}{PMC5256596}.
\item[{[102]}] Nguyen MH, Yu H, Kiew TY, Hadinoto K. Аморфный комплекс куркумина с хитозаном: высокая загрузка вещества и создание пересыщенных растворов как альтернатива наноинкапсуляции. \textit{\foreignlanguage{english}{European Journal of Pharmaceutics and Biopharmaceutics}}. 2015;96:1--10. DOI: \href{https://doi.org/10.1016/j.ejpb.2015.07.007}{\nolinkurl{10.1016/j.ejpb.2015.07.007}}. PMID: \href{https://pubmed.ncbi.nlm.nih.gov/26170159/}{26170159}.

\item[{[103]}] Izumi S, Nozaki Y, Lee W, Sugiyama Y. Экспериментальные и модельные данные в пользу транс-ингибирования как механизма зависящего от предварительной инкубации длительного подавления OATP1B1 циклоспорином A. \textit{\foreignlanguage{english}{Drug Metabolism and Disposition}}. 2022;50:541--551. DOI: \href{https://doi.org/10.1124/dmd.121.000783}{\nolinkurl{10.1124/dmd.121.000783}}. PMID: \href{https://pubmed.ncbi.nlm.nih.gov/35241487/}{35241487}.
\item[{[104]}] Gherbi K, Briddon SJ, Charlton SJ. Микрофармакокинетика: количественное определение локальной концентрации лекарственного вещества у мембран живых клеток. \textit{\foreignlanguage{english}{Scientific Reports}}. 2018;8:3479. DOI: \href{https://doi.org/10.1038/s41598-018-21100-x}{\nolinkurl{10.1038/s41598-018-21100-x}}. PMID: \href{https://pubmed.ncbi.nlm.nih.gov/29472588/}{29472588}. PMCID: \href{https://pmc.ncbi.nlm.nih.gov/articles/PMC5823863/}{PMC5823863}.
\item[{[105]}] de Witte WEA, Vauquelin G, van der Graaf PH, de Lange ECM. Влияние распределения лекарственного вещества и связывания с мишенью на ее занятость: приближение лимитирующей стадии. \textit{\foreignlanguage{english}{European Journal of Pharmaceutical Sciences}}. 2017;109 Suppl:S83--S89. DOI: \href{https://doi.org/10.1016/j.ejps.2017.05.024}{\nolinkurl{10.1016/j.ejps.2017.05.024}}. PMID: \href{https://pubmed.ncbi.nlm.nih.gov/28502676/}{28502676}.

\end{description}

\end{document}

